\RequirePackage[thinlines]{easybmat} %-- muss aufgrund eines Bugs vor etex und tikz geladen werden
\documentclass[a4paper, twoside, headsepline, index=totoc,toc=listof, fontsize=10pt, cleardoublepage=empty, headinclude, DIV=13, BCOR=13mm, titlepage,]{scrartcl}
\usepackage{scrtime} % KOMA, Uhrzeit ermoeglicht
\usepackage{scrpage2} % wie fancyhdr, nur optimiert auf KOMA-Skript, leicht andere Syntax
\usepackage{etoolbox}
\usepackage{letltxmacro}

% Compiler abhaengige IFs
\usepackage{ifxetex,ifluatex}
\newif\ifxetexorluatex
\ifxetex
  \xetexorluatextrue
\else
  \ifluatex
    \xetexorluatextrue
  \else
    \xetexorluatexfalse
  \fi
\fi

%--Farbdefinitionen
%-- muss vor tikz geladen werden
\usepackage[usenames, table, x11names]{xcolor}
\definecolor{dark_gray}{gray}{0.45}
\definecolor{light_gray}{gray}{0.6}

%--Zum Zeichnen
%-- muss vor polyglossia bzw. babel geladen werden
\usepackage{tikz}
\usepackage{tikz-cd}
\usetikzlibrary{external}
\tikzset{>=latex}
\usetikzlibrary{shapes,arrows,intersections}
\usetikzlibrary{calc,3d}
\usetikzlibrary{decorations.pathreplacing,decorations.markings}
\usetikzlibrary{angles}

%-- Konfiguration von tikzexternalize
\tikzexternalize[prefix=tikz/, up to date check=diff]

\ifxetexorluatex
\tikzset{external/system call={xelatex \tikzexternalcheckshellescape %-- verwende LuaLaTeX, wegen dynamischer Speicherallokation
    -halt-on-error -interaction=batchmode -jobname "\image" "\texsource"}}
\else
\tikzset{external/system call={pdflatex \tikzexternalcheckshellescape 
    -halt-on-error -interaction=batchmode -jobname "\image" "\texsource"}}
\fi

%-- tikzexternalize fuer tikzcd deaktivieren, da inkompatibel
\AtBeginEnvironment{tikzcd}{\tikzexternaldisable}
\AtEndEnvironment{tikzcd}{\tikzexternalenable}

%-- um Inkompatibilitaeten von quotes und polyglossia bzw. babel zu vermeiden
\tikzset{
  every picture/.append style={
    execute at begin picture={\shorthandoff{"}},
    execute at end picture={\shorthandon{"}}
  }
}
\usetikzlibrary{quotes}
\usepackage{pgfplots}
\usepgfplotslibrary{colormaps}
\pgfplotsset{compat=1.10}



%-- Mathe
\usepackage{mathtools} % beinhaltet amsmath
\mathtoolsset{showonlyrefs, centercolon}
\newtagform{brackets}[\textbf]{[}{]}
\usetagform{brackets}
\usepackage{wasysym}
\usepackage{amssymb} %zusätzliche Symbole
\usepackage{latexsym} %zusätzliche Symbole
\usepackage{stmaryrd} %für Blitz
\usepackage{nicefrac} %schräge Brüche
\usepackage{cancel} %Befehle zum Durchstreichen
\usepackage{mathdots}
% \usepackage{bm}
\DeclareSymbolFont{bbold}{U}{bbold}{m}{n}
\DeclareSymbolFontAlphabet{\mathbbold}{bbold}
\newcommand{\mathds}[1]{\mathbb{#1}} %Um Kompatibilitaet mit frueheren Benutzung von dsfont herzustellen
\newcommand{\ind}{\mathbbold{1}} % charakteristische-Funktion-Eins
\def\mathul#1#2{\color{#1}\underline{{\color{black}#2}}\color{black}} %farbiges Untersteichen im Mathe-Modus

%-- Alles was mit Schrift und XeteX zu tun hat
\ifxetexorluatex
    % XeLaTeX or LuaTeX
	\ifxetex
		\usepackage{mathspec} %beinhaltet fontspec 
	\else
		\usepackage[no-math]{fontspec}
	\fi
    \usepackage{polyglossia} %babel-ersatz
    \setmainlanguage[spelling=new,babelshorthands=false]{german}
	\newcommand\glqq{"}
	\newcommand\grqq{"}
	\defaultfontfeatures{Mapping=tex-text, WordSpace={1.4}} %
    \setmainfont[Ligatures=Common, BoldFont={* Bold}, ItalicFont={* Light Italic},SmallCapsFont={Linux Libertine O}, SmallCapsFeatures={Letters=SmallCaps}]{Source Sans Pro}
	\setsansfont[Scale=MatchLowercase,Ligatures=Common, BoldFont={* Medium}]{Ubuntu}
	\ifxetex
		\setallmonofonts[Scale=MatchLowercase, ItalicFont={*}]{Consolas} 
	\else
		\setmonofont[Scale=MatchLowercase, ItalicFont={*}]{Consolas}
	\fi
	\usepackage{xltxtra}
	\usepackage{fontawesome}
	\usepackage{microtype}
\else
    % default: pdfLaTeX
    \usepackage[ngerman]{babel}
    \usepackage[T1]{fontenc}
    \usepackage[utf8]{inputenc}
	\usepackage{textcomp} %verhindert ein paar Fehler bei den Fonts
    \usepackage[babel=true, tracking=true]{microtype}
	\usepackage{lmodern}
	\renewcommand{\familydefault}{\sfdefault} 
\fi


%--Mixed
\usepackage[neverdecrease]{paralist}
\usepackage[german=quotes]{csquotes}
\usepackage{makeidx}
\usepackage{booktabs}
\usepackage{wrapfig}
\usepackage{float}
\usepackage[margin=10pt, font=small, labelfont={sf, bf}, format=plain, indention=1em]{caption}
\captionsetup[wrapfigure]{name=Abb. }
\usepackage{stackrel}
\usepackage{ifthen}
\usepackage{multicol}
\usepackage{qrcode}
\flushbottom

%--Hyphenation
\hyphenation{di-men-si-o-na-le}


%--Unterstreichung
\usepackage[normalem]{ulem}
\setlength{\ULdepth}{1.8pt}

%--Indexverarbeitung
\newcommand{\bet}[1]{\textbf{#1}}
\newcommand{\Index}[1]{\textbf{#1}\index{#1}}
\makeindex
\setindexpreamble{{\noindent \itshape Die \emph{Seitenzahlen} sind mit Hyperlinks zu den entsprechenden Seiten versehen, also anklickbar} \ifxetexorluatex \faHandUp\fi \par \bigskip}
\renewcommand{\indexpagestyle}{scrheadings}


%--Marginnote & todonotes
\usepackage{marginnote}
\renewcommand*{\marginfont}{\color{dark_gray} \itshape \footnotesize }
\usepackage[textsize=small]{todonotes}

%--Todonotes mit tikz/externalize kompatibel machen
\LetLtxMacro{\oldtodo}{\todo}
\renewcommand{\todo}[2][]{\tikzexternaldisable\oldtodo[#1]{#2}\tikzexternalenable}
\LetLtxMacro{\oldmissingfigure}{\missingfigure}
\renewcommand{\missingfigure}[2][]{\tikzexternaldisable\oldmissingfigure[{#1}]{#2}\tikzexternalenable}

%--Konfiguration von Hyperref pdfstartview=FitH, 
\usepackage[hidelinks, pdfpagelabels,  bookmarksopen=true, bookmarksnumbered=true, linkcolor=black, urlcolor=SkyBlue2, plainpages=false,pagebackref, citecolor=black, hypertexnames=true, pdfauthor={Jannes Bantje}, pdfborderstyle={/S/U}, linkbordercolor=SkyBlue2, colorlinks=false]{hyperref}

% \let\orighref\href
\ifxetexorluatex
\newcommand{\hrefsym}[2]{\href{#1}{\texttt{#2\,\faExternalLink}}}
\newcommand{\hrefsymmail}[2]{\href{#1}{\texttt{\faEnvelope\,#2}}}
\renewcommand{\url}[1]{\hrefsym{#1}{\nolinkurl{#1}}}
\fi
\providecommand{\hrefsym}{\href}
% \usepackage{cleveref}
% \crefformat{footnote}{#2\footnotemark[#1]#3}

%--Römische Zahlen
\newcommand{\RM}[1]{\MakeUppercase{\romannumeral #1{}}}

%--Punkte (nach hyperref)
\usepackage{ellipsis}

%%--Abkürzungen etc.
\newcommand{\light}{\text{\Large $\lightning$}}

%-- Definitionen von weiteren Mathe-Befehlen
\DeclareMathOperator{\re}{Re} %Realteil
\DeclareMathOperator{\im}{Im} %Imaginaerteil
\DeclareMathOperator{\id}{id} %identische Abbildung
\DeclareMathOperator{\Sp}{Sp} %Spur
\DeclareMathOperator{\sgn}{sgn} %Signum
\DeclareMathOperator{\alt}{Alt} %Alternierende n-linearFormen
\DeclareMathOperator{\End}{End} %Endomorphismen
\DeclareMathOperator{\Vol}{Vol} %Volumen
\DeclareMathOperator{\dom}{dom} %Domain
\DeclareMathOperator{\Hom}{Hom} %Homomorphismen
\DeclareMathOperator{\bild}{Bild} %Bild
\DeclareMathOperator{\Kern}{Kern }
\DeclareMathOperator{\rg}{rg} %Rang
\DeclareMathOperator{\Rg}{Rg} %Rang
\DeclareMathOperator{\diam}{diam} %Durchmesser
\DeclareMathOperator{\dist}{dist} %Distanz
\DeclareMathOperator{\grad}{grad} %Gradient
\DeclareMathOperator{\dive}{div} %Gradient
\DeclareMathOperator{\rot}{rot} %Rotation
\DeclareMathOperator{\hess}{Hess} %Hesse-Matrix
\DeclareMathOperator{\koker}{Koker} %Kokern
\DeclareMathOperator{\aut}{Aut} %Automorphismen
\DeclareMathOperator{\ord}{ord} %Ordnung
\DeclareMathOperator{\ggT}{ggT} %ggT
\DeclareMathOperator{\kgV}{kgV} %kgV
\DeclareMathOperator{\Gr}{Gr} %Gerade
\DeclareMathOperator{\Kr}{Kr} %Kreis
\DeclareMathOperator{\Char}{char} %Charakteristik
\DeclareMathOperator{\Aut}{Aut} %Automorphismen
% \DeclareMathOperator{\D}{D} %formale Ableitung
\newcommand{\D}{\ensuremath{\mathrm{D}\mkern-1.5mu}}
\newcommand{\mathd}{\ensuremath{\mathrm{d}\mkern-1.0mu}}
\newcommand{\Tmap}{\ensuremath{\mathrm{T}\mkern-1.0mu}}
\newcommand{\opL}{\ensuremath{\mathrm{L}\mkern-0.6mu}}
\DeclareMathOperator{\cov}{cov} %Kovarianz
\DeclareMathOperator{\Gal}{Gal} %Galois
\DeclareMathOperator{\supp}{supp} %Träger
\DeclareMathOperator{\Sym}{Sym} %Symmetrische Gruppe
\DeclareMathOperator{\Zyl}{Zyl} %Zylinder
\DeclareMathOperator{\Mod}{Mod} %Moduln
\DeclareMathOperator{\EPK}{EPK} %Einpunktkompaktifizierung
\DeclareMathOperator{\conj}{conj} %Konjugation
\DeclareMathOperator{\ann}{ann} %Annulator
\DeclareMathOperator{\tor}{tor} %Torsionsmodul
\DeclareMathOperator{\Int}{Int} %Interior/Inneres
\DeclareMathOperator{\Br}{Br} %Brauerergruppe
\DeclareMathOperator{\GL}{GL} %allgemeine lineare Gruppe
\DeclareMathOperator{\SL}{SL} %spezielle lineare Gruppe
\DeclareMathOperator{\SU}{SU} %spezielle unitäre Gruppe
\DeclareMathOperator{\Tr}{Tr} %Spur/Trace
\DeclareMathOperator{\conv}{conv} %Konvexe Teilmenge
% \DeclareMathOperator{\tr}{tr}
\newcommand{\tr}{\mathrm{tr}}
\newcommand{\w}{\mathrm{w}}
\newcommand{\weakT}[1]{\ensuremath{\mathcal{T}_{#1}^{\mkern+1.0mu\text{\raisebox{0.4ex}{$\mathrm{w}$}}}}}
\newcommand{\weakTstar}[1]{\ensuremath{\mathcal{T}_{#1}^{\mkern+1.0mu\text{\raisebox{0.4ex}{$\mathrm{w}$}}^*}}}
\newcommand{\rand}[1]{\ensuremath{\partial^{\scriptscriptstyle #1}}}


% -- Kategorien
\DeclareMathOperator{\Ob}{Ob} %Objekte
\DeclareMathOperator{\Mor}{Mor} %Morphismen

% \DeclareMathOperator{\TOP}{TOP} %Topologische Räume
\newcommand{\TOP}{\textsc{Top}}
% \DeclareMathOperator{\HTOP}{HTOP} %Topologische Räume + Homotopieäquivalenz
\newcommand{\HTOP}{\textsc{HTop}}
% \DeclareMathOperator{\VR}{VR} %Vektorräume
\newcommand{\VR}{\textsc{VR}}
% \DeclareMathOperator{\MOD}{MOD} %Moduln
\newcommand{\MOD}{\textsc{Mod}}
% \DeclareMathOperator{\MENGEN}{MENGEN} %Mengen
\newcommand{\MENGEN}{\textsc{Mengen}}
% \DeclareMathOperator{\MAN}{MAN} %Mannigfaltigkeiten
\newcommand{\MAN}{\textsc{Man}}
% \DeclareMathOperator{\GRUPPEN}{GRUPPEN} %Gruppen
\newcommand{\GRUPPEN}{\textsc{Gruppen}}
% \DeclareMathOperator{\ABELGRUPPEN}{ABEL.GRUPPEN} %abelsche Gruppen
\newcommand{\ABELGRUPPEN}{\textsc{Abel.Gruppen}}
% \DeclareMathOperator{\KAT}{KAT} %Kategorien
\newcommand{\KAT}{\textsc{Kat}}
% \DeclareMathOperator{\FUN}{FUN} %Funktoren
\newcommand{\FUN}{\textsc{Fun}}


\newcommand{\op}{\mathrm{op}}

\newcommand{\Ker}{\ensuremath{\Kern \,}}
\newcommand{\zirkel}{\rotatebox[origin = c]{-90}{$\varangle$}}
\newcommand{\wirkung}{\!\curvearrowright\!}

\newcommand{\oE}{\OE}

%--Beweisende
\newcommand{\bewende}{\ifmmode \tag*{$\square$} \else \hfill $\square$ \fi}

%--Volumen
\newcommand{\vol}[1]{\ensuremath{\Vol_{#1}}}

%--Skalarprodukt
\DeclarePairedDelimiterX\sprod[2]{\langle}{\rangle}{#1\,\delimsize\vert\,#2}

%--Betrag, Gaußklammer
\DeclarePairedDelimiter{\abs}{\lvert}{\rvert}
\DeclarePairedDelimiter{\ceil}{\lceil}{\rceil}

%--Norm
\DeclarePairedDelimiter\doppelstrich{\Vert}{\Vert}
\newcommand{\norm}[2][\relax]{
\ifx#1\relax \ensuremath{\doppelstrich*{#2}}
\else \ensuremath{\doppelstrich*{#2}_{#1}}
\fi}


%--Umklammern
\DeclarePairedDelimiter\enbrace{(}{)}
\DeclarePairedDelimiter\benbrace{[}{]}
\DeclarePairedDelimiter\lenbrace{<}{>}

%--Mengen
\DeclarePairedDelimiterX\mengenA[1]{\lbrace}{\rbrace}{#1}
\DeclarePairedDelimiterX\mengenB[2]{\lbrace}{\rbrace}{#1\, \delimsize\vert \, #2}

\makeatletter
\newcommand{\set}[2][\relax]{
\ifx#1\relax \ensuremath{
\mengenA*{#2}}
\else \ensuremath{%
  \mengenB*{#1}{#2}}
\fi}
\makeatother


%--offen und abgeschlossen
\newcommand{\off}{\! \stackrel[\text{offen}]{}{\subset} \!}
\newcommand{\abg}{\! \stackrel[\text{abg.}]{}{\subset} \!}

%--Differential
\newcommand{\diff}[2]{\ensuremath{\frac{{\partial #1}}{{\partial #2}} }}
\newcommand{\diffd}[2]{\ensuremath{\frac{\mathd #1}{\mathd #2} }}

%--Diagonale Linien für Matrizen
\newcommand\tikzmark[1]{%
  \tikzexternaldisable\tikz[overlay,remember picture,baseline] \node [anchor=base] (#1) {};\tikzexternalenable}

\newcommand\MyLine[3][]{%
  \tikzexternaldisable\begin{tikzpicture}[overlay,remember picture]
    \draw[#1] (#2.south) -- (#3.north);
  \end{tikzpicture}\tikzexternalenable}
  
%--Jordankasten
\newcommand{\Jord}{ \tikzexternaldisable \begin{BMAT}{|ccc|}{|ccc|} %
				0 & & \\ %
				1\tikzmark{z} &  \diagdown & \\ %
				& \tikzmark{u}1 \MyLine[thick]{z}{u} &  0 %
			\end{BMAT} \tikzexternalenable}
			
%--Jordankasten mit Lambda
\newcommand{\JordLa}[1][\empty]{%
	\tikzexternaldisable\ifthenelse{\equal{#1}{\empty}}
		{\begin{BMAT}{|ccc|}{|ccc|} %
				\lambda & & \\ %
				1\tikzmark{z} &  \diagdown & \\ %
				& \tikzmark{u}1 \MyLine[thick]{z}{u} &  \lambda  %
			\end{BMAT}}
		{\begin{BMAT}{|ccc|}{|ccc|} %
				\lambda_{#1}  & & \\ %
				1\tikzmark{z} &  \diagdown & \\ %
				& \tikzmark{u}1 \MyLine[thick]{z}{u} &  \lambda_{#1}  %
			\end{BMAT}}\tikzexternalenable}

%--Inhaltsverzeichnis
\usepackage[tocindentauto]{tocstyle}
\usetocstyle{KOMAlike}			

\newcommand{\fach}{Analysis }
\newcommand{\semester}{WiSe 2013}
\newcommand{\homepage}{http://wwwmath.uni-muenster.de/u/wilhelm.winter/wwinter/analysis_III.html}
\newcommand{\prof}{Prof.\,Dr.\,Wilhelm Winter}
\providecommand{\verfasser}{Jannes Bantje}

%--Konfiguration von scrheadings
\setheadsepline{1pt}[\color{light_gray}]
\pagestyle{scrheadings}
%\fancyhf[H, F]{}
\clearscrheadfoot

\providecommand{\shortFach}{\fach}
\lehead{ \includegraphics[height=0.5 cm,keepaspectratio]{../!config/Bilder/Logo_WWU_Muenster_light_gray.pdf}}
\rehead{\rule{0cm}{0.5cm}\footnotesize \sffamily \color{light_gray} \verfasser -- Mitschrift \shortFach}
\lohead{\rule{0cm}{0.5cm} \footnotesize \sffamily \color{light_gray} Stand: \today \; \thistime[:]}
\rohead{\includegraphics[height=0.5 cm,keepaspectratio]{../!config/Bilder/fb10logo_gray.pdf}}


\ofoot[{ \color{dark_gray} \LARGE \sffamily \thepage}]{{ \color{dark_gray} \LARGE \sffamily \thepage}} %hier wir auch der plain Stil bearbeitet!
\automark{section}
\ifoot{ \color{dark_gray} \small \leftmark}

%--Metadaten
\author{\verfasser}
\titlehead{\includegraphics[height=1.5cm, keepaspectratio]{../!config/Bilder/Logo_WWU_Muenster.pdf}%
\hfill \includegraphics[height=1.3cm, keepaspectratio]{../!config/Bilder/fb10logo.pdf}}
\title{Skript \fach}
\subtitle{Mitschrift der Vorlesung  \enquote{\fach} von \prof}

\publishers{
Aktuelle Version verfügbar bei: \bigskip\\
\normalsize
\begin{minipage}[t]{0.45\textwidth}
	\centering{
	\qrcode[height=3.5cm, version=4]{https://github.com/JaMeZ-B/latex-wwu}
	\medskip\\
	\includegraphics[height=0.4cm, keepaspectratio]{../!config/Bilder/github_octo.pdf} 
	\includegraphics[height=0.4cm, keepaspectratio]{../!config/Bilder/GitHub_Logo.pdf} (inklusive Sourcecode)\\
	{\small\url{https://github.com/JaMeZ-B/latex-wwu}}}
\end{minipage}
\hfill
\begin{minipage}[t]{0.45\textwidth}
	\centering{
	\qrcode[height=3.5cm, version=4]{btsync://B6WH2DISQ5QVYIRYIEZSF4ZR2IDVKPN3I?n=latex_share}
	\medskip\\
	\includegraphics[height=0.4cm, keepaspectratio]{../!config/Bilder/bt_sync_logo.pdf} 
	{\large \textbf{Bittorrent} Sync} \\
	\texttt{B6WH2DISQ5QVYIRYIEZSF4ZR2IDVKPN3I}
	}
\end{minipage}
}



\begin{document}
\pagenumbering{Roman}
\maketitle
\begin{abstract}
\section*{Vorwort --- Mitarbeit am Skript}
Dieses Dokument ist eine Mitschrift aus der Vorlesung \enquote{\fach, \semester}, gelesen von \prof. Der Inhalt entspricht weitestgehend dem Tafelanschrieb. Für die
Korrektheit des Inhalts übernehme ich keinerlei Garantie! Für Bemerkungen und Korrekturen -- und seien es nur Rechtschreibfehler -- bin ich sehr dankbar. 
Korrekturen lassen sich prinzipiell auf drei Wegen einreichen: 
\begin{itemize}
	\item Persönliches Ansprechen in der Uni, Mails an \hrefsymmail{mailto:\mail}{\mail} (gerne auch mit annotieren PDFs) 
	\item \emph{Direktes} Mitarbeiten am Skript: Den Quellcode poste ich auf GitHub (siehe oben), also stehen vielfältige Möglichkeiten der Zusammenarbeit zur Verfügung:
	Zum Beispiel durch Kommentare am Code über die Website und die Kombination Fork + Pull Request. Wer sich verdient macht oder ein Skript zu einer Vorlesung, die 
	ich nicht besuche, beisteuern will, dem gewähre ich gerne auch Schreibzugriff.
	
	Beachten sollte man dabei, dass dazu ein Account bei \url{github.com} notwendig ist, der allerdings ohne Angabe von persönlichen Daten angelegt werden kann. 
	Wer bei GitHub (bzw. dem zugrunde liegenden Open-Source-Programm \enquote{\texttt{git}}) -- verständlicherweise -- Hilfe beim Einstieg braucht, dem helfe ich gerne 
	weiter. Es gibt aber auch zahlreiche empfehlenswerte Tutorials im Internet.\footnote{zB. \url{https://try.github.io/levels/1/challenges/1}, ist auf Englisch, aber dafür 
	interaktives LearningByDoing}
	\item \emph{Indirektes} Mitarbeiten: \TeX-Dateien per Mail verschicken. 
	
	Dies ist nur dann sinnvoll, wenn man einen ganzen Abschnitt ändern möchte (zB. einen alternativen Beweis geben), da ich die Änderungen dann per Hand einbauen muss! Ich freue mich aber auch über solche Beiträge!
\end{itemize}
\section*{Vorlesungshomepage}
{\centering 
\begin{minipage}[c][][c]{\textwidth}
	\centering \qrcode[height=3cm]{\homepage} \medskip\\
	\footnotesize \url{\homepage}
\end{minipage}
\par}

\end{abstract}

\setcounter{page}{1}
\tableofcontents
\cleardoubleoddemptypage
\pagenumbering{arabic}
\setcounter{page}{1}

\section{Gewöhnliche Differentialgleichungen – Existenz und Eindeutigkeit} % (fold)
\label{sec:1}

\subsection[Definition: Gewöhnliche Differentialgleichungen]{Definition} % (fold)
\label{sub:11}
\begin{enumerate}[a)]
	\item Sei $G \subset \mathds{R}^2, f : G \to \mathds{R}$ stetig. 
	\[
		y' = f(x,y) \tag{$\star$}
	\] 
	heißt \bet{gewöhnliche Differentialgleichung 1. Ordnung}. \index{Differentialgleichung!1. Ordnung} Eine Lösung von ($\star$) auf dem Intervall $I$ ist eine differenzierbare Funktion 
	$\varphi : I \to \mathds{R}$ mit
	\begin{enumerate}[(i)]
		\item $\set[\big(x, \varphi(x)\big)]{x \in I} \subset G $
		\item $\varphi ' (x) = f\big(x, \varphi(x)\big)$, $x  \in I$
	\end{enumerate}
	\item Sei $G \subset \mathds{R} \times \mathds{R}^n$, $f : G \to \mathds{R}^n$ stetig.
	\[
		y' = f(x,y) \tag{$\star \star$}
	\]
	heißt \bet{System von $n$ Differentialgleichungen 1. Ordnung}\index{Differentialgleichung!System von $n$}. Eine Lösung von $(\star \star)$ auf $I$ ist eine differenzierbare Funktion 
	$\varphi : I \to \mathds{R}^n$ mit
	\begin{enumerate}[(i)]
		\item $\set[\big(x, \varphi(x)\big)]{x \in I} \subset G$
		\item $\varphi'(x) = f \big(x,\varphi(x)\big)$, $x \in I$
	\end{enumerate}
	$(\star \star)$ und (ii) lassen sich auch schreiben als 
	\[
		\begin{pmatrix}
			y_1' \\ \vdots \\ y_n'
		\end{pmatrix} = \begin{pmatrix}
			f_1(x, y_1, \ldots , y_n) \\
			\vdots \\
			f_n(x, y_1, \ldots , y_n)
		\end{pmatrix} \qquad , \qquad \begin{pmatrix}
			\varphi_1'(x) \\
			\vdots \\
			\varphi_n'(x)
		\end{pmatrix} = \begin{pmatrix}
			f_1\big(x, \varphi_1(x), \ldots , \varphi_n(x)\big) \\
			\vdots \\
			f_n\big(x, \varphi_1(x), \ldots , \varphi_n(x)\big)
		\end{pmatrix}
	\]
	\item Sei $G \subseteq \mathds{R} \times \mathds{R}^n$, $f : G \to \mathds{R}$ stetig
	\[
		y^{(n)} = f \enbrace{x,y,y',y'', \ldots , y^{(n-1)}} \tag{$\star\star\star$}
	\]
	heißt \bet{Differentialgleichung $n$-ter Ordnung}\index{Differentialgleichung!$n$-ter Ordnung}. Eine Lösung von $(\star \star \star)$ ist eine $n$-mal differenzierbare
	Funktion $\varphi : I \to \mathds{R}$ mit
	\begin{enumerate}[(i)]
		\item $\set[\big(x, \varphi(x), \varphi'(x), \ldots , \varphi^{(n-1)}(x)\big)]{x \in I} \subset G $
		\item $\varphi^{(n)}(x) = f\Big(x, \varphi(x), \varphi'(x), \ldots , \varphi^{(n-1)} (x)\Big)$, $x \in I$
	\end{enumerate}
\end{enumerate}
% subsection 11 (end)

\subsection[Beispiel: Einfache Differentialgleichungen]{Beispiel} % (fold)
\label{sub:12}
\begin{enumerate}[(i)]
	\item \label{enum:12:1} Sei $y'= \frac{y}{x} $, $G= \mathds{R}_+^* \times \mathds{R}$. Lösungen sind Geraden der Form 
	$\varphi(x) = c \cdot x$, wo $c \in \mathds{R}$ konstant ist. Dann gilt mit $I = (0, \infty)$
	\[
		\varphi'(x) = \frac{c \cdot x}{x} = \frac{\varphi(x)}{x} 
	\]
	\item \label{enum:12:2} Sei $y'= - \frac{x}{y} $, $G = \mathds{R} \times \mathds{R}_+^*$. Lösungen sind Halbkreise der Form 
	$\varphi(x) = \sqrt{c-x^2} , x \in \enbrace{- \sqrt{c}, \sqrt{c}}=I , c \in \mathds{R}_+ $
	\[
		\varphi'(x) = \frac{-2x}{2 \sqrt{c-x^2}  } = - \frac{x}{\varphi(x)}  
	\]
	\item $m \cdot \ddot x = F(t,x, \dot x)$ \marginnote{$t$ Zeit, $x$ Ort, $\dot x$ Geschwindigkeit, $m$ Masse, $\ddot x$ Beschleunigung, $F$ Kraft}
\end{enumerate}
% subsection 12 (end)

\subsection[Bemerkung: geometrische Interpretationen und System von DGL $\Leftrightarrow$ DGL $n$-ter Ordnung]{Bemerkung} % (fold)
\label{sub:13}
\begin{enumerate}[(i)]
	\item geometrische Interpretationen:
	\begin{itemize}
		\item allgemeine Darstellung:
		\[
			\begin{tikzpicture}[scale=0.6]
				%\draw[help lines] (0,0) grid (10,6);
				\coordinate (x0) at (7,0);
				\coordinate (y0) at (0,3);
				\coordinate (xy) at (7,3);
				\coordinate (fpeak) at (10,4.5);
				\draw[ultra thick, <->] (10,0) -- (0,0) -- (0,6);
				\draw[ultra thick, Chartreuse4] (2,1) to [out=-5, in=200](6,1) to [out=20, in=260](8,3) to [out=80, in=-85](8,5) to [out=95, in=-20](4.5,5.5) 
				to [out=160, in=5](2,5) to [out=185, in=80](1,2) to [out=260, in=175](2,1);
				\draw[decorate, decoration={brace,amplitude=5pt, mirror, raise=4pt},yshift=0pt, thick] (10,3) -- (fpeak) node[right, midway, xshift=0.25cm]{$f(x_0,y_0)$};
				\draw[decorate, decoration={brace,amplitude=5pt, mirror, raise=4pt},yshift=0pt, thick] (xy) -- (10,3) node[below, midway, yshift=-0.27cm]{$1$};
				\node[Chartreuse4] at (2.5,5.8) {$G$};
				\draw[fill] (x0) circle (0.07) node[below]{$x_0$};
				\draw[fill] (y0) circle (0.07) node[left]{$y_0$};
				\draw[dashed, thick] (x0) -- (xy);
				\draw[dashed, thick] (y0) -- (xy) -- ++ (3,0) -- (fpeak);
				\draw (2,2.3) -- ++ (130:0.7);
				\draw (2,3) -- ++ (122:0.7);
				\draw (2,4) -- ++ (116:0.7);
				\draw (2,1.4) -- ++ (140:0.7);
				\draw (3,4) -- ++ (110:0.7);
				\draw (3,3) -- ++ (130:0.7);
				\draw (3,2) -- ++ (165:0.7);
				\draw (3,1.05) -- ++ (167:0.7);
				\draw (4,4) -- ++ (125:0.7);
				\draw (4,5) -- ++ (108:0.5);
				\draw (4,2.7) -- ++ (160:0.7);
				\draw (4,1.85) -- ++ (175:0.7);
				\draw (4,1) -- ++ (177:0.7);
				\draw (5,1) -- ++ (185:0.7);
				\draw (5,1.84) -- ++ (188:0.3);
				\draw (5,2.8) -- ++ (185:0.7);
				\draw (5,3.7) -- ++ (179:0.7);
				\draw (5,4.6) -- ++ (170:0.7);
				\draw (6,2.1) -- ++ (195:0.7);
				\draw (6,1.3) -- ++ (198:0.7);
				\draw (6,3) -- ++ (195:0.7);
				\draw (6,4.15) -- ++ (208:0.7);
				\draw (6,5.2) -- ++ (215:0.7);
				\draw (6.8,1.6) -- ++ (205:0.5);
				\draw (6.8,2.4) -- ++ (208:0.5);
				\draw (7,3.8) -- ++ (220:0.7);
				\draw (6.7,5.2) -- ++ (245:0.7);
				\draw (7.7,5) -- ++ (245:0.7);
				\draw[ultra thick, DarkOrange2] (8,4.3) to [out=-110,in=26.565](xy) to [out=206.565, in=0](4.5,2.2) node[below]{$\varphi$} to [out=180, in=-80](1.2,4);
				\draw[ultra thick, DarkOrange2] (7,5.6) to [out=270, in=0](4.5,3.2) to [out=180, in=-80](3,5.3);
				\draw[ultra thick, Firebrick3] (6.5,2.75) -- ++(26.565:4.5) node[right]{$y'=f(x,y)$};
			\end{tikzpicture}
		\]
		\item zu \ref{sub:12} (\ref{enum:12:1}) und \ref{sub:12} (\ref{enum:12:2})
		\[
			\vcenter{\hbox{
			\begin{tikzpicture}[scale=0.8]
				%\draw[help lines] (0,-3) grid (5,3);
				\draw[ultra thick] (0,-3) -- (0,3);
				\draw[ultra thick] (0,0) -- (3.1,0);
				\foreach \x in {1,...,16}
					\draw[dashed] (0,0) -- ++ (90-11.25*\x:3); 
				\draw[ultra thick, Sienna3] (0,0) -- ++ (22.5:3.2) node[right]{$\varphi$};
			\end{tikzpicture}}}
			\qquad 
			\vcenter{\hbox{
			\begin{tikzpicture}[scale=0.8]
				%\draw[help lines] (-3,0) grid (3,3);
				\draw[ultra thick] (-3.2,0) -- (3.2,0);
				\draw[ultra thick] (0,0) -- (0,3.2);
				\foreach \x in {1,...,4}
					\draw[thick, dashed] (\x*0.75,0) arc [radius=\x*0.75, start angle=0, end angle=180];
				\draw[ultra thick, Sienna3](1.5,0) arc [radius=1.5, start angle=0, end angle=180] node[below]{$\varphi$};
			\end{tikzpicture}}}
		\]
		Symmetrie bei \ref{sub:12} (\ref{enum:12:2}): $f(\lambda x, \lambda y)= f(x,y)$ für $(x,y) \in G, \lambda  \in \mathds{R}$
	\end{itemize}
	\item\label{enum:13.2} Sei $G \subset \mathds{R} \times \mathds{R}^n$, $f : G \to \mathds{R}$ stetig und sei
	\[
		y^{(n)} = f \enbrace{x, y, y', \ldots , y^{(n-1)}} \tag{$\star$}
	\]
	eine DGL $n$-ter Ordnung. Betrachte nun die Gleichung
	\[
		\begin{pmatrix}
			y_0' \\ \vdots  \\ y_{n-2}' \\ y_{n-1}'
		\end{pmatrix} = \begin{pmatrix}
			y_1 \\ \vdots \\ y_{n-1} \\ f(x,y_0, \ldots , y_{n-1})
		\end{pmatrix} \tag{$\star\star$}
	\]
	mit $Y = \enbrace{ \begin{smallmatrix}
		y_0 \\ : \\ y_{n-1}
	\end{smallmatrix}} , F(x,Y) = \enbrace{ \begin{smallmatrix}
		y_1 \\ : \\ y_{n-1} \\ f(x, y_0, \ldots , y_{n-1})
	\end{smallmatrix}} $ schreibt sich $(\star\star)$ als 
	\[
		Y' = F(x,Y) \tag{$\star\star\star$}
	\]
	ein System von DGL erster Ordnung. Falls $\varphi : I \to \mathds{R}$ Lösung von $(\star)$ ist, so ist $\Phi : I \to \mathds{R}^n$ mit
	\[
		\Phi (x) = \begin{pmatrix}
			\varphi(x) \\ \varphi'(x) \\ \vdots \\ \varphi^{(n-1)}(x)
		\end{pmatrix}
	\]
	eine Lösung von $(\star\star\star)$. Beweis:
	\[
		\Phi'(x) = \begin{pmatrix}
			\varphi'(x) \\ \vdots \\ \varphi^{(n-1)}(x) \\ \varphi^{(n)}(x)
		\end{pmatrix} = \begin{pmatrix}
			\varphi'(x) \\ \vdots \\ \varphi^{(n-1)}(x) \\ f\big(x, \varphi(x), \ldots , \varphi^{(n-1)}(x)\big) 
		\end{pmatrix} = F\big(x,\Phi(x)\big)
	\]
	Falls umgekehrt $\Phi = \begin{pmatrix}
		\varphi_0 \\ \vdots \\ \varphi_{n-1}
	\end{pmatrix} : I  \to \mathds{R}^n$ Lösung von $(\star \star \star)$ ist, so ist $\varphi_0$ Lösung von $(\star)$. \marginnote{(warum?)}
	
\end{enumerate}
% subsection 13 (end)


\subsection[Definition: Lipschitz bezüglich der 2. Variablen]{Definition} % (fold)
\label{sub:14}
Sei $G \subset \mathds{R} \times \mathds{R}^n, f : G \to \mathds{R}^n$. 
\begin{enumerate}[(i)]
	\item $f$ ist Lipschitz (bezüglich der zweiten Variablen), falls $L \ge 0$ existiert, sodass für $(x,y) , (x,\overline{y} ) \in G$ gilt:
	\[
		\norm[2]{f(x,y) - f(x,\overline{y} )} \le L \cdot \norm[2]{y- \overline{y} }  
	\]
	\item $f$ ist lokal Lipschitz (bezüglich der zweiten Variablen), falls gilt: Jedes $(a,b) \in G$ besitzt eine Umgebung $U$, sodass $f|_{U \cap G}$ Lipschitz ist.
\end{enumerate}
% subsection 14 (end)

\subsection[Satz: Kriterium für lokal Lipschitz (stetig partiell differenzierbar bzgl. der 2. Variablen)]{Satz} % (fold)
\label{sub:15}
Sei $G \subset \mathds{R} \times \mathds{R}^n$ offen und $f : G \to \mathds{R}^n, (x,y) \mapsto f(x,y)$, stetig partiell differenzierbar bezüglich der 2. Variablen, d.h.
\[
	(x,y) \mapsto \diff{f}{y} (x,y) = \enbrace{\diff{f_i}{y_j} (x,y)}_{1 \le i,j \le n}
\]
ist stetig. Dann ist $f$ lokal Lipschitz.
\minisec{Beweis}
Sei $(a,b) \in G$. Es gibt $r>0$ sodass $K:= [a-r,a+r]\times \overline{B}_{\norm[2]{.} }(b,r) \subset G $ \hfill (warum?)\\
Dann ist $K$ kompakt nach Heine-Borel. Für $(x,y) \in G$ ist $\diff{f}{y} (x,y) \in M_n(\mathds{R})$ und $(x,y) \mapsto \norm{\diff{f}{y} (x,y)} $ stetig. $K$ kompakt
$\overset{\text{Ana II., 4.10}}{\Rightarrow} $ $L := \sup \set[ \norm{\diff{f}{y} (x,y)} ]{(x,y) \in K} < \infty$. Für $(x,y), (x,\overline{y} ) \in K$ gilt nach dem
Mittelwertsatz (Ana II., 7.11)
\[
	\norm[2]{f(x,y) - f(x,\overline{y} )} \le L \cdot \norm[2]{y-\overline{y} }  \bewende
\]
% subsection 15 (end)

\subsection[Satz: Eindeutigkeit von Lösungen, wenn $f$ lokal Lipschitz]{Satz (Eindeutigkeit)} % (fold)
\label{sub:16}
Sei $G \subset \mathds{R} \times \mathds{R}^n$, $f :G  \to \mathds{R}^n$ stetig und lokal Lipschitz. Seien $\varphi, \psi : I \to \mathds{R}^n$ Lösungen von $y'= f(x,y)$.
Falls $\varphi(x_0)= \psi(x_0)$ für ein $x_0 \in I$, so gilt $\varphi(x)= \psi(x)$ für alle $x \in I$.
\minisec{Beweis}
\begin{description}
	\item[1. Behauptung] Falls $\varphi(a)= \psi(a)$ für $a \in I$, so existiert $\varepsilon >0 : \varphi(x)= \psi (x)$ für $x \in [a- \varepsilon, a+\varepsilon]\cap I$.
	\minisec{Beweis}
	Es gilt $\varphi'(x)= f\big(x, \varphi(x)\big), \psi'(x)= f\big(x,\psi(x)\big)$, also 
	\begin{align*}
		\int_{a} ^{x} \! \enbrace{ f\big(t,\varphi(t)\big) - f\big(t,\psi(t)\big)}  \, \mathrm{d}t &= \int_{a} ^{x} \! \enbrace{\varphi'(t) - \psi'(t)}  \, \mathrm{d}t = \varphi(x)-\psi(x) 
	\end{align*}
	$f$ lokal Lipschitz $\Rightarrow \exists L >0$ und $\delta >0$ mit 
	\begin{align*}
		\norm[2]{f\big(t, \varphi(t) \big) - f\big(t, \psi(t)\big)} \le L \cdot \norm[2]{ \varphi(t) - \psi(t) }  \quad, \enspace t \in [a- \delta , a+ \delta ]\cap I
	\end{align*} 
	Setze $\varepsilon := \min \set{\delta , \frac{1}{2 L} } $. Für $x \in  [a- \delta , a+ \delta ] \cap I$ gilt nun
	\begin{align*}
		\norm[2]{\varphi(x)- \psi(x)} = \norm[2]{\int_{a} ^{x} \! \enbrace{f\big(t, \varphi(t)\big) - f\big(t, \psi(t)\big)}  \, \mathrm{d}t} &\le \int_{a} ^{x} \! 
		\norm[2]{f\big(t, \varphi(t)\big) - f\big(t, \psi(t)\big)}  \, \mathrm{d}t \\
		&\le L \cdot \int_{a} ^{x} \! \norm[2]{\varphi(t)- \psi(t)}  \, \mathrm{d}t 
	\end{align*}
	Für $x \in [a- \delta , a+ \delta ] \cap I$ setze $M(x) := \sup \set[{\norm[2]{\varphi(t)- \psi(t)}}]{t \in I, \abs{t-a}\le \abs{x-a}  } < \infty  $.
	Für $x, \xi \in [a-\delta , a+\delta ] \cap I$ mit $\abs{\xi-a}\le \abs{x-a}  $ gilt dann 
	\begin{align*}
		\norm[2]{\varphi(\xi) - \psi(\xi)} \le L \cdot \int_{a} ^{\xi} \!  \norm[2]{\varphi(t)- \psi(t)}  \, \mathrm{d}t \le L \cdot \abs{\xi -a}\cdot M(x) 
		\le L \cdot \abs{x-a} \cdot M(x) 
	\end{align*}
	also $M(x) \le L \cdot \abs{x-a} M(x) $. Für $\abs{x-a} < \varepsilon\le \frac{1}{2 \cdot L} $ ergibt sich 
	\[
		0 \le M(x) \le L \cdot \varepsilon \cdot M(x) \le \frac{1}{2} M(x)
	\]
	also muss $M(x)=0$ gelten und $\varphi(x)= \psi(x)$ für $x \in [a-\varepsilon, a + \varepsilon] \cap I$.
	\item[2. Behauptung] $\varphi(x) = \psi(x)$ für $x \in [x_0,\infty) \cap I$.
	\minisec{Beweis}
	Setze $x_1 := \sup \set[\xi \in I]{\varphi|_{[x_0, \xi]} = \psi|_{[x_0, \xi]}} $, \marginnote{Lücke schließen mit $\varphi, \psi$ stetig}
	dann gilt $x_1 > x_0$. Falls $x_1 \not\in I$, ist die Behauptung bewiesen. Falls
	$x_1 \in I$ und $[x_0, x_1]= [x_0, \infty) \cap I$, so ist die Behauptung auch wieder bewiesen. Sei nun $x_1 \in I$ und $[x_0, x_1] \not= [x_0, \infty) \cap I$. Dann 
	existiert $\eta >0$ mit $[x_0, x_1+ \eta] \subset I$. Nach der 1. Behauptung existiert $\varepsilon >0$ sodass $\varphi(x)=\psi(x)$ für 
	$x \in [x_1 -\varepsilon, x_1 + \varepsilon	] \cap I$. Setze $\zeta := \min \set{\varepsilon, \eta} $, dann $\varphi|_{[x_0, x_1 + \zeta]} = \psi|_{[x_0, x_1+ \zeta]}$
	\light
	\item[3. Behauptung] $\varphi(x) = \psi(x)$ für $x \in (- \infty, x_0] \cap I$
	\minisec{Beweis}
	Analog zu Behauptung 2.
\end{description}
$\Rightarrow \varphi(x)= \psi(x)$ für $x \in (- \infty, \infty) \cap I$. \bewende 
% subsection 16 (end)

\subsection[Satz von Picard-Lindelöf (Existenz einer Lösung)]{Satz (Picard-Lindelöf)} % (fold)
\label{sub:17}
Sei $G \subset \mathds{R} \times \mathds{R}^n$ offen, $f : G \to \mathds{R}^n$ stetig und lokal Lipschitz, $(x_0, y_0) \in G$. Dann existiert $\varepsilon >0$, sodass 
\[
	y' = f(x,y) \tag{$\star$}
\]
auf $[x_0 - \varepsilon, x_0 + \varepsilon]$ eine Lösung $\varphi$ mit der Anfangsbedingung $\varphi(x_0)= y_0$ hat. Dabei ist 
$\varphi(x)= \lim_{ k \to \infty} \varphi_k(x)$, wo $\varphi_0 \equiv y_0$ und
\[
	\varphi_{k+1} (x) = y_0 + \int_{x_0} ^{x} \! f\big(t, \varphi_k(t)\big)  \, \mathrm{d}t
\]
\minisec{Beweis}
\begin{description}
	\item[1. Schritt:] Es gibt $r>0$ sodass $K:= [x_0 -r, x_0 +r] \times \overline{B}_{\norm[2]{.} }(y_0,r) \subset G$ und $f|_{K}$ Lipschitz ist mit der Konstanten $L$. 
	$K$ kompakt, $f$ stetig $\Rightarrow M := \sup \set[{\norm[2]{f(x,y)} }]{(x,y) \in K} < \infty $. Setze $\varepsilon := \min \set{r, \frac{r}{M} }>0 $ (dürfen $M>0$
	annehmen). 
	\item[2. Schritt] $\varphi : [x_0 -\varepsilon, x_0 + \varepsilon] \to \mathds{R}^n$ erfüllt $(\star)$ mit Anfangsbedingung $\varphi(x_0)=y_0$ genau dann, wenn 
	\[
		\varphi(x) = y_0 + \int_{x_0} ^{x} \! f\big(t, \varphi(t)\big)  \, \mathrm{d}t \tag{$\star\star$}
	\]
	gilt für alle $x \in  [x_0 - \varepsilon, x_0+ \varepsilon]$ (Fundamentalsatz). Wir definieren induktiv Funktionen 
	$\varphi_k : [x_0 - \varepsilon, x_0 + \varepsilon] \to \mathds{R}^n$ durch $\varphi_0 (x)= y_0$, $x \in [x_0 - \varepsilon, x_0 + \varepsilon]$,
	\[
		\varphi_{k+1} := y_0 + \int_{x_0} ^{x} \! f\big(t, \varphi_k(t)\big)  \, \mathrm{d}t \quad , \qquad x \in[x_0-\varepsilon, x_0 + \varepsilon]\tag{$\star\star\star$}
	\]
	Zu zeigen: $\varphi_k$ ist wohldefiniert für $k \in \mathds{N}$, $\varphi_k$ konvergiert gleichmäßig auf $[x_0 -\varepsilon, x_0 +\varepsilon]$ gegen eine Lösung 
	$\varphi$.
	\item[3. Schritt] $(\star\star\star)$ ist wohldefiniert, falls $\norm[2]{\varphi_k(x)-y_0} \le r $ für $x \in [x_0 -\varepsilon, x_0 +\varepsilon]$. $k=0$ trivial.
	Falls die Aussage für $k$ gilt, so folgt 
	\begin{align*}
		\norm[2]{\varphi_{k+1} -y_0} \le \int_{x_0} ^{x} \! \norm[2]{f\big(t, \varphi_k(t)\big)}  \, \mathrm{d}t \le   M \cdot \abs{x-x_0} \le M \cdot \varepsilon
		\le r \qquad x \in [x_0 -\varepsilon, x_0 +\varepsilon]
	\end{align*}
	$\Rightarrow $ $(\star\star\star)$ ist wohldefiniert für alle $k \in \mathds{N}$
	\item[4. Schritt] Wir zeigen 
	\begin{align*}
		\norm[2]{\varphi_{k+1}(x)- \varphi_k(x)} \le M \cdot L^k \cdot \frac{\abs{x-x_0}^{k+1} }{(k+1)!} 
	\end{align*}
	für $k \in \mathds{N}$, $x \in [x_0 -\varepsilon, x_0 +\varepsilon]$. $k=0$: 
	\[
		\norm[2]{\varphi_1(x) - \varphi_0(x)} = \norm[2]{\int_{x_0} ^{x} \! f(t,y_0)  \, \mathrm{d}t}  \le \int_{x_0} ^{x} \! \norm[2]{f(t,y_0)}  \, \mathrm{d}t 
		\le M \cdot \frac{\abs{x-x_0}^{0+1} }{(0+1)!}
	\]
	$k-1 \mapsto k$:
	\begin{align*}
		\norm[2]{\varphi_{k+1}(x) - \varphi_k(x)} &= \int_{x_0} ^{x} \! \norm[2]{\Big(f\big(t, \varphi_k(t)\big) - f\big(t, \varphi_{k-1}(t)\big)\Big)  \, \mathrm{d}t} \\
		&\le \int_{x_0} ^{x} \! \norm[2]{f\big(t,\varphi_k(t)\big) -f\big(t, \varphi_{k-1}(t)\big) }  \, \mathrm{d}t \\
		&\le  \int_{x_0} ^{x} \! L \cdot \norm[2]{\varphi_k(t) - \varphi_{k-1}(t)}  \, \mathrm{d}t \\
		&\le M \cdot L^k \int_{x_0} ^{x} \! \frac{\abs{t-x_0}^k}{k!}  \, \mathrm{d}t \\
		&= M \cdot L^k \frac{\abs{x-x_0}^{k+1} }{(k+1)!} 
	\end{align*}
	\item[5. Schritt] Die Reihe $\sum_{k=0}^{\infty} (\varphi_{k+1} - \varphi_{k})$ wird majorisiert durch 
	$\sum_{k=0}^{\infty} M \cdot \frac{L^k \cdot \varepsilon^{k+1}}{(k+1)!}  $
	\begin{align*}
		\enbrace{\norm[2]{\sum_{k=0}^{k_0} \varphi_{k+1}(x) - \varphi_{k}(x)} \le \sum_{k=0}^{k_0}  \norm[2]{\varphi_{k+1}(x) - \varphi_{k}(x)}  \le  \sum_{k=0}^{k_0}
		\frac{M \cdot L^k \cdot \abs{x-x_0}^{k+1} }{(k+1)!}  \le \sum_{k=0}^{k_0} \frac{M \cdot L^k \cdot \varepsilon^{k+1}}{(k+1)!} }
	\end{align*}
	und konvergiert daher gleichmäßig; Der Limes ist stetig auf $[x_0 -\varepsilon, x_0 +\varepsilon]$. Wir erhalten \marginnote{Teleskopsumme}
	\[
		\varphi := \lim_{ k \to \infty} \varphi_k =  \lim_{ k \to \infty} y_0 + \sum_{l=0}^{k} \varphi_{l+1}- \varphi_l = y_0 + \lim_{ k \to \infty} \sum_{l=0}^{k}
		\varphi_{l+1} - \varphi_l = y_0 + \sum_{k=0}^{\infty} \varphi_{k+1} - \varphi_{k}
	\]
	$\varphi : [x_0 -\varepsilon, x_0 +\varepsilon] \to \mathds{R}^n$ stetig. Für $k \in \mathds{N}$, $x \in [x_0 -\varepsilon, x_0 +\varepsilon]$ gilt 
	\begin{align*}
		\norm[2]{f(x, \varphi(x)) - f(x, \varphi_k(x))} \le L \cdot \norm[2]{\varphi (x) - \varphi_k(x)} 
	\end{align*}
	$\varphi_l \xrightarrow{k \to \infty} \varphi$ gleichmäßig
	$\Rightarrow $ $\Big(x \mapsto f\big(x, \varphi_k(x)\big)\Big) \xrightarrow{k \to \infty} \Big(x \mapsto f\big(x, \varphi(x)\big)\Big)$ 
	gleichmäßig \marginnote{warum?}
	\begin{align*}
		\varphi(x) \stackrel{k \to \infty}{\longleftarrow} \varphi_{k+1}(x) = y_0 + \int_{x_0} ^{x} \! f\big( t, \varphi_k(t)\big)  \, \mathrm{d}t 
		\xrightarrow{k \to \infty} y_0 + \int_{x_0} ^{x} \! f\big(t, \varphi(t)\big)  \, \mathrm{d}t
	\end{align*}
	$\Rightarrow (\star\star)$ \bewende
\end{description}
% subsection 17 (end)

\subsection[Beispiel: Anwendung von Picard-Lindelöf]{Beispiel} % (fold)
\label{sub:18}
Betrachte $y' = 2 \cdot x \cdot  y$, $G = \mathds{R}\times \mathds{R}$. Gesucht: Lösung $\varphi$ mit $\varphi(0)=y_0$. Setze $\varphi_0(x) := y_0$, $x \in \mathds{R}$.
$\varphi_{k+1}(x) := y_0 + \int_{0} ^{x} \! 2 t \varphi_k(t)  \, \mathrm{d}t$, $x \in \mathds{R}$. Dann
\begin{align*}
	\varphi_1(x) &= y_0 + 2 \int_{0} ^{x} \! t \cdot  y_0  \, \mathrm{d}t = y_0 + y_0 \cdot x^2 \\
	\varphi_2(x) &= y_0 + \int_{0} ^{x} \! 2 t (y_0 + y_0 \cdot x^2)  \, \mathrm{d}t = y_0 + y_0 \cdot x^2 + y_0 \cdot \frac{x^4}{2} 
\end{align*}
Induktion liefert dann: $\varphi(x) = y_0(1+ x^2 + \frac{x^4}{2!} + \frac{x^6}{3!} + \frac{x^8}{4!} + \ldots + \frac{x^{2k}}{k!}    )$ also 
$\varphi_k (x) \xrightarrow{k \to \infty} y_0 \cdot e^{x^2}$
% subsection 18 (end)

\subsection[Corollar: Folgerungen aus dem Eindeutigkeits- und  Existenzsatz]{Corollar} % (fold)
\label{sub:19}
Sei $G \subset \mathds{R} \times \mathds{R}^n$ offen, $f : G \to \mathds{R}$ stetig und lokal Lipschitz.
\begin{enumerate}[(i)]
	\item Seien $\varphi, \psi : I \to \mathds{R}$ zwei Lösungen von 
	\[
		x^{(n)}= f\big(x, y, y', \ldots , y^{(n-1)}\big) \tag{$\star$}
	\]
	über dem Intervall $I$. Falls für $x_0 \in I$ gilt:
	\[
		\varphi(x_0)= \psi(x_0), \ldots , \varphi^{(n-1)}(x_0) = \psi^{(n-1)}(x_0)
	\]
	so ist $\varphi=\psi$ über $I$.
	\item Zu $(x_0, y_0, y_1, \ldots , y_{n-1}) \in G$ existiert $\varepsilon>0$ und eine Lösung $\varphi : [x_0 -\varepsilon, x_0 +\varepsilon] \to \mathds{R}$ vom 
	($\star$) mit 
	\[
		\varphi(x_0) = y_0, \enspace\varphi'(x_0)=y_1, \enspace\ldots , \enspace \varphi^{(n-1)}(x_0) = y_{n-1} 
	\]
	\minisec{Beweis}
	Bemerkung \ref{sub:13} (\ref{enum:13.2}), Satz \ref{sub:16}, Satz \ref{sub:17} \bewende
\end{enumerate}
% subsection corollar (end)

\subsection[Beispiel: Winkelfunktionen als Lösungen]{Beispiel} % (fold)
\label{sub:110}
$y''=-y$ ist eine Differentialgleichung 2. Ordnung. Lösungen sind $\sin(x), \cos(x), 2 \sin(x)+ \cos(x)$. Allgemeiner gilt: Für Konstanten $c_1, c_2 \in \mathds{R}$ betrachte 
\[
	\varphi(x) = c_1 \sin(x)+ c_2 \cos(x)
\]
Dann gilt $\varphi$ löst $y''=-y$ und $\varphi(0)=c_2$ und $\varphi'(0)=c_1$. Mit Eindeutigkeit folgt, dass alle diese Funktionen Lösungen sind.
% subsection 110 (end)
% section 1 (end)
\newpage
\section{Einige Lösungsmethoden} % (fold)
\label{sec:2}
\subsection{Trennung der Variablen} % (fold)
\label{sub:21}
Betrachte die Differentialgleichung
\[
	y' = f(x) \cdot g(y) \tag{$\star$}
\]
An dieser Stelle benehmen wir uns mal wie Physiker:
\[
	\diffd{y}{x} = f(x) \cdot g(y) \Longleftrightarrow \frac{1}{g(y)} \mathrm{d}y = f(x) \mathrm{d}x \Longleftrightarrow \int \! \frac{1}{g(y)} \mathrm{d}y = \int \! f(x) \mathrm{d}x
\]
\bet{Satz:} Seien $I,J \in \mathds{R}$ Intervalle $f:  I \to \mathds{R},\, g : J \to \mathds{R}^*$ stetig. Für $(x_0, y_0) \in I \times J$ definiere 
$F : I \to \mathds{R},\, G : J \to \mathds{R}$ durch
\[
	F(x) = \int_{x_0} ^{x} \! f(t)  \, \mathrm{d}t \qquad \qquad G(y) = \int_{y_0} ^{y} \! \frac{1}{g(s)}  \, \mathrm{d}s 
\]
Sei $I' \subset I$ ein Intervall mit $x_0 \in I'$ und $F(I') \subset G(J)$. Dann besitzt 
\[
	y' =  f(x) \cdot g(y) \tag{$\star$}
\]
eine eindeutige Lösung $\varphi : I' \to J$ mit $\varphi(x_0)= y_0$. Diese Lösung erfüllt: 
\[
	G\big(\varphi(x)\big) = F(x) \qquad \text{ für } x \in I' \tag{$\star\star$}
\]
\minisec{Beweis}
\begin{description}
	\item[Schritt 1:] Sei $\varphi : I' \to \mathds{R}$ eine Lösung von $(\star)$ mit $\varphi(x_0)=y_0$. Wir zeigen, dass $\varphi$ ($\star\star$) erfüllt:
	Wir haben: $\varphi'(x) = f(x) g\big(\varphi(x)\big)$, $x \in I'$. Also 
	\begin{align*}
		F(x) = \int_{x_0} ^{x} \! f(t)  \, \mathrm{d}t = \int_{x_0} ^{x} \! \frac{\varphi'(t)}{g\big( \varphi(t)\big)}  \, \mathrm{d}t 
		= \int_{\varphi(x_0)} ^{\varphi(x)} \! \frac{ \mathrm{d}s}{g(s)} = \int_{y_0} ^{\varphi(x)} \! \frac{ \mathrm{d}s}{g(s)} = G \big( \varphi(x)\big)   
	\end{align*} 
	\item[Schritt 2 (Eindeutigkeit):] $G'(y)= \frac{1}{g(y)} \not= 0 $ ist stetig $\Rightarrow $ $G$ ist stetig differenzierbar und streng monoton. 
	Daraus folgt: $G$ besitzt eine stetig
	differenzierbare Umkehrfunktion $H : G(J) \to \mathds{R}$. Aus $(\star\star)$ folgt
	\[
		\varphi(x) = H\Big(G\big(\varphi(x)\big)\Big) = H\big(F(x)\big) \qquad x \in I'
	\]
	$\Rightarrow \varphi$ eindeutig bestimmt durch $F$ und $H$.
	\item[Schritt 3 (Existenz):] Betrachte $H$ aus 2. und definiere $\varphi : I'  \to \mathds{R}$ durch $\varphi(x) = H \big( F(x)\big)$, $x \in I'$. $\varphi$ ist
	wohldefiniert, da $F(I') \subset G(J)$. $H,F$ stetig differenzierbar, also auch $\varphi$. Anfangsbedingung:
	\[
		\varphi(x_0) = H \big(F(x_0) \big) = H(0) = y_0
	\]
	Weiter gilt $G \big(\varphi(x) \big) = F(x)$
	\[
		\Rightarrow f(x)= F'(x) = (G \circ \varphi)' (x)= G'\big(\varphi(x)\big) \cdot \varphi'(x) = \frac{1}{g\big(\varphi(x)\big)} \cdot \varphi'(x) 
	\]
	Daraus folgt $\varphi'(x)= f(x) \cdot g\big(\varphi(x)\big)$. Also löst $\varphi$ $(\star)$. \bewende
\end{description}
% subsection 21 (end)

\subsection[Beispiel: Anwendung von \ref{sub:21} auf $y' =y^2$]{Beispiel} % (fold)
\label{sub:22}
\[
	y' = y^2 = f(x) \cdot g(y) \tag{$\star$}
\]
Dann $f(x)= 1, g(y)= y^2$. Anfangsbedingung $\varphi(0)=y_0 \in \mathds{R}$
\begin{description}
	\item[1. $y_0=0$:] Aus $y_0=0$ folgt, dass $\varphi(x) \equiv 0$ löst $(\star)$
	\item[2. $y_0 >0$:] Aus $y_0 >0$ folgt: Sei $\varphi: I \to \mathds{R}$ eine Lösung von $(\star)$ mit $\varphi(0)= y_0 > 0$. Falls es $x \in I$ gibt 
	mit $\varphi(x)<0$, dann
	existiert nach dem Zwischenwertsatz ein $\overline{x} \in I $ mit $\varphi(\overline{x} )=0$. Dies ist ein Widerspruch zur Eindeutigkeit der Lösung \light
	
	$\Rightarrow \varphi(x)>0$ für alle $x \in I$. $f : I=\mathds{R} \to \mathds{R}$, $g : J=R \to \mathds{R}_+^*$. $f(x)=1, g(x)=x^2$. Setze nun:
	$F : I \to \mathds{R}$, $G : J  \to \mathds{R}$
	\begin{align*}
		F(x) &= \int_{0} ^{x} \! f(t)  \, \mathrm{d}t = x \\
		G(y) &= \int_{y_0} ^{y} \! \frac{1}{g(s)}  \, \mathrm{d}s = \int_{y_0} ^{y} \! \frac{1}{s^2}  \, \mathrm{d}s = \frac{1}{y_0} - \frac{1}{y}   
	\end{align*}
	Für $I'=(- \infty, \frac{1}{y_0} )$ gilt $F(I')= (-\infty, \frac{1}{y_0} )$ und somit $F(I') \subset G(J)$. Nach \ref{sub:21} besitzt $(\star)$ eine eindeutige Lösung
	$\varphi : I' \to \mathds{R}_+^*$ mit $\frac{1}{y_0}- \frac{1}{\varphi(x)}= x  $ $\leadsto \varphi(x)= \frac{y_0}{1- x \cdot y_0} $
	\item[3. $y_0<0$:] Analog zu 2. $I'= (\frac{1}{y_0}, \infty )$ $\varphi(x)= \frac{y_0}{1- x \cdot y_0} $ 
\end{description}
% subsection 22 (end)

\subsection{Satz (homogene, lineare Differentialgleichung)} % (fold)
\label{sub:23}
Sei $I$ ein Intervall, $a : I \to \mathds{R}$ stetig, $x_0 \in I, y_0 \in \mathds{R}$. Dann besitzt \index{Differentialgleichung!homogene,lineare}
\[
	y' = a(x) y \tag{$\star$}
\]
genau eine Lösung $\varphi : I \to \mathds{R}$ mit $\varphi(x_0)= y_0$ gegeben durch:
\[
	\varphi(x) = y_0 \cdot e^{\int_{x_0} ^{x} \! a(t)  \, \mathrm{d}t} \tag{$\star\star$}
\]
\minisec{Beweis}
$\varphi(x_0)= y_0$ \checkmark \\
$\varphi$ wie in $(\star \star)$ ist eine Lösung von $(\star)$ (Nachrechnen). Da $(x,y) \mapsto a(x) \cdot y$ Lipschitz $\Rightarrow$ Lösung ist eindeutig (\ref{sub:16}) \bewende
% subsection 23 (end)

\subsection[Beispiel: homogene, lineare DGL mit konstanten Koeffizienten]{Beispiel} % (fold)
\label{sub:24}
\[
	y' = c \cdot y \tag{$\star$}
\]
hat Lösungen der Form $\varphi(x)= y_0 e^{c \abs{x-x_0} }$ mit $\varphi(x_0)=y_0$. $(\star)$ ist eine homogene lineare Differentialgleichung mit konstanten 
Koeffizienten\index{Differentialgleichung!homogene,lineare!mit konstanten Koeffizienten}. 
% subsection 24 (end)

\subsection{Satz (inhomogene lineare DGL; Variation der Konstanten)} % (fold)
\label{sub:25}
Sei $I \subset \mathds{R}$ ein Intervall, $a,b : I \to \mathds{R}$ stetig, $x_0 \in I$, $y_0 \in \mathds{R}$. Dann besitzt 
\[
	y' = a(x)\cdot y + b(x) \tag{$\star$}
\]
\emph{genau eine} Lösung $\psi : I \to \mathds{R}$ mit $\psi(x_0)=y_0$, nämlich
\[
	\psi(x) = \underbrace{y_0 \cdot \varphi(x)}_{\text{homogene Lösung}}+ \underbrace{\varphi(x) \cdot \int_{x_0} ^{x} \! \varphi(t) ^{-1} b(t)  \,
	 \mathrm{d}t}_{\text{inhomogene Lösung}} = \varphi(x) \underbrace{\enbrace{y_0 + \int_{x_0} ^{x} \! 
	\varphi(t) ^{-1} b(t)  \, \mathrm{d}t}}_{\text{Variation der Konstanten}} 
\]
wobei $\varphi(x)= e^{\int_{x_0} ^{x} \! a(t)  \, \mathrm{d}t}$ eine Lösung der homogenen linearen DGL $y'=a(x)y$ mit $\varphi(x_0)= 1$ ist.
\minisec{Beweis}
$\psi$ löst $(\star)$: einsetzen. 
$f(x,y)= a(x)y + b(x)$ ist lokal Lipschitz $\Rightarrow $ Eindeutigkeit. \bewende
% subsection 25 (end)

\subsection[Bemerkung: Variation der Konstanten bei System von linearen DGL]{Bemerkung} % (fold)
\label{sub:26}
Der Ansatz lässt sich (zum Teil) auf Systeme von linearen Differentialgleichungen der Form
\[
	y' = Ay+b \qquad \quad y,b : I \to \mathds{R}^n, A : I \to M_n(\mathds{R}) \tag{$\star$}
\]
übertragen. $(\star)$ schreibt sich auch als
\[
	\begin{pmatrix}
		y_1' \\ \vdots \\ y_n'
	\end{pmatrix} = \begin{pmatrix}
		a_{11}\cdot y_1 + \ldots + a_{1n}\cdot y_n \\
		\vdots \\
		a_{n1}\cdot y_1 + \ldots + a_{nn}\cdot y_n
	\end{pmatrix} + \begin{pmatrix}
		b_1 \\ \vdots \\ b_n
	\end{pmatrix}
\]
Es gilt: 
\[
	\set{\text{Lösungen der inhomogenen linearen DGL}} = \psi_0 + \set{\text{Lösungen der homogenen linearen DGL}}  
\]
wobei $\psi_0$ (irgend-)eine Lösung der inhomogenen linearen DGL ist. Die Lösungen der homogenen linearen DGL bilden einen Vektorraum.
% subsection 26 (end)

\subsection{Homogene Differentialgleichungen} % (fold)
\label{sub:27}
Betrachte
\[
	y' = f \enbrace{\frac{y}{x} } \tag{$\star$}
\]
Mit $z= \frac{y}{x} $ ergibt sich $y= x \cdot z$, also $y' = x \cdot z' + z$ und damit $x \cdot z' +z = f(z)$, also
\[
	z'=  \frac{1}{x} \big(f(z)-z\big) \tag{$\star\star$}
\]
$(\star\star)$ lässt sich mit Trennung der Variablen behandeln (siehe \ref{sub:21}). 
% subsection 27 (end)

\subsection[Satz: Lösung von homogenen Differentialgleichungen]{Satz} % (fold)
\label{sub:27.2}
Seien $I \subset \mathds{R}^*$, $J \subset \mathds{R}$ Intervalle, $f : J \to \mathds{R}$ stetig und 
$G := \set[(x,y) \in \mathds{R}^* \times \mathds{R}]{\frac{y}{x} \in J } $. Sei $(x_0,y_0) \in G$ ein Punkt mit $x_0 \in I$. Dann gilt: Eine Funktion 
$\varphi : I \to \mathds{R}$ ist Lösung von $y'= f(\frac{y}{x} ) \enspace(\star)$ genau dann, wenn $\psi : I  \to \mathds{R}$, $\psi(x) := \frac{\varphi(x)}{x} $ Lösung ist 
von $z'= \frac{1}{x} \big(f(z)-z\big) \enspace (\star\star)$ mit $\psi(x_0)= \frac{y_0}{x_0} $.
\minisec{Beweis}
Nachrechnen \bewende
% subsection 27 (end)

\subsection[Beispiel zur Lösung einer homogenen DGL]{Beispiel} % (fold)
\label{sub:28}
Betrachte 
\[
	y' = 1+ \frac{y}{x} + \enbrace{\frac{y}{x} }^2 
\]
auf $\mathds{R}^* \times \mathds{R}$ mit Anfangsbedingung $(x_0,y_0)$. Die Substitution $z=\frac{y}{x} $ liefert 
\[
	z'= \underbrace{\frac{1}{x}}_{=:f(x)} \underbrace{(1+z^2)}_{=:g(z)} 
\]
mit Anfangsbedingung $(x_0, \frac{y_0}{x_0} )$. Trennung der Variablen (Satz \ref{sub:21}) liefert:
\begin{align*}
	\arctan (\psi(x)) - \arctan \enbrace{\frac{y_0}{x_0} }  = \int_{\frac{y_0}{x_0} } ^{\psi(x)} \! \frac{1}{1+s^2}  \, \mathrm{d}s= G(\psi(x)) = F(x) = \int_{x_0} ^{x} \! \frac{1}{t}  \, \mathrm{d}t  = \ln x - \ln x_0 = \ln \frac{x}{x_0} 
\end{align*}
$\Rightarrow \psi(x) = \tan \enbrace{\ln \frac{x}{x_0} + \arctan \frac{y_0}{x_0} } $ und $\varphi(x)= x \cdot \psi(x)$ löst das Anfangswertproblem 
$y'= 1+ \frac{y}{x} + \enbrace{\frac{y}{x} }^2, (x_0, y_0) $.
% subsection 28 (end)

\subsection[Beispiel aus der Physik (?)]{Beispiel} % (fold)
\label{sub:29}
Sei $U : \mathds{R}^3 \to \mathds{R}$ zweimal stetig differenzierbar. Betrachte 
\[
	\ddot x = - \grad U(x) \tag{$\star$}
\]
d.h. wir suchen $x: I \to \mathds{R}^3$, $x(t)= \begin{pmatrix}
	x_1(t) \\ x_2(t) \\ x_3(t)
\end{pmatrix}$ mit $\ddot x(t)= - \grad U\big(x(t)\big)$, $t \in I$, also 
\[
	\begin{pmatrix}
		\ddot x_1(t) \\ \ddot x_2(t) \\ \ddot x_3(t)
	\end{pmatrix} = - \begin{pmatrix}
		\diff{U}{x_1} \big( x_1(t), x_2(t), x_3(t)\big) \\
		\diff{U}{x_2} \big( x_1(t), x_2(t), x_3(t)\big) \\
		\diff{U}{x_3} \big( x_1(t), x_2(t), x_3(t)\big) 
	\end{pmatrix}
\]
Wir nehmen an, dass das Potential $U$ in $0 \in \mathds{R}^3$ ein lokales Maximum hat und interessieren uns für Lösungen $x$ von ($\star$) in einer kleinen Umgebung von $0$.


$x(t)\equiv 0$ ist Lösung , d.h. das System ist im Minimum im Gleichgewicht. Wir nehmen weiter an, dass 
\[
	A := ( \hess U)(0) = \enbrace{ \frac{\partial^2 U}{\partial x \partial x_j} (0)}_{i,j} \in M_3(\mathds{R})  
\]
positiv definit ist, d.h. es existiert $S \in M_3(\mathds{R})$ orthogonal mit 
\[
	B := S ^{-1} A S = \begin{pmatrix}
		\lambda_1 & & 0 \\
		& \lambda_2 & \\
		0 & & \lambda_3
	\end{pmatrix} \qquad \quad , \lambda_1, \lambda_2, \lambda_3 >0
\]
Ana. II $\leadsto$
\[
	U(\xi)= U(0 + \xi) = U(0) + \frac{1}{2} \sprod{\xi}{A \xi} + \beta(\xi)  
\]
für ein $\beta : \mathds{R}^3 \to \mathds{R}$ mit $\lim_{ \xi \to 0} \frac{\beta(\xi)}{\norm[2]{\xi} } = 0 $. Wir approximieren nun $U(\xi)$ durch 
\[
	\tilde U (\xi) := U(0)+ \frac{1}{2} \sprod{\xi}{A \cdot \xi} 
\]
Wie in Ana. II 8.6 (i) erhält man 
\[
	\grad\enbrace{\xi \mapsto \sprod{\xi}{A \cdot \xi} }  = 2 A \cdot x
\]
Also $\grad \tilde U (x) = A \cdot x$
\[
	\leadsto \ddot x = - A x
\]
$y := S ^{-1} x$ $\Rightarrow \ddot y = S ^{-1} \ddot x  = S ^{-1} (-A)x =- S ^{-1} A S S ^{-1} x = - B S ^{-1}  x = -B y$. Also folgt:
\[
	\ddot y = - B y = - \begin{pmatrix}
		\lambda_1 & & \\
		& \lambda_2 & \\
		& & \lambda_3
	\end{pmatrix}y, \qquad \quad \begin{pmatrix}
		\ddot y_1 \\
		\ddot y_2 \\
		\ddot y_3
	\end{pmatrix} = \begin{pmatrix}
		-\lambda_1 y \\
		-\lambda_2 y \\
		- \lambda_3 y
	\end{pmatrix}
\]
mit den Lösungen $y_i(t) = \alpha_i \cdot  \cos \enbrace{\sqrt{\lambda_i}\cdot t  } + \beta_i \cdot \sin \enbrace{\sqrt{\lambda_i} \cdot t  } $, $\alpha_i, \beta_i \in \mathds{R}$
$i=1,2,3$
% subsection 29 (end)
% section 2 (end)
\newpage

\section{Treppenfunktionen: Die $L^1$-Halbnorm} % (fold)
\label{sec:3}

\subsection[Definition: Quader, Volumen, Treppenfunktion]{Definition} % (fold)
\label{sub:31}
\begin{enumerate}[a)]
	\item Ein \Index{Quader} $Q \subset \mathds{R}^n$ ist eine Teilmenge der Form $Q = I_1 \times I_2 \times \ldots \times I_n$, wo $I_j \subset \mathds{R}, j=1,\ldots ,n$
	ein beschränktes, nicht leeres Intervall ist. \\
	(Wir lassen offene, halboffene, abgeschlossene und Intervalle mit einem Punkt zu.)
	\item Das \Index{Volumen} eines Quaders wie in a) definieren wir als $v(Q) := \abs{I_1} \cdot \abs{I_2} \cdot \ldots \cdot \abs{I_n}$ \\
	($\abs{I_j} $ ist die Länge von $I_j$)
	\item $\varphi : \mathds{R}^n \to \mathds{R}$ heißt \Index{Treppenfunktion}, falls paarweise disjunkte Quader $Q_1, \ldots , Q_s \subset \mathds{R}^n$ existieren mit
	$\varphi|_{Q_j}$ konstant für $j=1, \ldots ,s$ und 
	\[
		\varphi\Big|_{\mathds{R}^n \setminus \bigcup_{j=1}^s Q_j} \equiv 0
	\]
	$\varphi$ lässt sich dann schreiben als
	\[
		\varphi = \sum_{j=1}^{s} c_j \cdot \chi_{Q_j}
	\]
	wo $\chi_{Q_j} : \mathds{R}^n \to \mathds{R}$ die charakteristische Funktion von $Q_j$ ist, \marginnote{auch mit $\mathds{1}_{Q_j}$ bezeichnet}
	\[
		\chi_{Q_j}(x) = \begin{cases}
			1, &\text{ falls }x \in Q_j\\
			0, &\text{ falls }x \not\in Q_j 
		\end{cases}
	\]
	Wir schreiben $\mathcal{T}_{\mathds{R}^n} := \set{\text{Treppenfunktionen auf }\mathds{R}^n} $.
	\item Für eine Treppenfunktion $\varphi$ wie in c) definieren wir das \bet{Integral}\index{Integral!einer Treppenfunktion}
	\[
		\int_{\mathds{R}^n} \! \varphi(x)  \, \mathrm{d}x  :=  \sum_{j=1}^{s} c_j \cdot v(Q_j)
	\]
	Wir schreiben auch $\int\! \varphi  \, \mathrm{d}x$.
\end{enumerate}
% subsection 31 (end)

\subsection[Bemerkung: Zerlegung von Quadern in disjunkte Quader]{Bemerkung} % (fold)
\label{sub:32}
Falls $Q_1, \ldots , Q_s \subset \mathds{R}^n$ Quader sind (nicht notwendig paarweise disjunkt), $c_1, \ldots , c_s \in \mathds{R}$, so ist 
$\varphi = \sum_{j=1}^{s} c_j \cdot \chi_{Q_j}$ eine Treppenfunktion: Es existieren paarweise disjunkte Quader $Q_1', \ldots , Q_{s'}' \subset \mathds{R}^n$ mit
$\varphi|_{Q_i'} \equiv c_i$ und $\bigcup Q_j = \bigcup Q_i'$ $\leadsto \varphi = \sum_{i=1}^{s'} c_i' \cdot \chi_{Q_i'}$
% subsection 32 (end)

\subsection[Proposition: Eigenschaften von Treppenfunktionen]{Proposition} % (fold)
\label{sub:33}
\begin{enumerate}[a)]
	\item $\mathcal{T}_{\mathds{R}^n}$ ist ein Vektorraum über $\mathds{R}$.
	\item $\int_{\mathds{R}^n} \! \varphi(x)  \, \mathrm{d}x$ hängt nicht von der Darstellung $\varphi = \sum_{k=1}^{s} c_k \cdot \chi_{Q_k}$ ab (wo 
	$Q_1, \ldots , Q_s \subset \mathds{R}^n$ Quader sind und $c_1, \ldots , c_s \in \mathds{R}$)
	\item Es gilt 
	\begin{enumerate}[(i)]
		\item $\int(\alpha \cdot \varphi + \beta \cdot \psi)  \, \mathrm{d}x  = \alpha \cdot  \int \! \varphi  \, \mathrm{d}x  + \beta \cdot \int\! \psi  \, \mathrm{d}x  $,
		$\varphi, \psi \in \mathcal{T}_{\mathds{R}^n}$, $\alpha, \beta \in \mathds{R}$. Mit anderen Worten $(\varphi \mapsto \int\! \varphi  \, \mathrm{d}x )$ ist lineares
		Funktional auf $\mathcal{T}_{\mathds{R}^n}$
		\item $\abs{\int\! \varphi  \, \mathrm{d}x } \le \int \! \abs{\varphi}  \, \mathrm{d}x   $
		\item $\varphi \le \psi \Longrightarrow \int \! \varphi  \, \mathrm{d}x  \le \int\! \psi  \, \mathrm{d}x  $
	\end{enumerate}
\end{enumerate}
\minisec{Beweis}
\begin{enumerate}[a)]
	\item 1. $Q, Q' \subset \mathds{R}^n$ Quader $\Rightarrow Q \cap Q'$ Quader \\
	2. $\varphi = \sum_{k=1}^{s} c_k \cdot \chi_{Q_k}$, $\psi = \sum_{l=1}^{s'}c_l' \cdot \chi_{Q_l'} \in \mathcal{T}_{\mathds{R}^n} $, $\alpha, \beta \in \mathds{R}$. Dann
	gilt o.E. $\cup Q_k  = \cup Q_l'$
	\[
		\alpha \cdot \varphi + \beta \cdot \psi = \sum_{\substack{j=1, \ldots ,s \\ l=1, \ldots ,s'}} (\alpha \cdot c_k + \beta \cdot c_l') \cdot \chi_{Q_k \cap Q_l'} \in
		\mathcal{T}_{\mathds{R}^n} 
	\]
	\item Für $n=1$: Analysis \RM{1}. \\
	Sei die Behauptung für $n-1$ bewiesen. Betrachte 
	\[
		\varphi = \sum_{k=1}^{s} c_k \cdot \chi_{Q_k} \in \mathcal{T}_{\mathds{R}^n} \tag{$\star$}
	\]
	Jeder Quader $Q_k$ lässt sich schreiben als $Q_k = Q_k ' \times Q_k'' $ mit Quadern $Q_k'  \subset \mathds{R}$ und $Q_k'' \subset \mathds{R}^{n-1}$ 
	$Q_k = (I_1) \times \enbrace{I_2 \times \ldots \times I_n} $. Es gilt für $(x,y) \in R \times R^{n-1}= \mathds{R}^n$. Dann gilt
	\[
		\chi_{Q_k} ((x,y)) = \chi_{Q_k'}(x) \cdot \chi_{Q_k''}(y)
	\]
	Für $y \in \mathds{R}^{n-1}$ definieren wir $\varphi_y : \mathds{R} \to \mathds{R}$ durch $\varphi(y) := \varphi(x,y)$ ($\varphi_y$ hängt nicht von der Darstellung
	in $(\star)$ ab). Dann gilt:
	\begin{align*}
		\varphi_y(x) &= \sum_{k=1}^{s} c_k \cdot \chi_{Q_k} (x,y) = \sum_{k=1}^{s} c_k \cdot \chi_{Q_k''}(y) \cdot \chi_{Q_k'}(x) \\
		\intertext{Also} \varphi_y &= \sum_{k=1}^{s} c_k \cdot \chi_{Q_k''} \cdot \chi_{Q_k'} \in \mathcal{T}_{\mathds{R}} \tag{$\star\star$}
	\end{align*}
	nach Bemerkung \ref{sub:32}. Aus Analysis I folgt $\int \varphi_y(x)  \, \mathrm{d}x $ ist wohldefiniert und hängt nicht von der Darstellung $(\star\star)$ ab:
	\[
		\Phi (y) := \int \! \varphi_y(x)  \, \mathrm{d}x  = \sum_{k=1}^{s} c_k \cdot  \chi_{Q_k''}(y) \cdot v(Q_k') \tag{$\star\star\star$}
	\]
	Dann ist $\Phi : \mathds{R}^{n-1} \to \mathds{R}$ definiert durch $y \mapsto \Phi(y)$, eine Treppenfunktion, denn nach $(\star\star\star)$
	\[
		\Phi = \sum_{k=1}^{s} c_k \cdot v(Q_k') \cdot \chi_{Q_k''}
	\]
	I.V. $\Rightarrow $ $\int \! \Phi(y)  \, \mathrm{d}y$ ist wohldefiniert und hängt nicht von der Darstellung ($\star\star\star$) ab; es gilt
	\[
		\int_{\mathds{R}^{n-1}} \! \Phi (y)  \, \mathrm{d}y = \sum_{k=1}^{s} c_k \cdot v(Q_k') \cdot v(Q_k'')
	\]
	also 
	\[
		\int_{\mathds{R}^{n-1}}\! \enbrace{\int \! \varphi_y(x)  \, \mathrm{d}x } \mathrm{d}y  = \sum_{k=1}^{s} c_k \cdot v(Q_k) \tag{$\star\star\star\,\star$}
	\]
	und die linke Seite hängt nicht von der Darstellung ($\star$) ab. 
	\item Übung, benutze ($\star\star\star\,\star$). \bewende
\end{enumerate}
% subsection 33 (end)

\subsection{Corollar (Fubini für Treppenfunktionen)} % (fold)
\label{sub:34}
Für $\varphi \in \mathcal{T}_{\mathds{R}^n}$ gilt mit $\mathds{R}^n = R^{p} \times \mathds{R}^{n-p} $
\[
	\int_{\mathds{R}^n} \varphi(z)  \, \mathrm{d}z = \int_{\mathds{R}^{n-p}} \enbrace{\int_{\mathds{R}^p} \varphi(x,y)  \, \mathrm{d}x } \mathrm{d}y
\]
\minisec{Beweis}
Benutze wiederholt ($\star\star\star\,\star$) aus \ref{sub:33} b) \bewende
% subsection 34 (end)

\subsection[Notation für den Funktionswert $\infty$]{Notation} % (fold)
\label{sub:35}
Wir versehen $(- \infty, \infty] = \mathds{R} \cup \set{\infty} $ mit Ordnung und Operationen von $\mathds{R}$ und 
\begin{multicols}{2}
	\begin{itemize}
		\item $c < \infty$, $\forall c \in \mathds{R}$,
		\item $\infty+c = c + \infty = \infty$, $c \in \mathds{R} \cup \set{\infty} $,
		\item $\infty \cdot c = c \cdot \infty = \infty$, $c \in \mathds{R}^*$,
		\item $\infty \cdot  0 = 0 \cdot \infty = 0$.
	\end{itemize}
\end{multicols}
% subsection 35 (end)

\subsection[Definition: Treppe]{Definition} % (fold)
\label{sub:36}
Sei $f : \mathds{R}^n \to (- \infty, \infty]$ eine Abbildung. Eine \Index{Treppe} über $f$ ist eine Reihe
\[
	\Phi = \sum_{k=0}^{\infty} c_k \cdot \chi_{Q_k}
\] 
mit folgenden Eigenschaften: 
\begin{enumerate}[(i)]
	\item $Q_k \subset \mathds{R}^n$ sind offene Quader, $c_k \in \mathds{R}_+$
	\item $\abs{f(x)} \le \Phi(x) := \sum_{k=0}^{\infty} c_k \cdot \chi_{Q_k}(x)  \in (- \infty, \infty] $, $x \in \mathds{R}^n$
\end{enumerate}
Wir setzen $I(\Phi) := \sum_{k=0}^{\infty} c_k \cdot v(Q_k) \in (- \infty, \infty]$ \marginnote{$I$ für "{}Inhalt"}
% subsection 36 (end)

\subsection[Definition: $L^1$-Halbnorm]{Definition} % (fold)
\label{sub:37}
Sei $f : \mathds{R}^n \to (-\infty, \infty]$. Wir setzen\index{L1-Halbnorm@$L^1$-Halbnorm} 
\[
	\norm[1]{f} :=  \inf \set[I(\Phi)]{\Phi \text{ Treppe über }f } 
\]
% subsection 37 (end)

\subsection[Bemerkungen zur $L^1$-Halbnorm]{Bemerkung} % (fold)
\label{sub:38}
\begin{enumerate}[1.]
	\item Eine Treppe über $f$ existiert für jedes $f$ (warum?), daher $\norm[1]{f} \in [0, \infty] $ für jedes $f$.
	\item Für $f,g : \mathds{R}^n \to (- \infty, \infty]$, $\lambda  \in \mathds{R}$ gilt:
	\begin{enumerate}[(i)]
		\item $\norm[1]{\lambda  \cdot f}  = \abs{\lambda } \cdot  \norm[1]{f}  $
		\item $\norm[1]{f+g} \le \norm[1]{f} + \norm[1]{g}   $
		\item $\abs{f} \le \abs{g} \Longrightarrow \norm[1]{f} \le \norm[1]{g}   $
	\end{enumerate}
	\uline{Beweis:} Übung.
	\item Es gilt \emph{nicht}: $\norm[1]{f} = 0 \Longrightarrow f \equiv 0 $ (Übung)\marginnote{für stetige $f$ allerdings schon}\\
	$\norm[1]{.} $ ist also \Index{Halbnorm} auf $\set[f : \mathds{R}^n \to \mathds{R}]{\norm[1]{f} < \infty } $ VR: Siehe \ref{sub:61} und \ref{sub:62}
\end{enumerate}
% subsection 38 (end)

\subsection[{Proposition: Dreiecksungleichung für $\norm[1]{.} $ gilt auch für unendliche Reihen}]{Proposition} % (fold)
\label{sub:39}
Für $f_0, f_1, \ldots  : \mathds{R}^n \to [0, \infty]$ definiert $x \mapsto \sum_{k=0}^{\infty} f_k(x) = \lim_{ K \to \infty} \sum_{k=0}^{K} f_k(x) $, $x \in \mathds{R}^n$
eine Abbildung $\sum_{k=0}^{\infty} f_k : \mathds{R}^n \to [0, \infty]$. Es gilt 
\[
	\norm[1]{\sum_{k=0}^{\infty} f_k} \le \sum_{k=0}^{\infty} \norm[1]{f_k}  
\]
\minisec{Beweis}
Sei $\varepsilon >0$. Für jedes $k \in \mathds{N}$ wähle eine Treppe $\Phi_k = \sum_{i=0}^{\infty} c_{k,i} \cdot \chi_{Q_{k,i}}$ über $f_k$ mit 
$I(\Phi_k) \le \norm[1]{f_k} + \frac{\varepsilon}{2^{k+1}} $. Wähle eine Bijektion $\gamma : \mathds{N} \to \mathds{N} \times \mathds{N}$ und setze 
\[
	\Phi := \sum_{l \in \mathds{N}} c_{\gamma(l)} \cdot \chi_{Q_{\gamma(l)}} = \sum_{\substack{k \in \mathds{N} \\ i \in \mathds{N}}} c_{k,i} \cdot \chi_{Q_{k,i}}
\]
Dann gilt 
\begin{itemize}
	\item $\Phi = \sum_{k=0}^{\infty} \Phi_k$, insbesondere ist $\Phi$ nicht abhängig von $\gamma$
	\item $\Phi$ ist Treppe über $\sum_{k \in \mathds{N}} f_k$
	\item \begin{align*}
		\norm[1]{\sum_{k=0}^{\infty} f_k} \le I(\Phi) = \sum_{l=0}^{\infty} c_{\gamma(l)} \cdot v(Q_{\gamma(l)}) &\stackrel{\star}{=} \sum_{k=0}^{\infty} I(\Phi_k) \\
		&\le \sum_{k=0}^{\infty} \enbrace{\norm[1]{f_k} + \frac{\varepsilon}{2^{k+1}}  } \\
		&= \sum_{k=0}^{\infty} \norm[1]{f_k} + \varepsilon  
	\end{align*}
	$\Rightarrow \norm[1]{\sum_{k=0}^{\infty} f_k} \le \sum_{k=0}^{\infty} \norm[1]{f_k}  $
\end{itemize}
Für $\star$ benutze: $(a_k)_{k \in \mathds{N}} \subset [0,\infty]$, $\beta : \mathds{N} \to \mathds{N}$ Bijektion, dann gilt\marginnote{gilt nur, da $a_k\ge 0$;\\ $\Rightarrow $ Riemannscher Umordnungssatz}
\[
	S_K = \sum_{k=0}^{K} a_k \xrightarrow{K \to \infty}  \sum_{k=0}^{\infty} a_k = \sum_{k=0}^{\infty} a_{\beta(k)} \xleftarrow{L \to \infty} \sum_{k=0}^{L} a_{\beta(k)}
	= V_L \bewende 
\]
% subsection 39 (end)

\subsection[{Lemma: Volumen eines abgeschlossenen Quaders $Q$ entspricht $\norm[1]{\chi_Q} $}]{Lemma} % (fold)
\label{sub:310}
Sei $Q \subset \mathds{R}^n$ ein abgeschlossener Quader. Dann gilt 
\[
	\norm[1]{\chi_Q} = v(Q) \enbrace{=: \int\! \chi_Q  \, \mathrm{d}x  }
\]
\minisec{Beweis}
Zu jedem $\varepsilon>0$ existiert ein offener Quader $Q' \subset \mathds{R}^n$ mit $Q \subset Q'$ und $v(Q') \le v(Q) + \varepsilon$. $\chi_{Q'}$ ist dann Treppe über 
$\chi_Q$. Es gilt
\[
	I(\chi_Q') = v(Q') \le v(Q)+ \varepsilon
\]
also $\norm[1]{\chi_Q} = \inf \set[I(\Phi)]{\Phi \text{ Treppe über }\chi_Q} \le v(Q) $. Sei nun $\Phi = \sum_k c_k \cdot \chi_{Q_k}$ eine Treppe über $\chi_Q$. Zu 
$\varepsilon>0$ und $x\in Q$ existieren wegen $\chi_Q (x) \le \Phi(x)$ ein $N(x) \in \mathds{N}$ mit 
\[
	1- \varepsilon = \chi_Q(x) - \varepsilon \le \sum_{k=0}^{N(x)} c_k \cdot \chi_{Q_k}(x)
\]
Die $Q_k$ sind offen, daher existiert eine Umgebung $U(x)$, so dass gilt: $x \in Q_k \Rightarrow U(x) \subset Q_k$ für $k=0,\ldots ,N(x)$. Für $y \in U(x)$ erhalten wir
\[
	1- \varepsilon \le \sum_{k=0}^{N(x)} c_k \cdot \chi_{Q_k}(x) \le \sum_{k=0}^{N(x)} c_k \cdot \chi_{Q_k}(y)  
\]
$(U(x))_{x \in Q}$ ist eine offene Überdeckung von $Q$. $Q$  ist kompakt, also gilt nach Heine-Borel
\[
	\exists x_1, \ldots , x_L \in Q : Q \subset \bigcup_{i=1}^L U(x)
\]
Setze $N := \max \set{N(x_1), \ldots , N(x_L)} $, dann gilt für $y \in Q$
\[
	\sum_{k=0}^{N} c_k \cdot \chi_{Q_k}(y) \ge 1- \varepsilon = (1- \varepsilon) \chi_{Q}(y) 
\]
also $\sum_{k=0}^{N} c_k \cdot \chi_{Q_k} \ge (1-\varepsilon) \cdot \chi_Q $. Mit Proposition \ref{sub:33} folgt dann
\begin{align*}
	I(\Phi) \ge \sum_{k=0}^{N} c_k \cdot v(Q_k) = \int \! \enbrace{\sum_{k=0}^{N} c_k \cdot \chi_{Q_k}(y) } \mathrm{d}y  \ge \int\! (1-\varepsilon) \chi_{Q}(y)  \, \mathrm{d}y = (1-\varepsilon) v(Q)
\end{align*}
Es folgt $I(\Phi) \ge v(Q)$ und damit $\norm[1]{\chi_Q} \ge v(Q)$ \bewende
% subsection 310 (end)

\subsection[{Lemma: Das Integral einer Treppenfunktion $\varphi$ ist gleich $\norm[1]{\varphi} $}]{Lemma} % (fold)
\label{sub:311}
Für jedes $\varphi \in \mathcal{T}_{\mathds{R}^n}$ gilt 
\[
	\norm[1]{\varphi} = \int\! \abs{\varphi}  \, \mathrm{d}x   
\]
\minisec{Beweis}
\begin{enumerate}[1.]
	\item Es existieren paarweise disjunkte Quader $Q_1, \ldots , Q_s, P_1, \ldots , P_r \subset \mathds{R}^n$ und $c_1, \ldots , c_s, d_1, \ldots , d_r \in \mathds{R}$
	mit 
	\begin{itemize}
		\item $\varphi = \sum_{k=1}^{s} c_k \cdot \chi_{Q_k} + \sum_{l=1}^{r} d_l \cdot \chi_{P_l}$
		\item $Q_k \subset \mathds{R}^n$ ist offen für $k=1, \ldots , s$
		\item $v(P_l) = 0$, $l=1,\ldots ,r$ (d.h. $P_l$ ist von der Form $I_{l,1} \times \ldots \times I_{l,n}$ \marginnote{(Übung)}
		 und mindestens ein $I_{l,j}$ besteht nur aus einem Punkt)
	\end{itemize}
	Es gilt dann $\abs{\varphi} = \sum_{k=1}^{s} \abs{c_k} \cdot \chi_{Q_k} + \sum_{l=1}^{r} \abs{d_l} \cdot \chi_{P_l}  $ (die Quader sind paarweise disjunkt). Sei nun
	$\varepsilon>0$. Zu jedem $P_l$ existiert ein offener Quader $P_l'$ mit $P_l \subset P_l'$ und $v(P_l') < \varepsilon$. Dann ist 
	\[
		\Phi := \sum_{k=0}^{s} \abs{c_k} \cdot \chi_{Q_k} + \sum_{l=1}^{r} \abs{d_l} \cdot \chi_{P_l'}   
	\]
	eine Treppe über $\abs{\varphi} $ also auch über $\varphi$. Es ist 
	\[
		I(\Phi) = \sum_{k=1}^{s} \abs{c_k} \cdot v(Q_k) + \sum_{l=1}^{r} \abs{d_l} \cdot v(P_l') \le \sum_{k=1}^{s} \abs{c_k} \cdot v(Q_k) + \varepsilon \cdot 
		\sum_{l=1}^{r} \abs{d_k}      
	\]
	also $\norm[1]{\varphi}  \le \sum_{k=1}^{s} \abs{c_k} \cdot v(Q_k) = \int \abs{\varphi}  \, \mathrm{d}x   $
	\item Sei nun $Q \subset \mathds{R}^n$ ein abgeschlossener Quader mit $\enbrace{\bigcup Q_k} \cup \enbrace{\bigcup P_l} \subset Q$, d.h. 
	$\varphi\big|_{R^n \setminus Q} \equiv 0$. Setze $M := \max \set{\abs{c_k}, \abs{d_l}  } = \max \abs{\varphi} \in \mathds{R} $. Setze 
	$\psi := M \cdot \chi_Q - \abs{\varphi}$ dann ist $0 <\psi = \abs{\psi} \in \mathcal{T}_{\mathds{R}^n}$. Nach Teil 1. gilt 
	\[
		\norm[1]{\psi} \le \int\! \abs{\psi}  \, \mathrm{d}x = \int \!\psi  \, \mathrm{d}x   
	\]
	Wir erhalten 
	\begin{align*}
		\int \! \abs{\varphi} = \int \!  \enbrace{M \cdot \chi_Q -\psi}  \, \mathrm{d}x  &= M \cdot \int\! \chi_Q  \, \mathrm{d}x - \int\! \psi  \, \mathrm{d}x  \\
		&= M \cdot \norm[1]{\chi_Q} - \int\! \psi  \, \mathrm{d}x \tag{\ref{sub:310}} \\
		&= \norm[1]{M \cdot \chi_Q} - \int\! \psi  \, \mathrm{d}x  \\
		&\le \norm[1]{M \cdot \chi_Q} - \norm[1]{\psi}   \\
		&=  \norm[1]{\abs{\varphi} + \psi } - \norm[1]{\psi} \\
		&\le \norm[1]{\abs{\varphi} }   \\
		&= \norm[1]{\varphi} \bewende 
	\end{align*}
\end{enumerate}
% subsection 311 (end)
% section 3 (end)
\newpage

\section{Das Lebesgue-Integral} % (fold)
\label{sec:4}

\subsection[Definition: Lebesgue-integrierbar]{Definition} % (fold)
\label{sub:41}
$f : \mathds{R}^n \to (- \infty, \infty]$ heißt \Index{Lebesgue integrierbar}, falls eine Folge $(\varphi_k)_{k \in \mathds{N}} \subset \mathcal{T}_{\mathds{R}^n}$ 
existiert mit
\[
	\norm[1]{f- \varphi_k} \xrightarrow{k \to \infty} 0 
\]
Wir setzen 
\[
	\int_{\mathds{R}^n} \! f(x)  \, \mathrm{d}x := \lim_{ k \to \infty} \int_{\mathds{R}^n}\! \varphi_k (x) \, \mathrm{d}x \tag{$\star$}
	\marginnote{siehe auch \ref{sub:31}}
\]
Wir schreiben auch $\int\! f  \, \mathrm{d}x $.
% subsection 41 (end)

\subsection[Bemerkung: Limes existiert und Wahl der Approximationsfolge irrelevant]{Bemerkung} % (fold)
\label{sub:42}
\begin{enumerate}[1)]
	\item Der Limes in $(\star)$ existiert, denn $\enbrace{\int\! \varphi_k(x)  \, \mathrm{d}x }_{k \in \mathds{N}} $ ist eine Cauchyfolge:
\begin{align*}
	\abs{\int\! \varphi_k  \, \mathrm{d}x - \int\!\varphi_l  \, \mathrm{d}x } \le \int\! \abs{\varphi_k- \varphi_l}  \, \mathrm{d}x = \norm[1]{\varphi_k - \varphi_l} 
	\le \norm[1]{\varphi_k -f} + \norm[1]{f- \varphi_l }  \le \varepsilon \marginnote{für geeignete $k,l$}    
\end{align*}
Ebenso: Der Limes in $(\star)$ hängt nicht von der Wahl der Approximationsfolge $(\varphi_k)_{k \in \mathds{N}}$ ab. (D.h. für eine weitere Folge 
$(\psi_k)_{k \in \mathds{N}} \subset \mathcal{T}_{\mathds{R}^n}$ mit $\norm[1]{\psi_k - f} \xrightarrow{k \to \infty} 0$, so ist 
$\lim_k\int\! \psi_k  \, \mathrm{d}x  = \int\! \varphi_k  \, \mathrm{d}x $  )
\item Sei $\varphi \in \mathcal{T}_{\mathds{R}^n}$. Dann ist $\varphi$ Lebesgue integrierbar und das Integral stimmt mit dem aus (\ref{sub:31}) überein.
\end{enumerate}
% subsection 42 (end)

\subsection[Proposition: Linearität und andere Eigenschaften des Lebesgue-Integrals]{Proposition} % (fold)
\label{sub:43}
Seien $f,g : \mathds{R}^n \to (-\infty, \infty]$ integrierbar, $\alpha, \beta \in \mathds{R}$. Dann gilt
\begin{enumerate}[a)]
	\item $\alpha \cdot f + \beta \cdot g$ ist integrierbar; es gilt 
	\[
		\int\! (\alpha \cdot f + \beta \cdot g)  \, \mathrm{d}x  = \alpha \cdot \int\! f  \, \mathrm{d}x + \beta \cdot \int\! g  \, \mathrm{d}x. 
	\]
	Insbesondere ist $\mathcal{L}^{1}(\mathds{R}^n) := \set[f : \mathds{R}^n \to \mathds{R}]{f \text{ Lebesgue integrierbar}} $ ein Vektorraum über $\mathds{R}$ und
	$\int : \mathcal{L}^{-1}(\mathds{R}^n) \to \mathds{R}$ ist linear.
	\item $\abs{f}$ ist integrierbar und es gilt:
	\[
		\abs{\int\! f  \, \mathrm{d}x } \le \int\! \abs{f}  \, \mathrm{d}x   = \norm[1]{f} 
	\]
	\item Ist $f \le g$, so folgt $\int\! f  \, \mathrm{d}x  \le \int\! g  \, \mathrm{d}x $ \hfill(Monotonie)
	\item Aus $g$ beschränkt folgt $f \cdot g$ ist integrierbar.
\end{enumerate}
\minisec{Beweis}
\begin{enumerate}[a)]
	\item Seien $(\varphi_k)_\mathds{N}$, $(\psi_k)_\mathds{N} \subset \mathcal{T}_{\mathds{R}^n}$ mit 
	$\norm[1]{f- \varphi_k}, \norm[1]{g - \psi_k} \xrightarrow{k \to \infty} 0  $. Dann ist $\alpha \cdot \varphi_k + \beta \cdot \psi_k \in \mathcal{T}_{\mathds{R}^n}$
	und 
	\begin{align*}
		\norm[1]{\alpha \cdot f + \beta \cdot g - \alpha \cdot \varphi_k - \beta \cdot \psi_K}  \le \abs{\alpha} \cdot \norm[1]{f- \varphi_k} + \abs{\beta} \cdot \norm[1]{g- \psi_k} \xrightarrow{k \to \infty} 0    
	\end{align*}
	\item benutze $\big| \abs{f}- \abs{\varphi_k} \big|   \le \abs{f- \varphi_k} $, also $\norm[1]{\abs{f}- \abs{\varphi_k}} \le \norm[1]{f- \varphi_k}$, sowie 
	$\norm[1]{f} - \norm[1]{f- \varphi_k} \le \norm[1]{\varphi_k} \le \norm[1]{f} + \norm[1]{f- \varphi_k}$
	\item Übung
	\item Übung \bewende
\end{enumerate}
% subsection 43 (end)

\subsection[Bemerkung: Zerlegung einer Funktion in positiven und negativen Teil]{Bemerkung} % (fold)
\label{sub:44}
\begin{wrapfigure}{R}[2cm]{8cm}
	\begin{tikzpicture}[domain=-3:5, xscale=0.5, yscale=0.6]
			%\draw[help lines] (-3,-2) grid (5,2);
			\draw[thick, ->, color=Snow4] (-3.15,0) -- (5.5,0);
			%\draw[thick, ->] (0,-2) -- (0,2);
			\draw[color=DodgerBlue1, thick] plot (\x,{sin(\x r)+0.3*\x }) node[right]{$f=f_+ - f_-$};
			\draw[color=Firebrick3, ultra thick, dotted] (-3,0) node[left]{$f_+$}-- (0,0) plot[domain=0:5] (\x,{sin(\x r)+0.3*\x});
			\draw[color=SeaGreen4, ultra thick, dotted]  plot[domain=-3:0] (\x,{abs(sin(\x r)+0.3*\x)}) (0,0) --(5,0) ;
			\node[left, color=SeaGreen4] at (-3,1) {$f_-$};
	\end{tikzpicture}
	\caption{Zerteilung von $f$ in $f_+$ und $f_-$}
\end{wrapfigure}
Falls $f$ integrierbar ist, so auch
\begin{align*}
	f_+ &:= \max \set{f,0} = \frac{1}{2} \cdot  (f+ \abs{f})  \\
	f_- &:= \max \set{-f,0} = \frac{1}{2} \cdot (-f + \abs{f} )  
\end{align*}

Dann gilt $f= f_+ - f_-$ und $f_+ \cdot f_- =0$. Entsprechend für 
\begin{align*}
	\max \set{f,g} &= \frac{1}{2} (f+g + \abs{f-g}), \\
	\min \set{f,g} &= \frac{1}{2} (f+g - \abs{f-g} )  
\end{align*}
% subsection 44 (end)

\subsection[Definition: Integral einer nur auf $A \subset \mathds{R}^n$ definierten Funktion]{Definition} % (fold)
\label{sub:45}
Sei $A \subset \mathds{R}^n$, $f : A \to (- \infty, \infty]$. $f$ heißt integrierbar über $A$, falls $f_A : \mathds{R}^n \to (-\infty, \infty]$ gegeben durch\marginnote{triviale Fortsetzung} 
\[
	f_A (x) = \begin{cases}
		f(x), &\text{ falls }x \in A\\
		0 , &\text{ sonst}
	\end{cases}
\]
integrierbar ist.\index{triviale Fortsetzung} Wir schreiben 
\begin{align*}
	\int_A\! f  \, \mathrm{d}x  &:= \int_{\mathds{R}^n}\! f_A  \, \mathrm{d}x  \quad , \quad \norm[1,A]{f} := \norm[1]{f_A}  \\
	\mathcal{L}^{1}(A) &:= \set[g : A \to \mathds{R}]{g \text{ integrierbar über }A} 
\end{align*}
Proposition \ref{sub:43} und Bemerkung \ref{sub:44} gelten entsprechend.
% subsection 45 (end)

\subsection[Proposition: Regelfunktionen sind Lebesgue-integrierbar]{Proposition} % (fold)
\label{sub:46}
Sei $f : [a,b] \to \mathds{R}$ eine Regelfunktion. Dann ist $f$ Lebesgue integrierbar und die Integrale stimmen überein.
\minisec{Beweis}
Sei $(\varphi_k)_{k \in \mathds{N}} \subset \mathcal{T}([a,b])$ eine Folge von Treppenfunktionen mit $\norm[{\infty,[a,b]}]{f- \varphi_k} \xrightarrow{k \to \infty} 0$.
Dann 
\begin{align*}
	\norm[1]{f_{[a,b]} -( \varphi_k)_{[a,b]}} = \norm[1]{\abs{f- \varphi_k}_{[a,b]} } &= \norm[1]{\abs{(f- \varphi_k)_{[a,b]}} \cdot \chi_{[a,b]}}  \\
	&\le \norm[1]{\norm[{\infty, [a,b]}]{f- \varphi_k} \cdot  \chi_{[a,b]} } \\ 
	&= \norm[{\infty,[a,b]}]{f- \varphi_k} \cdot (b-a)\\
	&\xrightarrow{k \to \infty} 0
\end{align*}
$\Rightarrow f$ ist Lebesgue integrierbar und 
\[
	\int_{[a,b]}\! f  \, \mathrm{d}x \stackrel{\text{\ref{sub:45}}}{=} \int_{\mathds{R}}\! f_{[a,b]}  \, \mathrm{d}x \stackrel{\text{\ref{sub:41}}}{=}
	\lim_{ k \to \infty} \int_\mathds{R}\! (\varphi_k)_{[a,b]}  \, \mathrm{d}x = \lim_{ k \to \infty} \int_{a} ^{b} \! \varphi_k  \, \mathrm{d}x = \int_{a} ^{b} \! f  \, \mathrm{d}x \marginnote{die letzten beiden Integrale sind Regel-Integrale}\bewende
\]
% subsection 46 (end)

\subsection{Satz (Translationsinvarianz)} % (fold)
\label{sub:47}
Sei $f : \mathds{R}^n \to (-\infty, \infty]$ integrierbar und $a \in \mathds{R}^n$. Dann ist auch $f_a : \mathds{R}^n \to (-\infty, \infty]$, gegeben durch 
$f_a(x) := f(x-a)$ integrierbar und es gilt
\[
	\int_{\mathds{R}^n}\! f  \, \mathrm{d}x = \int_{\mathds{R}^n}\! f_a  \, \mathrm{d}x 
\]
\minisec{Beweis}
Für jeden Quader $Q \subset \mathds{R}^n$ ist $a+ Q := \set[x \in \mathds{R}^n]{x-a \in Q} $ wieder ein Quader. \hfill (warum?)\\
Es gilt 
\[
	\int\! \chi_Q  \, \mathrm{d}x  = v(Q) = v(a+ Q) = \int\! \chi_{a+Q}  \, \mathrm{d}x = \int\! (\chi_Q)_{a}  \, \mathrm{d}x 
\]
Der Satz gilt daher für Funktionen der Form $\chi_Q$; wegen der Linearität auch für $f \in \mathcal{T}_{\mathds{R}^n}$, denn 
$(\alpha \cdot  f + \cdot \beta g)_a = \alpha \cdot f_a + \beta \cdot g_a$. 

Ist nun $g : \mathds{R}^n \to (-\infty,\infty]$ eine Abbildung und 
$\Phi = \sum_{k=0}^{\infty} c_k \cdot \chi_{Q_k}$ eine Treppe über $g$, so ist $\Phi_a := \sum_{k=0}^{\infty} c_k \cdot (\chi_{Q_k})_a$ eine Treppe über $g_a$ und 
$I(\Phi)= I(\Phi_a)$ $\Rightarrow \norm[1]{g} = \norm[1]{g_a}  $. \\
Sei $(\varphi_k)_\mathds{N} \subset \mathcal{T}_{\mathds{R}^n}$ mit $\norm[1]{f- \varphi_k} \to 0$. Dann gilt auch 
\begin{align*}
	\norm[1]{f_a - (\varphi_k)_a} = \norm[1]{(f- \varphi_k)_a}  = \norm[1]{f- \varphi_k} \xrightarrow{k \to \infty} 0. 
\end{align*} 
Also ist $f_a$ integrierbar; es gilt weiter 
\begin{align*}
	\int\! f_a  \, \mathrm{d}x  = \lim_{ k \to \infty} \int\! (\varphi_k)_a  \, \mathrm{d}x = \lim_{ k \to \infty} \int\! \varphi_k = \int\! f  \, \mathrm{d}x. \bewende
\end{align*}
% subsection 47 (end)

\subsection{Satz (Beppo Levi mit Treppenfunktionen)} % (fold)
\label{sub:48}
Sei $f : \mathds{R}^n \to (- \infty, \infty]$ eine Abbildung und $(\varphi_k)_{k \in \mathds{N}} \subset \mathcal{T}_{\mathds{R}^n}$ eine Folge von Treppenfunktionen mit
\begin{enumerate}[(i)]
	\item $(\varphi_k)_{k \in \mathds{N}}$ ist monoton wachsend,
	\item $\varphi_k \xrightarrow{k \to \infty} f$ punktweise und
	\item $\enbrace{\int\! \varphi_k}_{k \in \mathds{N}}$ ist beschränkt (also konvergent, da monoton).
\end{enumerate}
Dann ist $f$ integrierbar und es gilt
\[
	\int\! f  \, \mathrm{d}x = \lim_{ k \to \infty} \int\! \varphi_k  \, \mathrm{d}x 
\]
\minisec{Beweis}
Für $x \in \mathds{R}^n$ gilt 
\begin{align*}
	f(x) - \varphi_m(x) &= \lim_{ k \to \infty} \enbrace{\varphi_k (x) - \varphi_m(x)} \\
	&= \lim_{ k \to \infty} \enbrace{\sum_{i=m}^{k-1} \varphi_{i+1}(x)- \varphi_i(x)} \marginnote{Teleskopsumme} \\
	&= \sum_{i=m}^{\infty} \varphi_{i+1}(x) - \varphi_i(x) \\
	\intertext{Es folgt} f- \varphi_m &= \sum_{i=m}^{\infty} \underbrace{\varphi_{i+1}- \varphi_i}_{\ge 0} \\
	\intertext{Damit ergibt sich} \Rightarrow \norm[1]{f- \varphi_m} &\stackrel{\text{\ref{sub:39}}}{\le} \sum_{i=m}^{\infty} \norm[1]{\varphi_{i+1}- \varphi_i} \\
	&\stackrel{\text{\ref{sub:311}}}{=}  \sum_{i=m}^{\infty} \int \! \abs{\varphi_{i+1} - \varphi_i}  \, \mathrm{d}x = \sum_{i=m}^{\infty} \int \! \enbrace{\varphi_{i+1} - 
	\varphi_i}  \, \mathrm{d}x  \\
	&= \sum_{i=m}^{\infty} \enbrace{\int\! \varphi_{i+1}  \, \mathrm{d}x - \int\! \varphi_i  \, \mathrm{d}x } \\
	&= \lim_{ k \to \infty} \int\! \varphi_k  \, \mathrm{d}x - \int\! \varphi_m  \, \mathrm{d}x  \xrightarrow{m \to \infty} 0 \bewende
\end{align*}
% subsection 48 (end)

\subsection[Proposition: Stetige Funktionen lassen sich durch Treppenfunktionen approximieren]{Proposition} % (fold)
\label{sub:49}
Sei $U \subset \mathds{R}^n$ offen und $f : U \to \mathds{R}_+$ und stetig. Dann existiert $(\varphi_k)_{k \in \mathds{N}} \subset \mathcal{T}_{\mathds{R}^n}$ mit
\begin{itemize}
	\item $\varphi_k : \mathds{R}^n \to \mathds{R}_+$,
	\item $(\varphi_k)_\mathds{N}$ monoton wachsend,
	\item $\varphi_k \xrightarrow{k \to \infty}  f_U$ punktweise.
\end{itemize}
\minisec{Beweis}
Sei $Q_{x,t} \subset \mathds{R}^n$ für $(x,t) \in \mathds{Q}^n \times \mathds{Q}$ durch $Q_{x,t} := [x_1 -t, x_1 +t] \times \ldots \times [x_n -t, x_n +t]$. Setze
\[
	Q'_{x,t} := \begin{cases}
		Q_{x,t} &\text{ falls } Q_{x,t} \subset U\\
		\emptyset  &\text{ sonst}
	\end{cases}
\]
und
\[
	c_{x,t} := \begin{cases}
		\min \set[f(y)]{y \in Q'_{x,t}}  &\text{ falls }Q'_{x,t} \not=\emptyset \\
		0 &\text{ sonst}
	\end{cases}
\]
(Das Minimum existiert aufgrund der Kompaktheit von $Q'_{x,t}$)\\
Dann gilt für $(x,t) \in \mathds{Q}^n \times \mathds{Q}$ folgendes: $c_{x,t} \cdot \chi_{Q'_{x,t}} \le f$. Außerdem gilt: Zu jedem $y \in U$, $\varepsilon>0$ existiert ein
$(x,t) \in \mathds{Q}^n \times \mathds{Q}$ mit 
\begin{align*}
	\abs{c_{x,t} \cdot \chi_{Q'_{x,t}}(y) - f(y)} < \varepsilon \tag{warum?} 
\end{align*}
Wähle $\gamma : \mathds{N} \to \mathds{Q}^n \times \mathds{Q}$ surjektiv und definiere $\varphi_k : \mathds{R}^n \to \mathds{R}$ durch 
\[
	\varphi_k (y) := \max \set[c_{\gamma(i)} \cdot \chi_{Q'_{\gamma(i)}}(y)]{i = 0, \ldots , k} 
\]
Dann ist $(\varphi_k)_\mathds{N} \subset \mathcal{T}_{\mathds{R}^n}$ wie gewünscht. \hfill (warum?)
% subsection 49 (end)

\subsection[Satz: Stetige, beschränkte Funktionen sind über offene, beschränkte Mengen integrierbar]{Satz} % (fold)
\label{sub:410}
Sei $U \subset \mathds{R}^n$ offen und beschränkt, $f : U \to \mathds{R}$ stetig und beschränkt. Dann ist $f$ über $U$ integrierbar.
\minisec{Beweis}
Wegen $f = f_+ - f_-$ (mit $f_+,f_- : U \to \mathds{R}_+$ stetig)\marginnote{siehe auch \ref{sub:44}} und Proposition (\ref{sub:43}) genügt es zu zeigen, dass $f_+$ und $f_-$ integrierbar sind. Wir dürfen also o.E.d.A. annehmen, dass $0\le f\le M \in \mathds{R}_+$. Wähle $(\varphi_k)_\mathds{N} \subset \mathcal{T}_{\mathds{R}^n}$ wie in
(\ref{sub:49}). (D.h. $\varphi_k \ge 0$, $\varphi_k \to f_U$ monoton und punktweise). Wähle $Q \subset \mathds{R}^n$ ein Quader mit $U \subseteq Q$, dann ist
\[
	\int\! \varphi_k  \, \mathrm{d}x  \le \int\! M \cdot \chi_U  \, \mathrm{d}x  \le \int\! M \cdot \chi_Q  \, \mathrm{d}x = M \cdot v(Q) < \infty
\]
Mit Satz \ref{sub:48} folgt: $f_U$ ist integrierbar über $\mathds{R}^n$, also ist $f$ integrierbar über $U$. \bewende
% subsection 410 (end)

\subsection{Satz (Fubini für stetige Funktionen auf offenen Teilmengen)} % (fold)
\label{sub:411}
Sei $U \subset \mathds{R}^p \times \mathds{R}^q$ offen und beschränkt, $f : U \to \mathds{R}$ stetig und beschränkt. Für $y \in \mathds{R}^q$ definiere 
$f_y : \mathds{R}^p \to \mathds{R}$ durch $f_y(x) := f_U(x,y)$. Dann ist $f_y$ integrierbar, die Funktion $ F : \mathds{R}^q \to \mathds{R}$, definiert durch
$F(y) := \int_{\mathds{R}^n}\! f_y(x)  \, \mathrm{d}x $ ist ebenfalls integrierbar und es ist 
\[
	\int_{\mathds{R}^n}\! f_U(x,y)  \, \mathrm{d}(x,y) = \int_{\mathds{R}^q} \! F(y) \, \mathrm{d}y = \int_{\mathds{R}^q}\! \enbrace{\int_{\mathds{R}^p} f_U(x,y)  
	\, \mathrm{d}x }  \, \mathrm{d}y.
\]
\minisec{Beweis}
O.E.d.A $0 \le f \le M$. Sei $(\varphi_k)_\mathds{N} \subset \mathcal{T}_{\mathds{R}^n}$ wie in Proposition \ref{sub:49}; wie im Beweis von \ref{sub:410} ist 
$(\int\! \varphi_k  \, \mathrm{d}x )_\mathds{N}$ beschränkt. Nach Satz \ref{sub:48} gilt 
\begin{align*}
	\int_{\mathds{R}^n}\! f_U (x,y)  \, \mathrm{d}(x,y) = \lim_{ k \to \infty} \int_{\mathds{R}^n}\! \varphi_k(x,y)  \, \mathrm{d}(x,y)
\end{align*}
Für $y \in \mathds{R}^q$ ist $\psi_{y,k} : \mathds{R}^p \to \mathds{R}$, definiert durch $\psi_{y,k}(x) := \varphi_k(x,y)$ in $\mathcal{T}_{\mathds{R}^p}$ (wie in
\ref{sub:33}). Dann gilt $\psi_{y,k} \xrightarrow{k \to \infty} f_y$ punktweise und monoton wachsend. Wie im Beweis von \ref{sub:410} sieht man, dass 
$\enbrace{\int_{R^p}\! \psi_{y,k}  \, \mathrm{d}x }_\mathds{N} $ beschränkt ist. Mit Satz \ref{sub:48} folgt 
\begin{align*}
	F(y) := \int_{\mathds{R}^p}\! f_U (x,y)  \, \mathrm{d}x  =  \int_{\mathds{R}^p}\! f_y(x)  \, \mathrm{d}x = \lim_{ k \to \infty} \int_{\mathds{R}^p}\! \psi_{y,k}(x)
	 \, \mathrm{d}x 
\end{align*}
Definiere $\Phi_k : \mathds{R}^q \to \mathds{R}$ durch $\Phi_k (y) := \int_{\mathds{R}^p} \! \varphi_k (x,y)  \, \mathrm{d}x  $ wie im Beweis \ref{sub:33} (bzw. \ref{sub:34}), sieht man $\Phi_k \in \mathcal{T}_{\mathds{R}^q}$. Es gilt
\begin{align*}
	\int_{\mathds{R}^q}\! \Phi_k (y)  \, \mathrm{d}y = \int_{\mathds{R}^n}\! \varphi_k (x,y)  \, \mathrm{d}(x,y) \le \int_{\mathds{R}^n}\! f_U(x,y)  \, \mathrm{d}(x,y)
	< \infty
\end{align*}
Die Folge $(\Phi_k)_\mathds{N}$ ist monoton (wegen der Monotonie der $\varphi_k$ und des Integrals). Es gilt $\Phi_k (y) \xrightarrow{k \to \infty} F(y)$, d.h.
$\Phi_k \nearrow F$ punktweise und monoton. Mit \ref{sub:48} folgt, dass $F$ integrierbar ist mit 
\[
	\int_{\mathds{R}^q}\! F(y)  \, \mathrm{d}y = \lim_{ k \to \infty} \int_{\mathds{R}^q}\! \Phi_k(y)  \, \mathrm{d}y = \lim_{ k \to \infty} \int_{\mathds{R}^n}\!
	\varphi_k(x,y)  \, \mathrm{d}(x,y) = \int_{\mathds{R}^n}\! f_U(x,y)  \, \mathrm{d}(x,y) \bewende
 \]
% subsection 411 (end)
% section 4 (end)
\newpage

\section{Messbarkeit in $\mathds{R}^n$, Nullmengen} % (fold)
\label{sec:5}

\subsection[Definition: Lebesgue-messbar]{Definition} % (fold)
\label{sub:51}
Eine Menge $A \subset \mathds{R}^n$ heißt \Index{Lebesgue-messbar}, falls $\chi_A$ integrierbar ist. Wir setzen $v(A) := \int_{\mathds{R}^n}\! \chi_A$.
% subsection 51 (end)

\subsection[Proposition: Einfache Mengenoperationen auf messbaren Mengen]{Proposition} % (fold)
\label{sub:52}
Seien $A,B \subset \mathds{R}^n$ messbar, dann sind $A \cup B$, $A \cap B$ messbar mit $v(A \cup B) = v(A) + v(B)- v(A \cap B)$. Falls $A \subset B$ gilt, so ist
$v(A) \le v(B)$.
\minisec{Beweis}
Benutze (\ref{sub:43}) und:
\begin{enumerate}[(i)]
	\item $\chi_{A \cup B} = \chi_A + \chi_B - \chi_{A \cap B}$
	\item $A \subset B \Longrightarrow \chi_A \le \chi_B$ \bewende
\end{enumerate}
% subsection 52 (end)

\subsection[Proposition: Beschränkte Teilmengen von $\mathds{R}^n$ sind messbar]{Proposition} % (fold)
\label{sub:53}
Sei $A \subset \mathds{R}^n$ beschränkt und offen oder kompakt. Dann ist $A$ messbar.
\minisec{Beweis}
\begin{enumerate}[1)]
	\item Sei $A \subset \mathds{R}^n$ beschränkt und offen. Dann ist $f \equiv 1$ stetig und beschränkt ($f : A \to \mathds{R}$).  Mit \ref{sub:410} folgt, dass
	$f_A$ integrierbar ist. $\Rightarrow f_A = \chi_A$ integrierbar $\Rightarrow $ A messbar.
	\item Sei $A$ kompakt. Dann existiert ein offener Quader $Q \subset \mathds{R}^n$ mit $A \subseteq Q$, dann ist $Q\setminus A \off \mathds{R}^n$ und beschränkt, also 
	sind $\chi_{Q\setminus A}, \chi_Q$ integrierbar, also $\chi_A = \chi_Q - \chi_{Q\setminus A}$ integrierbar. \bewende
\end{enumerate}
% subsection 53 (end)

\subsection[Beispiel: Berechnung des Volumens eines Kreises und einer Kugel]{Beispiel} % (fold)
\label{sub:54}
\begin{enumerate}[(i)]
	\item $U_r := B_{\mathds{R}^2}(0,r) \off \mathds{R}^2$
	\begin{align*}
		v(U_r) = \int_{\mathds{R}^2}\! \chi_{U_r}(x,y)  \, \mathrm{d}(x,y) &= \int_{\mathds{R}}\! \enbrace{\int_\mathds{R}\! \chi_{U_r}(x,y)  \, \mathrm{d}x }  \, 
		\mathrm{d}y \\
		&= \int_\mathds{R} \! \enbrace{\int_\mathds{R}\! \chi_{(-r,r)}(y) \cdot \chi_{\enbrace{-\sqrt{r^2-y^2}, \sqrt{r^2-y^2}}}(x)  \, \mathrm{d}x  }  \, \mathrm{d}y \\
		&= \int_\mathds{R}\! \chi_{(-r,r)}(y) \cdot \enbrace{\int_\mathds{R}\! \chi_{\enbrace{-\sqrt{r^2-y^2}, \sqrt{r^2-y^2}}}(x)  \, \mathrm{d}x }  \, \mathrm{d}y \\
		&= \int_\mathds{R}\! \chi_{(-r,r)} (y) \cdot 2 \cdot \sqrt{r^2-y^2}  \, \mathrm{d}y \\
		&= \int_{-r} ^{r} \! 2 \cdot \sqrt{r^2-y^2}  \, \mathrm{d}y   \\
		&= 2 \cdot \int_{-r} ^{r} \! r \cdot \sqrt{1-\left(\frac{y}{r} \right)^2}   \, \mathrm{d}y \tag{Substitution $s= \frac{y}{r}$, $\diffd{y}{s} = r$} \\
		&= 2 \cdot \int_{s(-r)} ^{s(r)} \! r \cdot r \cdot \sqrt{1-s^2}  \, \mathrm{d}s \\
		&= 2 \cdot r^2 \cdot \int_{-1} ^{1} \! \sqrt{1-s^2}  \, \mathrm{d}s \tag{Analysis \RM{1}} \\
		&= r^2 \cdot \pi   
	\end{align*}
	\item $S^1 \subset \mathds{R}^2$ kompakt also messbar nach \ref{sub:53}. $S_1 \subset U_{1+\varepsilon} \setminus U_1$. Mit \ref{sub:52} folgt 
	\marginnote{$S^1$ ist die Sphäre, also die Kugeloberfläche}
	\begin{align*}
		v(U_{1+\varepsilon}) &= v\big(U_{1}\cup (U_{1+\varepsilon} \setminus U_1)\big) = v(U_1) + v(U_{1+\varepsilon} \setminus U_1) - v(\emptyset) \\
		\intertext{damit folgt} v(S^1) &\le v(U_{1+\varepsilon}) - v(U_1) = \pi \cdot \big((1+\varepsilon)^2 - 1^2\big) \xrightarrow{\varepsilon \to 0} 0
	\end{align*}
	\item $K_r := B_{\mathds{R}^3}(0,r) \subset \mathds{R}^3$ offen.
	\[
		U_{r,z} := B_{\mathds{R}^2}(0,\sqrt{r^2- z^2}  )
		\qquad \quad 
		\vcenter{\hbox{\begin{tikzpicture}[scale=0.5]
			%\coordinate (A) at (-2,0);
			%\coordinate (B) at (2,0);
			\coordinate (a) at (-2.2,0);
			\coordinate (b) at (0,-2.2);
			\coordinate (c) at (2.2,0);
			\coordinate (d) at (0,2.2);
			\coordinate (M) at (0,0);
			\foreach \x in{a,b,c,d}%
				\draw[color=Grey0, dashed] (M) -- (\x);
			\draw[thick] (0,0) circle[radius=2cm];
			\draw[thick, fill=SeaGreen1, fill opacity=0.4] (0,0) ellipse (2cm and 0.45cm);
			\draw[thick,fill=Tomato2, fill opacity=0.5] (0,1) ellipse (1.72cm and 0.35cm);
			\node[color=Tomato2, right] at (1.72,1) {$U_{r,z}$};
		\end{tikzpicture}}}
	\]
	\begin{align*}
		v(K_r) = \int_{\mathds{R}^3}\! \chi_{K_r}(x,y,z)  \, \mathrm{d}(x,y,z) &= \int_\mathds{R}\! \enbrace{\int_{\mathds{R}^2}\! \chi_{K_r}(x,y,z)  \, \mathrm{d}(x,y)}  \, \mathrm{d}z \\
		&= \int_{\mathds{R}}\! \enbrace{\int_{\mathds{R}^2} \chi_{(-r,r)} (z) \cdot \chi_{U_{r,z}}(x,y)  \, \mathrm{d}(x,y)}  \, \mathrm{d}z \\
		&= \int_{\mathds{R}}\! \chi_{(-r,r)}(z) \enbrace{\int_{\mathds{R}^2}\! \chi_{U_{r,z}}(x,y)  \, \mathrm{d}(x,y) }  \, \mathrm{d}z \\
		&= \int_{\mathds{R}}\! \chi_{(-r,r)}(z) \cdot \big(\pi  \cdot (r^2-z^2)\big)  \, \mathrm{d}z \tag*{nach i)} \\
		&= \pi r^2 \int_{-r} ^{r} \! 1  \, \mathrm{d}z - \pi \int_{-r} ^{r} \! z^2  \, \mathrm{d}z \\
		&= \pi r^2 \cdot 2r - \pi \cdot 2 \cdot \frac{r^3}{3} \\
		&= \frac{4}{3}\cdot \pi \cdot  r^3 \bewende 
	\end{align*}
\end{enumerate}
% subsection 54 (end)

\subsection[Bemerkung: Prinzip von Cavalieri]{Bemerkung} % (fold)
\label{sub:55}
Seien $U,W \subset \mathds{R}^{n-1} \times \mathds{R}$ offen. Setze $U_y := \set[x \in \mathds{R}^{n-1}]{(x,y) \in U} $ für $y \in \mathds{R}$,
$W_y := \set[x \in \mathds{R}^{n-1}]{(x,y) \in W} $ für $y \in \mathds{R}$. Dann gilt $v_{\mathds{R}^n}(U) = v_{\mathds{R}^n}(W)$, falls gilt
\[
	v_{\mathds{R}^{n-1}}(U_y) = v_{\mathds{R}^{n-1}}(W_y) \quad \text{ für alle } y \in \mathds{R},
\]
denn $v_{\mathds{R}^n}(U) = \int_{\mathds{R}}\! v_{\mathds{R}^{n-1}} (U_y)  \, \mathrm{d}y = \int_{\mathds{R}}\! v_{\mathds{R}^{n-1}}(W_y)  \, \mathrm{d}y
= v_{\mathds{R}^n}(W)$ gilt nach Fubini (\ref{sub:411}). Mit Translationsinvarianz (\ref{sub:47}) ergibt sich das \Index{Prinzip von Cavalieri}:
\begin{figure}[h]
	\centering{{\begin{tikzpicture}[scale=0.7]
		%\draw[help lines] (0,0) grid (8,4);
		\coordinate (A) at (0,0);
		\coordinate (B) at (8,0);
		\coordinate (C) at (0,4);
		\coordinate (C') at (9,4);
		\draw[dashed, color=Grey0, thick] (C) --(C');
		
		\begin{scope}% wichtig damit die Wirkung von clip begrenzt wird
		% Schneiden und Füllen der eigentlichen Form
		  \clip[postaction={fill=Sienna3,fill opacity=0.2},](A)--(B)--(C);
		% Zeichnen der Streifen
		  \foreach \x in {0.5,1,1.5,2,2.5,...,4}%
		           \draw[Sienna3!80,line width=0.5mm](0,\x)--(9,\x);
		\end{scope}
		
		\begin{scope}% wichtig damit die Wirkung von clip begrenzt wird
		% Schneiden und Füllen der eigentlichen Form
		  \clip[postaction={fill=Chartreuse3,fill opacity=0.2},](A)--(B)--(C');
		% Zeichnen der Streifen
		  \foreach \x in {0.5,1,1.5,2,2.5,...,4}%
		           \draw[Chartreuse3!80,line width=0.5mm, dashed](0,\x)--(9,\x);
		\end{scope}
		\draw (A) -- (B) -- (C) -- cycle;
		\draw (A) -- (B) -- (C') -- cycle;
	\end{tikzpicture}
	\caption[Prinzip von Cavalieri]{Prinzip von Cavalieri: Beide Dreiecke haben den gleichen Flächeninhalt}}}
\end{figure}
% subsection 55 (end)

\subsection[Proposition und Definition: Nullmenge]{Proposition und Definition} % (fold)
\label{sub:56}
$N \subset \mathds{R}^n$ heißt \Index{Lebesgue-Nullmenge}, falls sie eine der folgenden äquivalenten Eigenschaften erfüllt:
\begin{enumerate}[(i)]
	\item $N$ ist messbar und $v(N)=0$
	\item $\norm[1]{\chi_N}=0$
\end{enumerate} 
Sei $E : \mathds{R}^n \to \set{\text{wahr}, \text{falsch}} $, so sagen wir $E$ gilt \Index{fast überall} (f.ü.), falls $N_E := \set[x \in \mathds{R}^n]{E(x)= \text{falsch}} $ eine
Nullmenge ist.
\minisec{Beweis}
\begin{description}
	\item[(i) $\Rightarrow $ (ii):]  $0 \le \norm[1]{\chi_N}  \stackrel{\text{\ref{sub:43} b)}}{=} \int\! \chi_N  \, \mathrm{d}x = v(N) = 0 $.
	\item[(ii) $\Rightarrow $ (i):] Setze $\varphi_k \equiv 0 \in \mathcal{T}_{\mathds{R}^n}$ für $k \in \mathds{N}$, dann 
	\[
		\norm[1]{\chi_N - \varphi_k} = \norm[1]{\chi_N} = 0  
	\]
	also ist $\chi_N$ integrierbar. Nach Proposition \ref{sub:43} b) gilt $0 \le \int\! \chi_N  \, \mathrm{d}x = \norm[1]{\chi_N} =0$ \bewende
\end{description}
% subsection 56 (end)

\subsection[Proposition: Teilmengen von Nullmengen und Vereinigungen von Nullmengen]{Proposition} % (fold)
\label{sub:57}
\begin{enumerate}[(i)]
	\item Jede Teilmenge einer Nullmenge ist eine Nullmenge.
	\item Jede abzählbare Vereinigung von Nullmengen ist eine Nullmenge.
\end{enumerate}
\minisec{Beweis}
\begin{enumerate}[(i)]
	\item Sei $M \subset N$. Es folgt $0 \le \chi_M \le \chi_N$. Mit Bemerkung \ref{sub:38} folgt $0 \le \norm[1]{\chi_M} \le \norm[1]{\chi_N}  $
	\item Sei $N= \bigcup_{k=0}^\infty N_k$. Dann gilt $0 \le \chi_N \le \sum_{k=0}^{\infty} \chi_{N_k}$. Nach Bemerkung \ref{sub:38} und Proposition \ref{sub:39} gilt
	\[
		0 \le \norm[1]{\chi_N}  \le \norm[1]{\sum_{k=0}^{\infty} \chi_{N_k}} \le \sum_{k=0}^{\infty} \norm[1]{\chi_{N_k}}  =0 \bewende
	\]
\end{enumerate}
% subsection 57 (end)

\subsection[\protect{Proposition: $f$ mit $\norm[1]{f}< \infty $ ist fast überall endlich}]{Proposition} % (fold)
\label{sub:58}
Sei $f : \mathds{R}^n \to (-\infty, \infty]$ eine Abbildung mit $\norm[1]{f}< \infty $. Dann ist $f$ fast überall endlich, d.h. $N:= \set[x \in \mathds{R}^n]{f(x)=\infty} $
ist eine Nullmenge.
\minisec{Beweis}
Für $\varepsilon >0 $ ist $0\le \chi_N\le \varepsilon \cdot \abs{f} $, also
\begin{align*}
	0 \le \norm[1]{\chi_N} \stackrel{\text{\ref{sub:38}}}{\le} \norm[1]{\varepsilon \cdot \abs{f} } \stackrel{\text{\ref{sub:38}}}{=} \varepsilon \cdot \norm[1]{f}
	\xrightarrow{\varepsilon \to 0} 0   
\end{align*}
Daraus folgt $\norm[1]{\chi_N} =0 $ \bewende
% subsection 58 (end)

\subsection[Proposition: Die Integrale fast gleicher Funktionen stimmen überein]{Proposition} % (fold)
\label{sub:59}
Seien $f,g : \mathds{R}^n \to (-\infty, \infty]$ Abbildungen mit $f=g$ fast überall. Falls $f$ integrierbar ist, dann auch $g$ und es gilt 
\[
	\int\! f  \, \mathrm{d}x  = \int\! g  \, \mathrm{d}x 
\]
\minisec{Beweis}
Sei $N := \set[x \in \mathds{R}^n]{f(x)\not= g(x)} $, dann gilt $\int\! \chi_N  \, \mathrm{d}x =0$ nach Vorraussetzung. Definiere $h : \mathds{R}^n \to (-\infty, \infty]$
durch 
\[
	h(x) := \infty \cdot \chi_N(x) = \begin{cases}
		\infty, &\text{ falls }x \in N\\
		0 , &\text{ falls } x \not\in N
	\end{cases}
\]
Dann gilt $h = \sum_{k=0}^{\infty} \chi_N$ und $\norm[1]{h} \stackrel{\text{\ref{sub:39}}}{\le} \sum_{k=0}^{\infty} \norm[1]{\chi_N} = 0  $. $f$ ist integrierbar, als gilt
$\exists (\varphi_k)_\mathds{N} \subset \mathcal{T}_{\mathds{R}^n}$ mit $\norm[1]{f- \varphi_k} \xrightarrow{k \to \infty} 0 $. Es gilt 
$\abs{g - \varphi_k} \le   \abs{f- \varphi_k} +h$ also 
\[
	\norm[1]{g- \varphi_k} \le \norm[1]{\abs{f- \varphi_k} +h } \le \norm[1]{f- \varphi_k} + \norm[1]{h} \xrightarrow{k \to \infty} 0    
\]
mit (\ref{sub:41}) folgt: $g$ integrierbar und $\int\! g  \, \mathrm{d}x  = \lim_{ k \to \infty} \int\! \varphi_k  \, \mathrm{d}x = \int\! f  \, \mathrm{d}x $ \bewende
% subsection 59 (end)

\subsection[Beispiele für Nullmengen]{Beispiele} % (fold)
\label{sub:510}
\begin{enumerate}[(i)]
	\item $\mathds{Q}$ ist Nullmenge in $\mathds{R}$: $v(\{x\}) =0$, $x \in \mathds{R}$. $\mathds{Q}$ ist abzählbare Vereinigung solcher Einpunktmengen, also ist 
	$\mathds{Q}$ nach \ref{sub:57} (ii) eine Nullmenge.
	
	(Alternativ: Benutze Satz \ref{sub:48} (Beppo-Levi))
	\item Sei $C \subset [0,1]$ die Cantormenge, dann ist $C$ eine Nullmenge: 
	\[
		[0,1] \setminus C =  \dot \bigcup_{l=0}^\infty U_l
	\]
	Setze $\varphi_k := \chi_{\bigcup_{l=0}^k U_l}$, $k \in \mathds{N}$. Dann ist $(\varphi_k)_\mathds{N} \subset \mathcal{T}_{\mathds{R}}$, $\varphi_k \xrightarrow{k \to \infty} \chi_{[0,1] \setminus C}$ monoton wachsend und punktweise und $\int\! \varphi_k  \, \mathrm{d}x  \le 1$, $k \in \mathds{N}$.
	
	Mit \ref{sub:48} folgt 
	\begin{align*}
		\int\! \chi_{[0,1] \setminus C}  \, \mathrm{d}x = \lim_{ k \to \infty} \int\! \varphi_k  \, \mathrm{d}x  &= \lim_{ k \to \infty} \sum_{l=0}^{k} v(U_l)  \, \mathrm{d}x \\
		&= \lim_{ k \to \infty} \sum_{l=0}^{k} \frac{2^l}{3^{l+1}} = \frac{1}{3} \sum_{l=0}^{\infty} \enbrace{\frac{2}{3} }^l = \frac{1}{3} \cdot \frac{1}{1-\frac{2}{3} }
		= 1       
	\end{align*}
	Es folgt 
	\begin{align*}
		\int\! \chi_C  \, \mathrm{d}x = \int\! \enbrace{\chi_{[0,1]} - \chi_{[0,1] \setminus C}}  \, \mathrm{d}x = \int\! \chi_{[0,1]}  \, \mathrm{d}x  - \int\!
		\chi_{[0,1] \setminus C}  \, \mathrm{d}x  = 1-1 = 0 
	\end{align*}
	Die Cantormenge ist also eine überabzählbare Nullmenge! \bewende
\end{enumerate}
% subsection 510 (end)

\subsection[Bemerkung: Alternative Charakterisierung von Nullmengen]{Bemerkung} % (fold)
\label{sub:511}
Nullmengen lassen sich auch wie folgt charakterisieren: \\
$N \subset \mathds{R}^n$ ist Nullmenge genau dann, wenn Folgendes gilt:
\[
	\forall \varepsilon > 0 \enspace \exists Q_0, Q_1, \ldots  \subset \mathds{R}^n  \text{ offene Quader mit } N \subset \bigcup_{i \in \mathds{N}}Q_i \text{ und }
	\sum_{i \in \mathds{N}} v(Q_i) < \varepsilon
\]
(ohne Beweis)
% subsection 511 (end)

\subsection[Proposition: Beschränkte Folge messbarer Mengen]{Proposition} % (fold)
\label{sub:512}
Sei $Q \subset \mathds{R}^n$ ein Quader und $A_0, A_1, \ldots  \subset Q$ eine Folge messbarer Mengen. Dann ist $\bigcup_{l \in \mathds{N}} A_l$ messbar.
\minisec{Beweis}
Wende Beppo-Levi\marginnote{spätere Variante von Beppo-Levi} an auf 
\[
	\varphi_k := \chi_{\bigcup_{l=0}^k A_l} \bewende
\]
% subsection 512 (end)

\subsection[Bemerkung: Zusammenfassung über Messbarkeit]{Bemerkung} % (fold)
\label{sub:513}
Wir haben jetzt: 
\begin{itemize}
	\item Beschränkte offene   Mengen sind messbar.
	\item Kompakte Mengen sind messbar.
	\item beschränkte abzählbare Vereinigungen von messbaren Mengen sind messbar.
	\item Komplemente messbarer Mengen in beschränkten messbaren Mengen sind messbar.
\end{itemize}
\uline{Frage:} Existieren überhaupt (beschränkte) nicht messbare Mengen?
% subsection 513 (end)

\subsection[Satz: Existenz einer beschränkten nicht messbaren Teilmenge in $\mathds{R}$]{Satz} % (fold)
\label{sub:514}
$\mathds{R}$ besitzt\footnote{in ZFC} eine beschränkte Teilmenge, die nicht messbar ist.
\minisec{Beweis}
Für $x \in [0,1]$ definiere $E_x := \set[{y \in [0,1]}]{x-y \in \mathds{Q}} \subset [0,1]$. Dann gilt für $x,x' \in [0,1]: E_x= E_{x'}$ oder $E_x \cap E_{x'} = \emptyset$.
($\exists y \in E_x \cap E_{x'} \iff \exists y \in [0,1] x-y \in \mathds{Q}$ und $x'-y \in \mathds{Q} \iff x-x' \in \mathds{Q} \iff E_x = E_{x'}$). Es gilt
$[0,1] = \bigcup_{x \in [0,1]} E_x$.

Setze $A := \set[E]{E=E_x \text{ für ein } x \in [0,1]} \subset \mathcal{P}([0,1])$. Sei $f : A  \to [0,1] = \bigcup_{E \in A} E$ eine Funktion mit $f(E) \in E$, $E \in A$.
$f$ existiert nach dem Auswahlaxiom. Setze $M := f(A) \subset [0,1]$. Sei $J := \mathds{Q} \cap [-1,1]$, dann ist $J$ abzählbar, d.h. es existiert 
$(j_i)_{i \in \mathds{N}}$ mit $J = \set{j_0, j_1, \ldots } $ und $j_i \not= j_{i'}$ für $i \not= i'$. Setze $M_i := M + j_i = \set[y + j_i]{y \in M} \subset[-1,2] $.\\
Behauptungen:
\begin{enumerate}[(i)]
	\item $[0,1] \subset \bigcup_{i \in \mathds{N}} M_i = \set[y+ j_i]{y \in M, i \in \mathds{N}} $
	\item $M_i \cap M_{i'} = \emptyset$ falls $i \not= i'$
\end{enumerate}
\uline{Beweis:}
\begin{enumerate}[(i)]
	\item $x \in [0,1]$, dann ist $f(E_x)  \in E_x$, d.h. $x - f(E_x) \in \mathds{Q} \cap [-1,1] = J$. Also ist $x- f(E_x) = j_i$ für ein $i \in \mathds{N}$. Daraus folgt
	$x= f(E_x) + j_i \in M_i$.
	\item Falls $z \in M_i \cap M_{i'}$, $i \not= i'$, so gilt $z=y+ j_i$, $z= y' + j_{i'}$ für $y,y' \in M$, $j_i \not= j_{i'}$. Es gilt 
	$y= f(E_x) \in E_x$, $y' = f(E_{x'}) \in E_{x'}$, $x,x' \in [0,1]$. Es folgt $y-y' \in \mathds{Q} \cap [-1,1]$ und $E_x \cap E_{x'} \not= \emptyset$, also 
	$E_x = E_{x'}$. Andererseits $j_i \not= j_{i'}$, also $y \not= y'$ %\lightning
\end{enumerate}
Wir haben nun $[0,1] \subset \dot \bigcup_{i \in \mathds{N}} M_i \subset [-1,2]$. Angenommen $M$ sei messbar. Wegen Translationsinvarianz ist dann auch $M_i$ messbar, 
$i \in \mathds{N}$ und $v(M) = v(M_i)$. Weiter gilt 
\[
	v(M_i) = \int\! \chi_{M_i}  \, \mathrm{d}x  = \norm[1]{\chi_{M_i}} 
\]
Wegen $[0,1] \subset \bigcup_{i \in \mathds{N}} M_i$ gilt außerdem 
\[
	0 \le \chi_{[0,1]} \le \chi_{\bigcup_{i \in \mathds{N}} M_i} 
\]
also 
\[
	1 = v([0,1]) = \int\! \chi_{[0,1]}  \, \mathrm{d}x \stackrel{\text{\ref{sub:43}}}{=} \norm[1]{\chi_{[0,1]}} \le \norm[1]{\sum_{i \in \mathds{N}} \chi_{M_i}} 
	\stackrel{\text{\ref{sub:39}}}{\le} \sum_{i \in \mathds{N}} \norm[1]{\chi_{M_i}} = \sum_{i \in \mathds{N}} v(M)
\]
$\Rightarrow v(M) >0$. Wähle $K \in \mathds{N}$ so groß, dass $(K+1) \cdot v(M) >3$ ist, dann 
\[
	3 < (K+1) \cdot v(M) = \sum_{i=0}^{K} v(M_i) = v\enbrace{\dot \bigcup_{i=0}^K } \le v ([-1,2]) = 3  \quad \lightning \bewende 
\]
% subsection 514 (end)

\subsection[Satz: Volumen eines Parallelotops]{Satz} % (fold)
\label{sub:515}
Seien $\underline{x}_1, \ldots , \underline{x}_n \in \mathds{R}^n$. Setze\index{Parallelotop} \marginnote{Parallelotope hatten wir auch schon in LA2}
\[
	P(\underline{x}_1, \ldots , \underline{x}_n) := \set[\underline{x} \in \mathds{R}^n]{\exists t_1, \ldots , t_n \in [0,1] : \underline{x} = \sum_{i=1}^{n} t_i \cdot \underline{x}_i } 
\]
Dann gilt 
\[
	v\big(P(\underline{x}_1, \ldots , \underline{x}_n)\big) = \abs{\det \enbrace{\underline{x}_1, \ldots , \underline{x}_n} } = \abs{ \det \begin{pmatrix}
		x_{1,1} & \cdots & x_{1,n} \\
		\vdots  & & \vdots  \\
		x_{n,1} & \cdots & x_{n,n}
	\end{pmatrix}} 
\]
\minisec{Beweis}
Sei 
\[
	D : \underbrace{\mathds{R}^n \times \ldots \times \mathds{R}^n}_{n \text{-mal}} \to \mathds{R} 
\]
eine Abbildung mit 
\begin{description}
	\item[$(D1)$] $D(\underline{a}_1, \ldots , \lambda \underline{a}_i, \ldots , \underline{a}_n) = \abs{\lambda } \cdot D(\underline{a}_1, \ldots , \underline{a}_i, \ldots , \underline{a}_n) $ für $i \in \{1, \ldots , n\}$, $\lambda  \in \mathds{R}$, $\underline{a}_1, \ldots, \underline{a}_n \in \mathds{R}^n $
	\item[$(D2)$] $D(\underline{a}_1, \ldots , \underline{a}_i+ \underline{a}_j, \ldots , \underline{a}_j, \ldots , \underline{a}_n) = D(\underline{a}_1, \ldots , \underline{a}_n)$ für $i<j$, $\underline{a}_1, \ldots, \underline{a}_n \in \mathds{R}^n $
	\item[$(D3)$] $D(\underline{e}_1, \ldots , \underline{e}_n) = 1$, wo $\underline{e}_i $ der $i$-te Einheitsvektor ist.
\end{description}
Für $\underline{a}_1, \ldots , \underline{a}_n \in \mathds{R}^n$ ist $D(\underline{a}_1, \ldots , \underline{a}_n)$ durch $(D1), (D2), (D3)$ mittels elementaren 
Spaltenumformungen bestimmt (Vergleiche Charakterisierung der Determinante in der Linearen Algebra). Daraus folgt: Es kann höchstens eine Abbildung mit $(D1), (D2), (D3)$
geben. Andererseits erfüllt 
\begin{align*}
	(\underline{a}_1, \ldots , \underline{a}_n) &\mapsto \abs{\det (\underline{a}_1, \ldots , \underline{a}_n)} \\
	\intertext{$(D1), (D2)$ und $(D3)$. Es genügt zu zeigen: Die Abbildung }
	(\underline{a}_1, \ldots , \underline{a}_n) &\mapsto v\big(P(\underline{a}_1, \ldots , \underline{a}_n)\big)
\end{align*}
erfüllt $(D1), (D2), (D3)$. 
\begin{itemize}
	\item $(D3)$: $v\big(P(\underline{e}_1, \ldots , \underline{e}_n)\big) = 1^n = 1$
	\item Falls $\underline{a}_1, \ldots , \underline{a}_n$ linear abhängig sind, dann gilt 
	\[
		v\big(P(\underline{a}_1, \ldots , \underline{a}_n)\big) \stackrel{\text{Übung}}{=}0 = \abs{\det (\underline{a}_1, \ldots , \underline{a}_n)} 
	\]
	und $(D1), (D2)$ sind erfüllt.
\end{itemize}
Seien also $\underline{a}_1, \ldots , \underline{a}_n$ linear unabhängig.
\begin{description}
	\item[$(D1):$] Für $i \in \{1, \ldots , n\}$, $\lambda  \in \mathds{R}$ setze $P_\lambda := P(\underline{a}_1, \ldots , \lambda \cdot \underline{a}_i, \ldots , \underline{a}_n)$; wir wollen 
	\[
		v(P_\lambda ) = \abs{\lambda } \cdot v(P_1) \tag{$\star$} 
	\]
	zeigen: 
	\begin{enumerate}[a)]
		\item Zunächst für $\lambda \in \mathds{N}^*:$ \\
		$\lambda =1$ \checkmark \\
		$\lambda  \mapsto \lambda +1$: Es gilt $P_{\lambda +1} = P_\lambda \cup (\lambda \cdot \underline{a}_i + P_1)$
		\[
			\begin{tikzpicture}[scale=0.5]
				%\draw[help lines] (0,0) grid (16,6);
				\draw[<->, thick, color=Grey0] (0,6) -- (0,0) -- (14,0);
				\coordinate (A1) at (1,3);
				\coordinate (Ai) at (4,1);
				\coordinate (2Ai) at (8,2);
				\coordinate (3Ai) at (12,3);
				\coordinate (A1+i) at (5,4);
				\coordinate (A1+2i) at (9,5);
				\draw[fill=DeepSkyBlue1, fill opacity =0.25] (0,0) -- (A1) -- ++ (Ai) -- (Ai) -- cycle;
				\node[left, color=DeepSkyBlue3] at (3,2) {$P_1$};
				\draw[->, ultra thick] (0,0) -- (Ai) node[midway, above]{$\underline{a}_i$};
				\draw[decorate, decoration={brace,amplitude=5pt, mirror, raise=3pt},yshift=0pt, thick] (9,5) -- (A1) %
				node[above left, midway, yshift=0.27cm, xshift=0.13cm]{$P_\lambda $};
				\draw (A1+i) -- ++ (Ai) -- (2Ai) -- (Ai);
				\draw[decorate, decoration={brace,amplitude=5pt, mirror, raise=2pt},yshift=0pt, thick] (0,0) -- (3Ai) %
				node[below right, midway, yshift=-0.25cm, xshift=-0.1cm]{$P_{\lambda+1} $};
				\draw[fill=DarkOrange1, fill opacity =0.4] (2Ai) -- (3Ai) -- ++ (A1) -- (A1+2i);
				\node[color=DarkOrange3, xshift=-0.2cm] at (11,4) {$\lambda \cdot \underline{a}_i + P_1$};
				\draw[dashed, very thick, ->, DarkOrange2] (0,0) -- (Ai) -- ++ (Ai) node[midway, above]{$\lambda \cdot \underline{a}_i$};
			\end{tikzpicture}
		\]
		Weiter gilt $P_\lambda  \cap (\lambda \cdot \underline{a}_i +P_1) \subset \lambda  \cdot  \underline{a}_i + \mathrm{span} \set{\underline{a}_1, \ldots , \underline{a}_i, 
		\ldots , \underline{a}_n} $ mit $\dim \le n-1$. Also ist 
		\[
			v \big( P_\lambda  \cap (\lambda  \cdot \underline{a}_i + P_1)\big) = 0 \tag{Übung}
		\]
		Also folgt
		\begin{align*}
			v (P_{\lambda +1})&= v(P_\lambda ) + v(\lambda \cdot \underline{a}_i + P_1) - v(P_\lambda \cap \lambda  \cdot \underline{a}_i + P_1) = v(P_\lambda ) +
			v(P_1) - 0 \\ &\stackrel{\text{(I.V.)}}{=} \lambda  \cdot v(P_1)+ v(P_1) = (\lambda +1)\cdot v(P_1)
		\end{align*}
		Induktion $\leadsto$ $(\star)$ für $\lambda  \in \mathds{N}^*$. Ebenso zeigt man 
		$v(P_{q \cdot \lambda }) = q \cdot v(P_\lambda )$, $q \in \mathds{N}^*, \lambda \in \mathds{R}$
		\item Jetzt für $\lambda = \frac{p}{q} $ mit $p,q \in \mathds{N}^*$: \\
		Nach a) gilt $v(P_{q \cdot \lambda })= v(P_p) = p \cdot v(P_1)$ und $v(P_{q \cdot \lambda })= q \cdot v(P_\lambda )$, also
		\[
			p \cdot v(P_1) = q \cdot v(P_\lambda ) \Longrightarrow v(P_\lambda ) = \frac{p}{q} \cdot v(P_1) = \abs{\lambda } \cdot v(P_1). 
		\]
		\item $\lambda \in \mathds{R}_+^*$: Zu $\varepsilon >0$ wähle $r_1, r_2 \in \mathds{Q}_+^*$ mit $r_1 \le \lambda \le r_2$ und 
		$r_2-r_1 < \frac{\varepsilon}{v(P_1)}$. \\(Da $\underline{a}_1, \ldots , \underline{a}_n$ linear unabhängig, enthält 
		$P_1=P(\underline{a}_1, \ldots , \underline{a}_n)$ einen offenen Quader, also ist $v(P_1)>0$).
		
		Dann gilt $P_{r_1} \subset P(\lambda ) \subset P_{r_2}$ und 
		\begin{align*}
			v(P_{r_1}) &\le v(P_{\lambda }) \le v(P_{r_2}) \tag*{und}\\
			r_1 \cdot v(P_1) &\le \lambda  \cdot v(P_1) \le r_2 \cdot v(P_1)
		\end{align*}
		Wir erhalten 
		\begin{align*}
			\abs{v(P_\lambda )- \lambda \cdot v(P_1)} \le r_2 \cdot v(P_1) - r_1 \cdot v(P_1) = (r_2-r_1)\cdot  v(P_1) < \varepsilon
		\end{align*}
		Also $v(P_\lambda )= \lambda  \cdot v(P_1), \lambda  \in \mathds{R}_+^*$
		\item $\lambda =0$: $P_0$ liegt in einer Hyperebene. Nach einer Übung folgt dann
		\[
			v(P_0)=0 = 0 \cdot v(P_1)
		\]
		\item $\lambda < 0$: $P_\lambda = \lambda  \cdot \underline{a}_i + P_{-\lambda }$ und wegen Translationsinvarianz (\ref{sub:47}) gilt dann
		\[
			v(P_\lambda ) = v(P_{-\lambda }) =(-\lambda) \cdot v(P_1) = \abs{\lambda } v(P_1) 
		\]
		Es folgt $(\star)$ für $\lambda \in \mathds{R}$ $\Rightarrow D1$
	\end{enumerate}
	\item[$(D2):$] (nur Idee)\\
	\[
		\begin{tikzpicture}[scale=0.8]
			%\draw[help lines] (0,0) grid (7,3);
			\draw[DarkOrange2!20, fill=DarkOrange2!20, fill opacity=0.4] (0,0)--(4,0)--(6,3)--(2,3)-- cycle;
			\node[DarkOrange2] at (3,2.6) {$P(\underline{a}_1, \underline{a}_2)$};
			\begin{scope}% wichtig damit die Wirkung von clip begrenzt wird
			% Schneiden und Füllen der eigentlichen Form
			  \clip[postaction={ultra thick, DodgerBlue2,fill=DodgerBlue2,fill opacity=0.2},](0,0)--(4,0)--(10,3)--(6,3)-- cycle;
			% Zeichnen der Streifen
			  \foreach \x in {0.5,1,1.5,2,2.5,...,4}%
			           \draw[DodgerBlue2!80,line width=1.3mm, opacity=0.7](0,\x)--(10,\x);
			\end{scope}
			\node[DodgerBlue2] at (7,0.6) {$P(\underline{a}_1+ \underline{a}_2, \underline{a}_2)$};
			\draw[->, ultra thick, Firebrick2] (0,0) --(4,0) node[midway, below]{$\underline{a}_2$};
			\draw[->, ultra thick, Chartreuse4] (0,0)--(2,3) node[midway, above left]{$\underline{a}_1$};
			\draw[->, ultra thick, DodgerBlue2] (0,0)-- (6,3) node[midway, above left]{$\underline{a}_1 + \underline{a}_2$};
			
		\end{tikzpicture}
	\]
	geschickte Zerlegung + Translationsinvarianz
\end{description}
% subsection 515 (end)

\subsection[Corollar: Volumen eines Quaders unter einer linearen Abbildung]{Corollar} % (fold)
\label{sub:516}
Sei $T : \mathds{R}^n \to \mathds{R}^n$ linear und $Q \subset \mathds{R}^n$ ein Quader. Dann gilt
\[
	v\big( T(Q)\big) =\abs{\det T} \cdot v(Q)
\]
\minisec{Beweis}
O.E. $Q \subset \mathds{R}^n$ abgeschlossen. Es gibt $\underline{a} \in \mathds{R}^n$, $\alpha_1, \ldots , \alpha_n \in \mathds{R}$ mit $Q = \underline{a} + P(\alpha \cdot \underline{e}_1, \ldots , \alpha_n \cdot \underline{e}_n)$. Dann ist $v(Q) = v\big(P(\alpha \cdot \underline{e}_1, \ldots , \alpha_n \cdot \underline{e}_n)\big) = \abs{\alpha_1 \cdot \ldots \cdot  \alpha_n}$ und
\begin{align*}
	T(Q) = T(\underline{a}) + T\big(P(\alpha \cdot \underline{e}_1, \ldots , \alpha_n \cdot \underline{e}_n)\big) = T(\underline{a}) 
	+ P\big(\alpha \cdot T(\underline{e}_1), \ldots , \alpha_n \cdot T(\underline{e}_n)\big)
\end{align*}
Also
\begin{align*}
	v\big(T(Q)\big) \stackrel{\text{D1}}{=} \abs{\alpha_1 \cdot \ldots \cdot \alpha_n} \cdot P\big(T(\underline{e}_1), \ldots , T(\underline{e}_n)\big) = \abs{\alpha_1 \cdot \ldots \cdot \alpha_n} \cdot \abs{\det T}   \bewende
\end{align*}
% subsection 516 (end)
% section messbarkeit_in_r_n_nullmengen (end)
\newpage
\section{$L^1(\mathds{R}^n)$-Konvergenzsätze} % (fold)
\label{sec:6}
$\mathcal{L}^1(\mathds{R}^n)$ Raum der integrierbaren Funktionen
\subsection[\protect{Proposition: Der durch $\norm[1]{f}=0$ definierte Unterraum von $\mathcal{L}^1(\mathds{R}^n)$}]{Proposition} % (fold)
\label{sub:61}
\[
	\mathcal{N} = \set[f \in \mathcal{L}^1(\mathds{R}^n)]{\norm[1]{f}=0 } 
\]
ist ein Untervektorraum von $\mathcal{L}^1(\mathds{R}^n)$.
\minisec{Beweis}
folgt aus $\norm[1]{\alpha f + \beta g} \le \abs{\alpha} \cdot \norm[1]{f} + \abs{\beta} \cdot \norm[1]{g}$
% subsection 61 (end)

\subsection[Definition: Raum der integrierbaren Funktionen als Quotientenraum]{Definition} % (fold)
\label{sub:62}
\[
	L^1(\mathds{R}^n) = \nicefrac{\mathcal{L}^1(\mathds{R}^n)}{\mathcal{N}}
\]
ist ein Quotientenvektorraum versehen mit der Norm
\[
	\norm[1]{f+ \mathcal{N}} := \norm[1]{f}  \tag{links: Norm auf $L^1$, rechts Halbnorm $\mathcal{L}^1$}
\]
Falls $f \in \mathcal{L}^1(\mathds{R}^n)$, so schreiben wir $[f]$ oder nur $f$ für $f+\mathcal{N}$.
% subsection 62 (end)

\subsection{Satz (Riesz-Fischer)} % (fold)
\label{sub:63}
$L^1(\mathds{R}^n)$ ist vollständig. Ist $([f_k])_{k \in \mathds{N}} \subset L^1(\mathds{R}^n)$ eine Cauchyfolge mit Limes $[f]$, so existiert eine Teilfolge 
$(f_{k_\nu})_{\nu \in \mathds{N}} \subset \mathcal{L}^1(\mathds{R}^n)$ mit $f_{k_\nu} \xrightarrow{\nu \to \infty} f$ punktweise fast überall und:
\[
	\lim_{ k \to \infty} \int\! f_k  \, \mathrm{d}x = \int\! f  \, \mathrm{d}x \tag{$\star$}
\] 
\minisec{Beweis}
Sei $([f_k])k \subset L^1(\mathds{R}^n)$ Cauchy, d.h. $(f_k)_k \subset \mathcal{L}^1(\mathds{R}^n)$ erfüllt:
\[
	\forall \varepsilon >0 \exists N \in \mathds{N} : \forall m,n \ge N : \norm[1]{f_m -f_n} < \varepsilon \tag{$\star\star$} 
\]
Wähle Indizes $k_0 < k_1 < \ldots $ mit
\[
	\norm[1]{f_k - f_{k_\nu}} \le 2^{-(\nu+1)} \qquad \text{ falls } k \ge k_\nu.
\]
Es gilt dann $\sum_{\nu=0}^{\infty} \norm[1]{f_{k_{\nu+1}} - f_{k_\nu}} \le 1 $. Setze 
\begin{align*}
	g_\nu &:= f_{k_{\nu+1}} - f_{k_\nu}\\
	g &:= \sum_{y=0}^{\infty} \abs{g_\nu}
\end{align*}
Dann sind $g_\nu \in \mathcal{L}^1(\mathds{R}^n), g: \mathds{R}^n \to [0,\infty]$ und $\norm[1]{g} \le 1 $ (nach \ref{sub:39}). Nach \ref{sub:58} ist $g$ fast überall 
endlich, d.h. $M= \set[x \in \mathds{R}^n]{g(x)=\infty} $ ist eine Nullmenge. Es folgt, dass die Reihe $\sum_{\nu=0}^{\infty} g_\nu (x) $ absolut konvergiert für
$x \in \mathds{R}^n \setminus M$. Definiere $f : \mathds{R}^n \to \mathds{R}$ durch
\[
	f(x) = \begin{cases}
		f_{k_0}(x) + \sum_{\nu=0}^{\infty} g_\nu (x) = \lim_{ \nu \to \infty} f_{k_\nu}(x), &\text{ falls }x \not\in M\\
		0 , &\text{ falls } x \in M
	\end{cases} \marginnote{Teleskopsumme}
\]
Noch zu zeigen: $f$ ist integrierbar, $\norm[1]{f_k -f} \xrightarrow{k \to \infty} 0 $ und ($\star$). \\
Sei $\varepsilon>0$. Wähle $\sigma \in \mathds{N}$ mit
\[
	\sum_{\nu=\sigma}^{\infty} \norm[1]{g_\nu} < \varepsilon \quad \text{ und } \quad \norm[1]{f_k - f_{k_\sigma}}< \varepsilon \enspace\forall k \ge k_\sigma 
\]
Sei $\varphi \in \mathcal{T}_{\mathds{R}^n}$ mit $\norm[1]{f_{k_\sigma} - \varphi} < \varepsilon $. Dann gilt
\begin{align*}
	\norm[1]{f- \varphi} \le \norm[1]{f- f_{k_\sigma}} + \norm[1]{f_{k_\sigma} - \varphi} \le 2 \varepsilon   
\end{align*}
Es gilt weiter
\begin{align*}
	f(x) &= f_{k_0}(x) + \sum_{\nu=0}^{\infty} g_\nu(x) = f_{k_\sigma} (x) + \sum_{\nu=\sigma}^{\infty}g_\nu (x) \tag{für $x \not\in M$} \\
	\Longrightarrow f- f_{k_\sigma} &= \sum_{\nu=\sigma}^{\infty} g_\nu \tag{auf $\mathds{R}^n \setminus M$} 
\end{align*}
$\Rightarrow f \in \mathcal{L}^1(\mathds{R}^n)$. Für $k \ge k_\sigma$ gilt
\[
	\norm[1]{f-f_k} \le \norm[1]{f- f_{k_\sigma}} + \norm[1]{f_{k_\sigma} - f_k} \le \varepsilon + \varepsilon   
\]
d.h. $\norm[1]{f-f_k}  \xrightarrow{k \to \infty} 0 $. Trivialerweise gilt $f_{k_\nu} \to f$ punktweise fast überall (für $x \not\in M$). Für $(\star)$:
\begin{align*}
	\abs{\int\! f  \, \mathrm{d}x  - \int\! f_k  \, \mathrm{d}x } \le \int\! \abs{f-f_k}  \, \mathrm{d}x  = \norm[1]{f-f_k} \xrightarrow{k \to \infty} 0 \bewende
\end{align*}
% subsection 63 (end)

\subsection[Korollar: Existenz approximierende Folge von Treppenfunktionen für integrierbare $f$]{Korollar} % (fold)
\label{sub:64}
\footnote{Ist $\norm[1]{f-g} =0 $ für $f,g \in \mathcal{L}^1(\mathds{R}^n)$, dann gilt $f(x)=g(x)$ fast überall}
Sei $f \in \mathcal{L}^1(\mathds{R}^n)$. Dann existiert $(\varphi_\nu)_{\nu \in \mathds{N}} \subset \mathcal{T}_{\mathds{R}^n}$ mit:
\begin{enumerate}[(i)]
	\item $\norm[1]{f- \varphi_\nu} \xrightarrow{\nu \to \infty}  0 $
	\item $\sum_{\nu=0}^{\infty} \norm[1]{\varphi_{\nu +1} - \varphi_\nu} < \infty $
	\item $\varphi_\nu \xrightarrow{\nu \to \infty}  f$ punktweise, fast überall
\end{enumerate}
\minisec{Beweis}
Da $f \in \mathcal{L}^1(\mathds{R}^n)$ existiert $(\psi_k)_k \subset \mathcal{T}_{\mathds{R}^n}$ mit $\norm[1]{f- \psi_k}  \xrightarrow{k \to \infty}  0$. Dann gilt:
\[
	\forall \varepsilon >0 :  \exists N \in \mathds{N} : \forall m,n \ge N : \norm[1]{\psi_n - \psi_m} < \varepsilon \tag{vgl. $(\star\star)$ in \ref{sub:63}}
\]
Nach Beweis von \ref{sub:63} existiert eine Teilfolge $(\psi_{k_\nu})_\nu$ mit $\sum_{\nu=0}^{\infty} \norm[1]{\psi_{k_{\nu+1}} - \psi_{k_\nu}} \le 1  $ und so dass
$\psi_{k_\nu} \xrightarrow{\nu \to \infty} \tilde f$ punktweise fast überall. Dann gilt $f=\tilde f$ fast überall (da $\norm[1]{f- \tilde f}=0 $). Also auch
$\psi_{k_\nu} \xrightarrow{\nu \to \infty}  f$ fast überall. \bewende
% subsection 64 (end)

\subsection[Bemerkung: Der Übergang zu einer Teilfolge in \ref{sub:63} ist wesentlich]{Bemerkung} % (fold)
\label{sub:65}
Der Übergang zu einer Teilfolge in \ref{sub:63} ist wesentlich. Betrachte dazu
\[
	f_0 : \!
	\vcenter{\hbox{\begin{tikzpicture}[scale=0.3]
		\draw (0,0) node[below]{$0$} -- (2,0)node[below]{$\frac{1}{2} $} -- (4,0) node[below]{$1$};
		\draw (0,0) -- (0,2.3);
		\node[left] at (0,2) {$1$};
		\draw[very thick, Firebrick3] (0,2)-- (2,2);
	  	\foreach \x in {0,1,2,...,4}%
	           \draw[line width=0.3mm](\x,-0.15)--(\x,0.15);
	  	
	\end{tikzpicture}}} \enspace
	f_1 : \!
	\vcenter{\hbox{\begin{tikzpicture}[scale=0.3]
		\draw (0,0) node[below]{$0$} -- (2,0)node[below]{$\frac{1}{2} $} -- (4,0) node[below]{$1$};
		\draw (0,0) -- (0,2.3);
		\node[left] at (0,2) {$1$};
		\draw[very thick, Firebrick3] (2,2)-- (4,2);
	  	\foreach \x in {0,1,2,...,4}%
	           \draw[line width=0.3mm](\x,-0.15)--(\x,0.15);
	  	
	\end{tikzpicture}}} \enspace
	f_2 : \!
	\vcenter{\hbox{\begin{tikzpicture}[scale=0.3]
		\draw (0,0) node[below]{$0$} -- (2,0)node[below]{$\frac{1}{2} $} -- (4,0) node[below]{$1$};
		\draw (0,0) -- (0,2.3);
		\node[left] at (0,2) {$1$};
		\draw[very thick, Firebrick3] (0,2)-- (1,2);
	  	\foreach \x in {0,1,2,...,4}%
	           \draw[line width=0.3mm](\x,-0.15)--(\x,0.15);
	  	
	\end{tikzpicture}}} \enspace
	f_3 : \!
	\vcenter{\hbox{\begin{tikzpicture}[scale=0.3]
		\draw (0,0) node[below]{$0$} -- (2,0)node[below]{$\frac{1}{2} $} -- (4,0) node[below]{$1$};
		\draw (0,0) -- (0,2.3);
		\node[left] at (0,2) {$1$};
		\draw[very thick, Firebrick3] (1,2)-- (2,2);
	  	\foreach \x in {0,1,2,...,4}%
	           \draw[line width=0.3mm](\x,-0.15)--(\x,0.15);
	  	
	\end{tikzpicture}}} \enspace
	f_4 : \!
	\vcenter{\hbox{\begin{tikzpicture}[scale=0.3]
		\draw (0,0) node[below]{$0$} -- (2,0)node[below]{$\frac{1}{2} $} -- (4,0) node[below]{$1$};
		\draw (0,0) -- (0,2.3);
		\node[left] at (0,2) {$1$};
		\draw[very thick, Firebrick3] (2,2)-- (3,2);
	  	\foreach \x in {0,1,2,...,4}%
	           \draw[line width=0.3mm](\x,-0.15)--(\x,0.15);
	  	
	\end{tikzpicture}}} \enspace
	f_5 : \!
	\vcenter{\hbox{\begin{tikzpicture}[scale=0.3]
		\draw (0,0) node[below]{$0$} -- (2,0)node[below]{$\frac{1}{2} $} -- (4,0) node[below]{$1$};
		\draw (0,0) -- (0,2.3);
		\node[left] at (0,2) {$1$};
		\draw[very thick, Firebrick3] (3,2)-- (4,2);
	  	\foreach \x in {0,1,2,...,4}%
	           \draw[line width=0.3mm](\x,-0.15)--(\x,0.15);
	  	
	\end{tikzpicture}}} \enspace
\]
\ldots Dann gilt: $\norm[1]{f_k -0} \xrightarrow{k \to \infty} 0 $, aber es gibt \emph{kein} $x \in [0,1]$ mit $f_k(x) \xrightarrow{k \to \infty} 0$. Wie sieht eine Teilfolge
$(f_{k_\nu})_{\nu \in \mathds{N}}$ aus mit $f_{k_\nu} \xrightarrow{\nu \to \infty} 0$ fast überall?
% subsection 65 (end)

\subsection{Satz (Beppo-Levi)} % (fold)
\label{sub:66}
Sei $(f_k)_{k \in \mathds{N}}$ eine monoton wachsende Folge integrierbarer Funktionen. $f_k : \mathds{R}^n \to (-\infty, \infty]$, so dass 
$(\int\! f_k  \, \mathrm{d}x )_{k \in \mathds{N}} \subset \mathds{R}$ beschränkt ist. Sei $f: \mathds{R}^n \to (-\infty, \infty]$ der punktweise Limes, d.h. $f(x) = \lim_{ k \to \infty} f_k(x)$. Dann ist $f$ integrierbar und 
\[
	 \lim_{ k \to \infty} \int\! f_k  \, \mathrm{d}x = \int\! f  \, \mathrm{d}x 
\]
\minisec{Beweis}
Es ist $f_k \le f_{k+1}$, also $\int\!f_k  \, \mathrm{d}x  \le \int\! f_{k+1}  \, \mathrm{d}x $, d.h. die Folge $(\int\! f_k  \, \mathrm{d}x )_{k \in \mathds{N}} \subset \mathds{R}$ ist monoton wachsend und beschränkt, also konvergent. Also gilt auch
\[
	\forall \varepsilon >0 : \exists N \in \mathds{N}: \forall m,k \ge N : \abs{\int\! f_k  \, \mathrm{d}x - \int\! f_m  \, \mathrm{d}x } < \varepsilon  
\]
Falls $m \ge k \ge N$, dann
\begin{align*}
	\norm[1]{f_m - f_k} = \int\! \abs{f_m - f_k}  \, \mathrm{d}x  = \int\! (f_m -f_k)  \, \mathrm{d}x  = \int\! f_m  \, \mathrm{d}x  - \int\! f_k  \, \mathrm{d}x \le
	\abs{\int\! f_m  \, \mathrm{d}x - \int\! f_k  \, \mathrm{d}x }   
\end{align*}
Damit folgt: $([f_k])_k \subset L^1(\mathds{R}^n)$ ist Cauchy. Sei $\tilde f \in \mathcal{L}^1(\mathds{R}^n)$ eine Funktion mit $\norm[1]{f_k- \tilde f} \xrightarrow{k \to \infty}  0 $ 
(existiert nach \ref{sub:63}). Nach \ref{sub:63} gibt es eine Teilfolge $(f_{k_\nu})_\nu$ mit $f_{k_\nu} \xrightarrow{\nu \to \infty}  \tilde f$ fast überall und 
\[
	\lim_{ \nu \to \infty} \int\! f_{k_\nu}  \, \mathrm{d}x  = \int\!\tilde f  \, \mathrm{d}x 
\]
$\Rightarrow f=\tilde f$ fast überall. Daher $f \in \mathcal{L}^1(\mathds{R}^n)$ und 
\begin{align*}
	\int\! f  \, \mathrm{d}x = \int\! \tilde f  \, \mathrm{d}x = \lim_{ \nu \to \infty} \int\! f_{k_\nu}  \, \mathrm{d}x = \lim_{ k \to \infty} \int\! f_k  \, \mathrm{d}x 
	\bewende 
\end{align*}
% subsection 66 (end)

\subsection[Corollar: Messbarkeit abzählbarer Vereinigungen messbarer Mengen]{Corollar} % (fold)
\label{sub:67}
\begin{enumerate}[a)]
\item Seien $A_1 \subset A_1 \subset \ldots  \subset \mathds{R}^n$ messbar. Dann ist $\bigcup_{k=0}^\infty A_k$ messbar genau dann, wenn die Folge 
$\big(v(A_k)\big)_{k \in \mathds{N}}$ beschränkt ist;
in diesem Fall ist 
\[
	v \enbrace{\bigcup_{k=0}^\infty A_k} = \sup_k v(A_k)
\]
\item Seien $B_0, B_1, \ldots  \subset \mathds{R}^n$ messbar so dass $B_i \cap B_j$ Nullmengen sind für $i \not= j$. Dann ist $\bigcup_{k \in \mathds{N}} B_k$ messbar genau
dann wenn $\sum_{k=0}^{\infty} v(B_k) < \infty$;  in diesem Fall ist 
\[
	v \enbrace{\bigcup_{k \in\mathds{N}} B_k} = \sum_{k=0}^{\infty} v(B_k) 
\]
\end{enumerate}
\minisec{Beweis}
\begin{enumerate}[a)]
	\item Beppo-Levi mit $f_k := \chi_{A_k}$,
	\[
		f = \lim_{ k \to \infty} f_k = \chi_{\bigcup_{k \in \mathds{N}} A_k}
	\]
	\item folgt aus a) mit $A_k := \bigcup_{i=0}^k B_i$ und
	\[
		v(A_k) = \sum_{i=0}^{k} v(B_i) \bewende
	\]
\end{enumerate}
% subsection 67 (end)

\subsection[Bemerkung: Zusammenfassung über Eigenschaften des Lebesgue-Maßes]{Bemerkung} % (fold)
\label{sub:68}
Setze $\Lambda := \set[E]{E \subset \mathds{R}^n \text{ messbar}} \subset \mathcal{P}(\mathds{R}^n)$. Die Abbildung $v :  \Lambda \to \mathds{R}_+$ heißt auch
\Index{Lebesgue-Maß} und es gilt:
\begin{description}
	\item[(M0)] $\Lambda$ enthält alle kompakten und beschränkten offenen Mengen. 
	\item[(M1)] $A,B \in \Lambda$, $A \subseteq B  \Longrightarrow v(A) \le v(B)$.
	\item[(M2)] Für $a \in \mathds{R}^n$, $E \in \Lambda$ ist $a+ E \in \Lambda$ und $v(a+ E)= v(E)$.
	\item[(M3)] Für $B_0, B_1, \ldots  \in \Lambda$ paarweise disjunkt mit $\sum_{k \in \mathds{N}} v(B_k) < \infty$ ist $\bigcup_{k \in \mathds{N}} B_k \in \Lambda$ und
	es gilt
	\[
		v \enbrace{\bigcup_{k \in \mathds{N}} B_k} = \sum_{k \in \mathds{N}} v(B_k).
	\]
	\item[(M4)] $v(\emptyset)=0$ und $v([0,1]^n)=1$.
\end{description} 
Man kann zeigen: Das Lebesgue-Maß ist durch diese Eigenschaften eindeutig bestimmt.
% subsection 68 (end)

\subsection{Satz von Lebesgue über majorisierte Konvergenz} % (fold)
\label{sub:69}
Sei $f_k :\mathds{R}^n \to (-\infty, \infty], k \in \mathds{N}$ eine Folge integrierbarer Funktionen. Sei $f : \mathds{R}^n \to (-\infty, \infty]$ eine weitere Funktion mit 
$f_k \xrightarrow{k \to \infty} f$ punktweise fast überall. Sei $F :\mathds{R}^n \to [0, \infty]$ integrierbar mit $F\ge \abs{f_k}, k \in \mathds{N} $. Dann ist $f$
integrierbar und es gilt
\[
	\int\! f  \, \mathrm{d}x = \lim_{ k \to \infty} \int\! f_k  \, \mathrm{d}x 
\]
\minisec{Beweis}
Mittels Proposition \ref{sub:58} dürfen wir annehmen, dass $f_k, F$ nur Werte in $\mathds{R}$ annehmen\marginnote{warum?}. Wir dürfen außerdem annehmen, dass $f_k \to f$
punktweise überall.

Setze 
\[
	g_{k, \nu} := \max \set{f_k, \ldots , f_{k+\nu}} \qquad, \qquad g_k := \sup_{i \ge k} f_i \tag*{$k, \nu \in \mathds{N}$}
\]
Dann sind $g_{k, \nu}, g_k : \mathds{R}^n \to \mathds{R}$ Funktionen mit  $g_{k, \nu} \xrightarrow{\nu \to \infty} g_k$ punktweise monoton wachsend. Die $g_{k, \nu}$ sind
integrierbar und $\int\! g_{k, \nu}  \, \mathrm{d}x  \le \int\! F  \, \mathrm{d}x < \infty$. Mit Beppo-Levi (\ref{sub:66}) folgt: $g_k $ ist integrierbar und es gilt
\[
	\int\! g_k  \, \mathrm{d}x = \lim_{ \nu \to \infty} \int\! g_{k, \nu}  \, \mathrm{d}x 
\]
Weiter gilt $\abs{\int\! g_{k, \nu}  \, \mathrm{d}x } \le \int\! F  \, \mathrm{d}x $, also auch $\abs{\int\! g_k  \, \mathrm{d}x } \le \int\! F  \, \mathrm{d}x  $, $k \in \mathds{N}$. Es gilt $g_k \xrightarrow{k \to \infty} f$ punktweise monoton fallend, also $- g_k \xrightarrow{k \to \infty} -f$ punktweise monoton wachsend. Da 
\[
	\int\!-g_k  \, \mathrm{d}x  \le \abs{\int\! g_k  \, \mathrm{d}x} \le \abs{\int\! F  \, \mathrm{d}x} < \infty 
\]
ist $(\int\! -g_k  \, \mathrm{d}x )_{k \in \mathds{N}}$ beschränkt und wieder nach Beppo-Levi (\ref{sub:66}) ist $-f$ integrierbar mit 
\[
	\int\! -f  \, \mathrm{d}x = \lim_{ k \to \infty} \int\! g_k  \, \mathrm{d}x \enspace\Longrightarrow \enspace \int\! f  \, \mathrm{d}x  = \lim_{k \to \infty} \int\! g_k  \, \mathrm{d}x \tag{$\star$}
\]
Definiere nun 
\[
	h_{k, \nu} := \min \set{f_k , \ldots , f_{k+\nu}} \qquad , \qquad h_k := \inf_{i \ge k} f_i \tag*{$k, \nu \in \mathds{N}$}
\]
Dann gilt $h_{k, \nu} \xrightarrow{\nu \to \infty} h_k$ punktweise monoton fallend und $h_k \xrightarrow{k \to \infty} f$ punktweise monoton wachsend. Wie oben zeigt man
\[
	\int\! f  \, \mathrm{d}x = \lim_{k \to \infty} \int\! h_k  \, \mathrm{d}x \tag{$\star\star$}
\]
Wegen $h_k \le f_k \le g_k$ gilt 
\[
	\int\! f  \, \mathrm{d}x =\lim_{ k \to \infty} \int\! h_k  \, \mathrm{d}x  \le \lim_{ k \to \infty}  \int\! f_k  \, \mathrm{d}x  \ge \lim_{ k \to \infty} \int\! g_k  \, \mathrm{d}x = \int\! f  \, \mathrm{d}x  \bewende
\]
% subsection 69 (end)

\subsection[Corollar: Der Limes einer beschränkten Folge integrierbarer Funktionen ist integrierbar]{Corollar} % (fold)
\label{sub:610}
Sei $A \subset \mathds{R}^n$ messbar, $f_k : A  \to \mathds{R}$, $k \in \mathds{N}$ beschränkte Folge integrierbarer Funktionen. Falls $(f_k)_{k \in \mathds{N}}$ punktweise
fast überall gegen $f: A \to \mathds{R}$ konvergiert, so ist $f$ integrierbar und es gilt 
\[
	\int_A\! f  \, \mathrm{d}x = \lim_{ k \to \infty} \int_A\! f_k  \, \mathrm{d}x 
\]
\minisec{Beweis}
Mit \ref{sub:69} und $\chi_A$ mal der Konstanten aus der Beschränktheit als Majorante. \bewende
% subsection 610 (end)

\subsection[Corollar: Der Hauptsatz der Integralrechnung gilt auch für das Lebesgue-Integral]{Corollar} % (fold)
\label{sub:611}
Sei $f: [a,x] \to \mathds{R}$ differenzierbar mit beschränkter Ableitung $f'$. Dann ist $f'$ Lebesgue-integrierbar und so gilt
\[
	f(x) - f(a) = \int_{[a,x]}\! f'  \, \mathrm{d}x 
\]
\minisec{Beweis}
\begin{itemize}
	\item $f$ ist differenzierbar, also stetig. Definiere $f_k : [a,x] \to \mathds{R}$ durch 
	\[
		f_k(t) := \begin{cases}\marginnote{vgl. Differenzenquotient}
			(k+1)\cdot \Big(f\big(t+ \frac{1}{k+1}\big)- f(t)\Big), &\text{ falls }t \in [a,x-\frac{1}{k+1} ]\\
			0 ,&\text{ sonst}
		\end{cases}
	\]
	Dann gilt $f_k \xrightarrow{k \to \infty}  f'$ punktweise auf $[a,x)$. 
	\item $(f_k)_{k \in \mathds{N}}$ ist beschränkt: Nach Mittelwertsatz existiert ein $s \in (t,t+ \frac{1}{k+1})$ mit 
	\[
		\frac{f \big(t+ \frac{1}{k+1} \big) - f(t)}{\frac{1}{k+1} } = f'(s),
	\]
	also gilt $\abs{f_k(t)} \le \sup \set[\abs{f'(r)}]{r \in [a,x]} < \infty  $ (dieses Supremum hängt nicht von $k$ oder $t$ ab).
	\item $f_k$ ist integrierbar: $f_k$ ist stetig auf $[a,x- \frac{1}{k+1} ]$ und konstant 0 auf $(x-\frac{1}{k+1},x ]$, also Regelfunktion (warum?). Es gilt
	\[
		\int_{a} ^{x} \! f_k(t)  \, \mathrm{d}t =  \frac{1}{\frac{1}{k+1} } \cdot \enbrace{\int_{x-\frac{1}{k+1} } ^{x} \! f(t)  \, \mathrm{d}t - \int_{a} ^{a+ 
		\frac{1}{k+1} } \! f(t)  \, \mathrm{d}t}.
	\]
	MWS der Integralrechnung (AnaI, Satz 13.13) liefert: Es existiert ein $\overline{t}_k \in [x- \frac{1}{k+1},x]$ mit 
	\[
		f(\overline{t}_k ) \cdot \frac{1}{k+1} = \int_{x- \frac{1}{k+1} } ^{x} \! f(t)  \, \mathrm{d}t
	\]
	entsprechend für $\overline{s}_k \in[a,a+ \frac{1}{k+1} ]$. $f$ stetig $\Longrightarrow$
	\begin{align*}
		\int_{a} ^{x} \! f_k(t)  \, \mathrm{d}t &= \frac{1}{\frac{1}{k+1} } \cdot \enbrace{\int\! \ldots  - \int\! \ldots } = f(\overline{t}_k )- f(\overline{s}_k )
		\xrightarrow{k \to \infty} f(x)- f(a)
	\end{align*}
	mit \ref{sub:610} folgt: $f'$ ist Lebesgue-integrierbar und es gilt
	\[
		\int_{[a,x]}  \! f'  \, \mathrm{d}x  = \lim_{ k \to \infty} \int_{[a,x]} f_k  \, \mathrm{d}t = \lim_{ k \to \infty} \int_{a} ^{x} \! f_k(t)  \, \mathrm{d}t
		= f(x)- f(a) \bewende
	\]
\end{itemize}
% subsection 611 (end)

\subsection[Definition: $\sigma$-kompakt]{Definition} % (fold)
\label{sub:612}
$A \subset \mathds{R}^n$ heißt \bet{$\sigma$-kompakt}\index{Sigma-kompakt@$\sigma$-kompakt}, falls $A$ abzählbare Vereinigung kompakter Mengen ist.
% subsection 612 (end)

\subsection[Beispiel: $\sigma$-kompakte Mengen]{Beispiel} % (fold)
\label{sub:613}
Folgende Mengen sind alle $\sigma$-kompakt:
\begin{multicols}{2}
	\begin{itemize}
		\item offene Mengen
		\item abgeschlossene Mengen
		\item endliche Durchschnitte $\sigma$-kompakter Mengen
		\item abzählbare Vereinigungen $\sigma$-kompakter Mengen
	\end{itemize}
\end{multicols}
% subsection 613 (end)

\subsection[Definition: lokal integrierbar]{Definition} % (fold)
\label{sub:614}
Sei $A \subseteq \mathds{R}^n$ $\sigma$-kompakt. Dann heißt $f : A \to (-\infty, \infty]$ \Index{lokal integrierbar}, falls für jedes kompakte $K \subset A$ die Funktion
$f|_K : K \to (-\infty,\infty]$ integrierbar ist.
% subsection 614 (end)

\subsection{Corollar (Majorantenkriterium)} % (fold)
\label{sub:615}
Sei $A \subset \mathds{R}^n$ $\sigma$-kompakt, $f : A \to (-\infty, \infty]$ lokal integrierbar. Sei $F : A \to (-\infty, \infty]$ integrierbar mit $F\ge \abs{f} $.
Dann ist $f$ integrierbar über $A$.
\minisec{Beweis}
Sei $A= \bigcup_{k \in \mathds{N}} A_k$ mit $A_k \subseteq A_{k+1}$ kompakt für alle $k \in \mathds{N}$. Dann ist $f|_{A_k}$ integrierbar, also auch 
\[
	f_k := f_A \cdot \chi_{A_k} \enspace{\color{light_gray} = (f|_{A_k})_{A_k}}
\]
es gilt dann $f_k \xrightarrow{k \to \infty}  f$ punktweise und $\abs{f_k} \le F $. Mit Satz \ref{sub:69} folgt, dass $f$ integrierbar ist. \bewende
% subsection 615 (end)
% section 6 (end)
\newpage

\section{Der Satz von Fubini} % (fold)
\label{sec:7}
\subsection[Lemma über Nullmengen]{Lemma} % (fold)
\label{sub:71}
Sei $A \subset \mathds{R}^p \times \mathds{R}^q$ eine Nullmenge. Dann existiert eine Nullmenge $B \subset \mathds{R}^q$ so, dass für jedes $y \in \mathds{R}^q\setminus B$
die Menge 
\[
	A_y := \set[x \in \mathds{R}^p]{(x,y) \in A} \subset \mathds{R}^p
\]
eine Nullmenge im $\mathds{R}^p$ ist.
\minisec{Beweis}
$A$ ist Nullmenge, d.h. $\norm[1]{\chi_A} =0$. Sei $\varepsilon>0$. Dann existiert eine Treppe\footnote{$c_k \in \mathds{R}_+$, $Q_k$ offene Quader in $\mathds{R}^p \times \mathds{R}^q$}
$\Phi = \sum_{k=0}^{\infty} c_k \cdot \chi_{Q_k}$ über $\chi_A$ mit $I(\Phi ) = \sum_{k=0}^{\infty} c_k v(Q_k) < \varepsilon$. Es ist $Q_k = Q_k' \times Q_k''$, wo 
$Q_k' \subset \mathds{R}^p$ und $Q_K'' \subset \mathds{R}^q$ offene Quader sind. Definiere
$a : \mathds{R}^q \to (-\infty, \infty]$ durch 
\[
	a(y) := \norm[1, \mathds{R}^p]{\chi_{A_y}} 
\]
Für $x \in \mathds{R}^p$ gilt 
\[
	\chi_{A_y}(x)  \le \chi_A(x,y) \le \Phi (x,y) = \sum_{k=0}^{\infty} c_k \cdot \chi_{Q_k'}(x) \cdot \chi_{Q_k''}(y) 
\]
also $\chi_{A_y} \le \sum_{k=0}^{\infty} c_k \cdot \chi_{Q_k''}(y) \cdot \chi_{Q_k'}$ und nach \ref{sub:39} gilt
\[
	a(y) = \norm[1]{\chi_{A_y}} \le \sum_{k=0}^{\infty} c_k \cdot \chi_{Q_k''}(y) \cdot \norm[1,\mathds{R}^p]{\chi_{Q_k'}}   = \sum_{k=0}^{\infty} c_k \cdot v(Q_k') \cdot 
	\chi_{Q_k''}(y),
\]
somit $a \le \sum_{k=0}^{\infty} c_k \cdot c(Q_k') \cdot \chi_{Q_k''} =: \Psi$. $\Psi$ ist Treppe über $a$. Aber 
\[
	\norm[1,\mathds{R}^q]{a} \le I(\Psi) = \sum_{k=0}^{\infty} c_k \cdot v(Q_k') \cdot v(Q_k'') = I(\Phi) <\varepsilon 
\]
also gilt $\norm[1,\mathds{R}^q]{a}=0 $. Für $l \in \mathds{N}$ setze $B_l := \set[y \in \mathds{R}^q]{a(y) \ge \frac{1}{l+1} } $. Dann ist $B_l \subset \mathds{R}^q$
Nullmenge (warum?). Daraus folgt mit \ref{sub:57}
\[
	B := \bigcup_{l \in \mathds{N}} B_l \text{ Nullmenge}
\]
aber $B= \set[y \in \mathds{R}^q]{a(y)>0}$ also  
\[
	\forall y \in \mathds{R}^q \setminus B : 0 = a(y) = \norm[1]{\chi_{A_y}} 
\]
$\Rightarrow A_y$ Nullmenge. \bewende
% subsection 71 (end)

\subsection[Bemerkung, warum $B$ in \ref{sub:71} ausgeschlossen werden muss]{Bemerkung} % (fold)
\label{sub:72}
In \ref{sub:71} ist es wesentlich, die Nullmenge $B$ auszuschließen:
\[
	\begin{tikzpicture}
		%\draw[help lines] (0,0) grid (3,2);
		\draw[<->, thick] (0,2) node[left]{1} -- (0,0) -- (3,0);
		\draw[ultra thick, DodgerBlue2] (1,0)node[below]{$y$} -- (1,2) node[right]{$A_y$};
		\draw[Salmon3] (1,0) -- (1,2);
		\draw[Salmon3] (0.2,0.4) --(2.3,1.3);
		\draw[Salmon3] (0.6,1.7) --(0.3,0.3);
		\draw[Salmon3] (2,1.3) --(0.9,1.9);
		\draw[Salmon3] (3,0.5) --(0.4,1.1);
		\draw[Salmon3] (0.2,1.6) --(1.3,0.1);
		\draw[Salmon3] (0.1,1.1) --(2.7,1.8);
		\draw[Salmon3] (1.6,0.15) --(2.9,1.8) node[right, midway]{$A$};
	\end{tikzpicture}
\]
$A \subset \mathds{R}^2$ ist eine Nullmenge, aber $A_y = [0,1] \subset \mathds{R}$ ist keine Nullmenge!
% subsection 72 (end)

\subsection{Satz von Fubini} % (fold)
\label{sub:73}
Sei $f : \mathds{R}^p \times \mathds{R}^q \to (-\infty, \infty]$ integrierbar. 
\begin{enumerate}[a)]
	\item Die Funktion $f_y : \mathds{R}^p \to (-\infty, \infty]$ mit $f_y(x) := f(x,y)$ ist für fast alle $y \in \mathds{R}^q$ integrierbar, d.h. $f_y$ ist integrierbar für
	$y \in \mathds{R}^q \setminus N$, für eine Nullmenge $N \subset \mathds{R}^q$.
	\item Definiere $F : \mathds{R}^q \to (-\infty,\infty)$ durch
	\[
		F(y) := \begin{cases}
			\int_{\mathds{R}^p}\! f_y  \, \mathrm{d}x , &\text{ falls }y \in \mathds{R}^q \setminus N\\
			0 ,&\text{ falls } y \in N
		\end{cases}
	\]
	Dann ist $F$ integrierbar und es gilt
	\[
		\int_{\mathds{R}^p \times \mathds{R}^q}\! f(x,y)  \, \mathrm{d}(x,y) = \int_{R^q} F(y)  \, \mathrm{d}y
	\]
	"{}Kürzer"{}:
	\[
		\int_{\mathds{R}^p \times \mathds{R}^q}\! f(x,y)  \, \mathrm{d}(x,y) = \int_{\mathds{R}^q}\! \enbrace{\int_{\mathds{R}^p}\! f_y(x)  \, \mathrm{d}x }  
		\, \mathrm{d}y 
	\]
\end{enumerate}
\minisec{Beweis}
\begin{enumerate}[a)]
	\item Nach Corollar \ref{sub:64} existiert $(\varphi_k)_{k \in \mathds{N}} \subset \mathcal{T}_{\mathds{R}^p \times \mathds{R}^q}$ und eine Nullmenge 
	$A \subset \mathds{R}^p \times \mathds{R}^q$ mit $\norm[1]{f- \varphi_k} \xrightarrow{k \to \infty} 0 $ und 
	\begin{enumerate}[(i)]
		\item $\varphi_k(x,y) \xrightarrow{k \to \infty} f(x,y)$ für $(x,y) \in \mathds{R}^p \times \mathds{R}^q \setminus A$
		\item $ \sum_{k=0}^{\infty} \norm[1]{\varphi_{k+1} - \varphi_k} < \infty $
	\end{enumerate}
	Nach Lemma \ref{sub:71} existiert eine Nullmenge $N' \subset \mathds{R}^q$ so dass für $y \in \mathds{R}^q\setminus N'$ gilt:
	\[
		A_y := \set[x \in \mathds{R}^p]{(x,y) \in A} \subset \mathds{R}^p 
	\]
	ist eine Nullmenge. Aber für $x \in \mathds{R}^p \setminus A_y$ gilt $(x,y) \not\in A$, also 
	\begin{enumerate}[(i)]
		\setcounter{enumii}{2}
		\item Für $y \in \mathds{R}^q \setminus N'$ gilt
		\[
			{\color{light_gray}(x \mapsto \varphi_k (x,y) )} \enspace(\varphi_k)_y \xrightarrow{k \to \infty} f_y \quad \text{ punktweise fast überall, d.h. für } x \in 
			\mathds{R}^p \setminus A_y
		\]
	\end{enumerate}
	Setze 
	\[
		H_k(y) := \int_{\mathds{R}^p}\! \underbrace{\abs{\varphi_{k+1}(x,y) - \varphi_k(x,y)}}_{\in \mathcal{T}(\mathds{R}^p \times \mathds{R}^q)}  \, \mathrm{d}x   
		= \int_{\mathds{R}^p}\! \underbrace{\abs{(\varphi_{k+1})_y (x) - 
		(\varphi_k)_y (x)}}_{\in \mathcal{T}(\mathds{R}^p)}  \, \mathrm{d}x  = \norm[1,\mathds{R}^p]{(\varphi_{k+1})_y - (\varphi_k)_y} 
	\]
	Dann gilt nach Fubini für Treppenfunktionen(\ref{sub:34}):
	\begin{align*}
		\int_{\mathds{R}^q} H_k(y)  \, \mathrm{d}y = \int_{\mathds{R}^q}\! \enbrace{\int_{\mathds{R}^p} \abs{\varphi_{k+1} (x,y)- \varphi_k(x,y)}  \, \mathrm{d}x  }  
		\, \mathrm{d}y &= \int_{\mathds{R}^p \times \mathds{R}^q}\! \abs{\varphi_{k+1} - \varphi_k}  \, \mathrm{d}(x,y)\\
		&= \norm[1,\mathds{R}^p \times \mathds{R}^q]{\varphi_{k+1} - \varphi_k}
	\end{align*}
	Nach (ii) gilt dann
	\[
		\sum_{k=0}^{\infty} \int_{\mathds{R}^q}\! H_k(y)  \, \mathrm{d}y < \infty \tag{$\star$}
	\]
	Nun ist $\enbrace{\sum_{k=0}^{i} H_k}_{i \in \mathds{N}}$ eine monoton wachsende Folge von integrierbaren Funktionen; die Folge der Integrale 
	$\enbrace{\int_{\mathds{R}^q} \enbrace{\sum_{k=0}^{i} H_k}  \, \mathrm{d}y }_{i \in \mathds{N}} $ ist beschränkt durch 
	$\sum_{k=0}^{\infty} \int_{\mathds{R}^q}\! H_k(y)  \, \mathrm{d}y \stackrel{(\star)}{<} \infty$. nach Beppo-Levi (\ref{sub:66}) ist dann
	\[
		H := \sum_{k=0}^{\infty} H_k : \mathds{R}^q \to (-\infty, \infty]
	\]
	integrierbar und nach \ref{sub:58} ist $H$ außerhalb einer Nullmenge $N'' \subset \mathds{R}^q$ endlich. Wir erhalten
	\begin{enumerate}[(i)]
		\setcounter{enumii}{3}
		\item $\sum_{k=0}^{\infty} \norm[1, \mathds{R}^p]{(\varphi_{k+1})_y - (\varphi_k)_y}  = \sum_{k=0}^{\infty} H_k(y) = H(y) < \infty \quad$ für 
		$y \in \mathds{R}^q \setminus N''$
	\end{enumerate}
	Setze $N:= N' \cup N'' \subset \mathds{R}^q$.
	
	Für $y \in \mathds{R}^q \setminus N$ gilt nach (iv): $\big((\varphi_k)_y\big)_{k \in \mathds{N}}$ ist Cauchy-Folge bezüglich $\norm[1, \mathds{R}^p]{.} $ also bildet
	\[
		\Big(\big[(\varphi_k)_y\big]\Big)_{k \in \mathds{N}} \subset L^1(\mathds{R}^p)
	\]
	eine Cauchy-Folge und konvergiert nach \ref{sub:63} in $L^1(\mathds{R}^p)$ gegen ein $[g] \in L^1(\mathds{R}^p)$ für ein $g \in \mathcal{L}^1(\mathds{R}^p)$.
	Ebenfalls nach \ref{sub:63} gilt für eine Teilfolge $\big((\varphi_{k_\nu})_y\big)_{k \in \mathds{N}}$ mit $(\varphi_{k_\nu})_y \xrightarrow{\nu \to \infty} g$
	punktweise fast überall.
	
	Nach (iii) ist dann $g=f_y$ fast überall, also ist $f_y$ integrierbar für $y \in \mathds{R}^q \setminus N$.
	\item Nach Satz \ref{sub:63} gilt außerdem für $y \in \mathds{R}^q \setminus N$
	\[
		\lim_{ k \to \infty} \int_{\mathds{R}^p}\! \varphi_k(x,y)  \, \mathrm{d}x = \int_{\mathds{R}^p} f(x,y)  \, \mathrm{d}x \tag{$\star\star$}
	\]
	Wir setzen $\Phi_k (y) := \int_{\mathds{R}^p}\! \varphi_k(x,y) \, \mathrm{d}x $, $k \in \mathds{N}$. Dann ist $\Phi_k \in \mathcal{T}_{\mathds{R}^q}$ (vgl. Fubini für
	Treppenfunktionen, \ref{sub:34}) und es gilt 
	\begin{enumerate}[(i)]
		\setcounter{enumii}{4}
		\item $\Phi_k \xrightarrow{k \to \infty} F$ punktweise auf $\mathds{R}^q \setminus N \quad $ (nach $(\star\star)$)
		\item $ \sum_{k=0}^{\infty} \norm[1, \mathds{R}^q]{\Phi_{k+1} - \Phi_k} \le \sum_{k=0}^{\infty} \norm[1, \mathds{R}^q]{H_k} = \sum_{k=0}^{\infty} \norm[1, 
		\mathds{R}^q \times \mathds{R}^p]{\varphi_{k+1}- \varphi_k} \stackrel{\text{(ii)}}{<} \infty   $.
	\end{enumerate}
	Wegen (vi) ist $(\Phi _k)_{k \in \mathds{N}}$ Cauchy bzgl. $\norm[1,\mathds{R}^q]{.}$, also konvergiert nach \ref{sub:63} eine Teilfolge 
	$(\Phi _{k_\nu})_{\nu \in \mathds{N}}$ punktweise fast überall gegen ein $G \in \mathcal{L}^1(\mathds{R}^q)$. Wegen (v) gilt $F=G$ fast überall und $F$ ist 
	integrierbar.
	
	Wieder nach \ref{sub:63} gilt 
	\begin{align*}
		\int_{R^q}\! F(y)  \, \mathrm{d}y \stackrel{\text{\ref{sub:63}}}{=} \lim_{ k \to \infty} \int_{\mathds{R}^q}\! \Phi_k(y)  \, \mathrm{d}y
		&\stackrel{\text{\ref{sub:34}}}{=} \lim_{ k \to \infty} \int_{\mathds{R}^p \times \mathds{R}^q}\! \varphi_k(x,y)  \, \mathrm{d}(x,y) \\
		&\stackrel{\text{\ref{sub:42}}}{=} \int_{\mathds{R}^p \times \mathds{R}^q}\! f(x,y)  \, \mathrm{d}(x,y) \bewende
	\end{align*}
\end{enumerate}
% subsection 73 (end)

\subsection{Satz von Tonelli} % (fold)
\label{sub:74}
Sei $f : \mathds{R}^p \times \mathds{R}^q \to (-\infty, \infty]$ lokal integrierbar. $f$ ist genau dann integrierbar, wenn wenigstens eines der Integrale
\[
	\int_{\mathds{R}^q}\! \enbrace{\int_{\mathds{R}^p} \!\abs{f(x,y)}  \, \mathrm{d}x }  \, \mathrm{d}y \enspace(\star)\quad , \quad \int_{\mathds{R}^p}\! \enbrace{\int_{\mathds{R}^q} 
	\! \abs{f(x,y)}  \, \mathrm{d}y }  \, \mathrm{d}x \enspace (\star\star)
\]
existiert. In diesem Fall gilt die Formel aus \ref{sub:73}.
\minisec{Beweis}
\begin{description}
	\item["$\Rightarrow$":] $f$ integrierbar $\Rightarrow $ $\abs{f} $ integrierbar; Existenz von $(\star)$ und $(\star\star)$ folgt aus \ref{sub:73}.
	\item["$\Leftarrow$":] Nach \ref{sub:615} genügt es zu zeigen, dass $\abs{f} $ integrierbar ist. Sei $W_k := [-k,k]^{p} \times [-k,k]^q$; setze
	\[
		f_k := \min \set{\abs{f}, \abs{f} \cdot  \chi_{W_k} }, k \in \mathds{N} 
	\]
	Dann konvergiert $(f_k)_{k \in \mathds{N}}$ punktweise monoton wachsend gegen $\abs{f}$. Außerdem gilt 
	\begin{align*}
		0 \le\int_{\mathds{R}^p \times \mathds{R}^q}\! f_k(x,y)  \, \mathrm{d}(x,y) \stackrel{\text{\ref{sub:73}}}{=} \int_{\mathds{R}^q}\! 
		\enbrace{ \int_{\mathds{R}^p}\! f_k(x,y)  \, \mathrm{d}x }  \, \mathrm{d}y \le \int_{\mathds{R}^q}\! 
		\enbrace{ \int_{\mathds{R}^p}\! \abs{f(x,y)}  \, \mathrm{d}x  }  \, \mathrm{d}y \stackrel{\text{falls }(\star) \text{ ex.}}{<} \infty 
	\end{align*}
	$\Rightarrow \enbrace{\int_{\mathds{R}^p \times \mathds{R}^p} f_k(x,y)  \, \mathrm{d}(x,y)}_{k \in \mathds{N}} $ beschränkt. Mit Beppo-Levi (\ref{sub:66}) folgt
	$\abs{f} $ ist integrierbar. \bewende
\end{description}
% subsection 74 (end)

\subsection[Beispiel: Anwendung von Fubini]{Beispiel} % (fold)
\label{sub:75}
\[
	f(x,y) := \begin{cases}
		\frac{x-y}{(x+y)^3} , &\text{ falls }x+y\not= 0\\
		0 ,&\text{ sonst } 
	\end{cases}
\]
ist nicht integrierbar über $[0,1] \times [0,1]$, und 
\[
	\int_{0} ^{1} \! \enbrace{ \int_{0} ^{1} \! f(x,y)  \, \mathrm{d}x }  \, \mathrm{d}y \not= \int_{0} ^{1} \! \enbrace{\int_{0} ^{1} \! f(x,y)  \, \mathrm{d}y}  \, \mathrm{d}x   
\]
\minisec{Beweis}
Übung.
% subsection 75 (end)

\section*{Exkurs: Die Gammafunktion} % (fold)
\label{sec:exkurs_die_gammafunktion}


\subsection[Proposition/Definition: Die Gamma-Funktion]{Proposition/Definition} % (fold)
\label{sub:76}
Für $x>0$ ist $(t \mapsto t^{x-1} e^{-t})$ integrierbar über $(0,\infty)$. Wir definieren die \Index{Gammafunktion} 
\[
	\Gamma(x) := \int_{(0,\infty)} \! t^{x-1} e^{-t}  \, \mathrm{d}t
\]
\minisec{Beweis}
\begin{enumerate}[1.]
	\item $0 \le t^{x-1} e^{-t}\le t^{x-1}$. $t^{x-1}$ ist integrierbar über $(0,1]$, denn 
	\[
		\int_{(0,1]}\! t^s  \, \mathrm{d}t \stackrel{\text{\ref{sub:66}}}{=} \lim_{ k \to \infty} \int_{[\frac{1}{k+1},1 ]} \! t^s  \, \mathrm{d}t
		= \lim_{ k \to \infty} \int_{\frac{1}{k+1} } ^{1} \! t^s  \, \mathrm{d}t = \lim_{ k \to \infty}  \abs{\frac{1}{s+1} t^{s+1} }_{\frac{1}{k+1} }^{1}
		= \frac{1}{s+1} \tag*{für $s>-1$}  
	\]
	\item $0\le t^{x-1} e^{-t} \le \frac{1}{t^2} $, falls $t\ge t_0$ für ein geeignetes $t_0$, denn $t^2 \cdot t^{x-1} e^{-t} \xrightarrow{t \to \infty} 0$. Aber 
	$\frac{1}{t^2} $ ist integrierbar über $[t_0, \infty)$, also auch $t^{x-1} e^{-t}$.
	\item $t^{x-1} e^{-t}$ist stetig, also auch integrierbar auf $[1,t_0]$. \bewende
\end{enumerate}
% subsection 76 (end)

\subsection[Proposition: Rekursionsformel der Gammafunktion]{Proposition} % (fold)
\label{sub:77}
Für $x>0$ gilt $\Gamma(x+1)= x \cdot \Gamma(x)$. Weiter gilt $\Gamma(1)=1$, also $\Gamma(n+1)= n!$ , $n \in \mathds{N}^*$.
\minisec{Beweis}
\begin{enumerate}[1.)]
	\item \begin{align*}
		\Gamma(x+1) = \int_{(0,\infty)}\! t^{x+1-1} e^{-t}  \, \mathrm{d}t \stackrel{\text{\ref{sub:66}}}{=} \lim_{ k \to \infty} \int_{\frac{1}{k+1} } ^{k} \! 
		t^{x} e^{-t}  \, \mathrm{d}t &= \lim_{ k \to \infty} \enbrace{-t^x e^{-t} \Big|_{\frac{1}{k+1} }^k + \int_{\frac{1}{k+1} } ^{k} \! x \cdot t^{x-1} 
		\cdot e^{-t}  \, \mathrm{d}t} \\
		&= x \cdot \lim_{ k \to \infty} \int_{\frac{1}{k+1} } ^{k} \! t^{x-1} e^{-t}\\
		&= x \cdot \int_{(0,\infty)}\! t^{x-1} e^{-t}  \, \mathrm{d}t = x \cdot \Gamma(x)
	\end{align*}
	\item $\Gamma(1)= \int_{(0,\infty)}e^{-t}  \, \mathrm{d}t =1$.
	\item $\Gamma(n+1)=n!$ nach Induktion. \bewende
\end{enumerate}
% subsection 77 (end)

\subsection[Satz: Limesdarstellung der Gammafunktion]{Satz} % (fold)
\label{sub:78}
Für $x>0$ gilt 
\[
	\Gamma(x) = \lim_{ n \to \infty} \frac{n!\cdot  n^x}{x \cdot (x+1) \cdot \ldots \cdot (x+n)} 
\]
\minisec{Beweis}
Es ist $e^{-t}= \lim_{ n \to \infty} \enbrace{1- \frac{t}{n} }^n $ für $t>0$ (Analysis \RM{1}). Die Folge ist monoton wachsend in $n$. Setze 
\[
	f_n(t) := \begin{cases}
		t^{x-1} \cdot (1- \frac{t}{n} )^n, &\text{ falls }0 < t \le n\\
		0 , &\text{ sonst }
	\end{cases}
\]
Dann gilt für $t \in (0,\infty)$: $f_n(t) \to t^{x-1} e^{-t}$ monoton wachsend. Wir erhalten für $x>0$
\begin{align*}
	\Gamma(x) = \int_{(0,\infty)}\! t^{x-1} e^{-t}  \, \mathrm{d}t &\overset{n \to \infty}{\longleftarrow} \int_{(0,\infty)}\! f_n(t)  \, \mathrm{d}t
	= \lim_{ k \to \infty} \int_{\frac{1}{k+1} } ^{n} \! t^{x-1} \cdot \enbrace{1- \frac{t}{n}}^n  \, \mathrm{d}t\\
	&= \lim_{ k \to \infty} \enbrace{\frac{t^x}{x} \cdot \enbrace{1- \frac{t}{n} }^n \Bigg|_{\frac{1}{k+1} }^n + \int_{\frac{1}{k+1} } ^{n} \! \frac{t^x}{x} \cdot 
		\enbrace{1- \frac{t}{n} }^{n-1}  \, \mathrm{d}t}\\
	&= 0+ \lim_{ k \to \infty} \enbrace{\frac{t^{x+1}}{(x+1)x} \enbrace{1- \frac{t}{n} }^{n-1} \Bigg|_{\frac{1}{k+1} }^n - \int_{\frac{1}{k+1} } ^{n} \! 
		\frac{t^{x+1}}{(x+1)x} \cdot \frac{n-1}{n} \enbrace{1- \frac{t}{n} }^{n-2}  \, \mathrm{d}t}  \\
	&= 0 + \frac{n-1}{n} \cdot \frac{1}{(x+1)x} \cdot   \int_{0} ^{n} \! t^{x+1} \cdot \enbrace{1- \frac{t}{n} }^{n-2}  \, \mathrm{d}t = \ldots\\
	& = \frac{(n-1) \cdot (n-2) \cdot \ldots  \cdot 1}{n \cdot n \cdot \ldots \cdot n} \cdot \frac{1}{x \cdot (x+1) \cdot \ldots \cdot (x+n-1)} \cdot  \int_{0} ^{n} \! 
		t^{x+n+1} \enbrace{1- \frac{t}{n} }^{n-n} \\
	&= \frac{n! \cdot n^x}{x \cdot (x+1) \cdot \ldots \cdot (x+n-1)} \bewende
\end{align*}
% subsection 78 (end)

\subsection[Corollar: $\Gamma(1/2) = \sqrt{\pi}$]{Corollar} % (fold)
\label{sub:79}
$\Gamma(\frac{1}{2} ) = \sqrt{\pi }  $\marginnote{Das ist ein Knaller!}
\minisec{Beweis}
\begin{align*}
	\Gamma(\frac{1}{2} )^2 &= \lim_{ n \to \infty} \frac{n! \cdot n^{\frac{1}{2} }}{\frac{1}{2}(1+ \frac{1}{2} ) \cdot \ldots \cdot (n+ \frac{1}{2} ) } \cdot 
	\frac{n! \cdot n^{\frac{1}{2}} }{(1-\frac{1}{2} ) \cdot (2-\frac{1}{2} ) \cdot \ldots \cdot (n- \frac{1}{2} ) (n+\frac{1}{2} )}  \\
	&= \lim_{ n \to \infty} \frac{2 \cdot  (n!)^2 \cdot n}{(n+\frac{1}{2} ) \cdot (1^2- \frac{1}{4} ) (2^2 - \frac{1}{4} ) \cdot \ldots \cdot (n^2- \frac{1}{4} )} \\
	&= 2 \cdot \lim_{ n \to \infty}  \frac{n}{n+\frac{1}{2} } \cdot \lim_{ n \to \infty} \frac{(n!)^2}{(1^2 - \frac{1}{4} ) \cdot \ldots \cdot (n^2- \frac{1}{4} )} \\
	&= 2 \cdot \lim_{ n \to \infty} \underbrace{\prod_{k=1}^{n} \frac{k^2}{k^2- \frac{1}{4} } }_{=: c_n}
\end{align*}
Setze $A_m := \int_{0} ^{\frac{\pi }{2} } \! \sin^m t  \, \mathrm{d}t$. Dann gilt $A_0 = \frac{\pi }{2}, A_1 =1 $, $A_m = \frac{m-1}{m} A_{m-2}, m \ge 2 $ 
(partielle Integration). \\
Induktion $\leadsto$
\begin{align*}
	A_{2n} &= \frac{(2n-1) (2n-3) \cdot \ldots \cdot 3 \cdot 1}{2n \cdot (2n-2) \cdot \ldots \cdot 4 \cdot 2} \cdot \frac{\pi}{2} \\
	A_{2n+1} &= \frac{2n \cdot (2n-2) \cdot \ldots \cdot 4 \cdot 2}{(2n+1) \cdot (2n-1) \cdot \ldots \cdot 5 \cdot 3} \cdot 1 
\end{align*}
Es gilt $\sin^{2n+2}t \le \sin^{2n+1} t \le \sin^{2n} t $, $t \in [0, \frac{\pi}{2}]$, also
\[
	A_{2n+2} \le A_{2n+1} \le A_{2n}.
\]
Damit gilt dann $\frac{A_{2n+2}}{A_{2n}} = \frac{2n+1}{2n+2} \xrightarrow{n \to \infty} 1  $, also
\begin{align*}
	1\overset{\infty \leftarrow n}{\longleftarrow} \frac{A_{2n+2}}{A_{2n}} \le \underbrace{\frac{A_{2n+1}}{A_{2n}}}_{= \enbrace{\prod_{k=1}^n \frac{4k^2}{4k^2-1}} \cdot \frac{2}{\pi} = c_n \cdot \frac{2}{\pi} \to 1  } \le \frac{A_{2n}}{A_{2n}} = 1  
\end{align*}
Also $c_n \xrightarrow{n \to \infty} \frac{\pi }{2} $, woraus die Behauptung folgt. \bewende
% subsection 79 (end)

\subsection[Beispiel: Gaußintegral]{Beispiel} % (fold)
\label{sub:710}
Es gilt
\[
	\int_{\mathds{R}^n}\! e^{-\norm[2]{x}^2}  \, \mathrm{d}x  = \pi^{\frac{n}{2} } \enspace \text{ insbesondere } \enspace \int_{\mathds{R}}\! e^{-t^2}  \, \mathrm{d}t= 
	\pi^{\frac{1}{2} } \tag{\Index{Gaußintegral}}
\]
\minisec{Beweis}
Mit $n=t^2, \diffd{n}{t} = 2t, \mathrm{d}t = \frac{n^{-\frac{1}{2} }}{2}  \, \mathrm{d}n $
\begin{align*}
	\int_{\mathds{R}} e^{-t^2}  \, \mathrm{d}t \stackrel{\text{?}}{=} 2 \cdot \int_{(0,\infty)} e^{-t^2}  \, \mathrm{d}t = \int_{(0,\infty)}\! n^{-\frac{1}{2} } e^{-n}  \, \mathrm{d}n = \Gamma\enbrace{\frac{1}{2}} = \sqrt{\pi } 
\end{align*}
Weiter gilt nach Fubini (\ref{sub:73}) und Tonelli (\ref{sub:74})
\begin{align*}
	\int_{\mathds{R}^n}\! e^{- \norm[2]{x}^2 }  \, \mathrm{d}x  = \int_{\mathds{R}^n}\! e^{-(x_1^2 + x_2^2+ \ldots + x_n^2)}   \mathrm{d}(x_1, \ldots , x_n)
	&= \int_{\mathds{R}^{n-1}}\! \enbrace{\int_{\mathds{R}}\! e^{-(x_1^2 + x_2^2+ \ldots + x_n^2)}  \, \mathrm{d}x_1 }   \mathrm{d}(x_2, \ldots , x_n) \\ 
	&= \int_{R^{n-1}}\! \!\enbrace{ \int_{\mathds{R}}\! e^{-x_1^2}  \, \mathrm{d}x_1}  e^{-(x_2^2+ \ldots + x_n^2)}  \mathrm{d}(x_2, \ldots , x_n)   \\
	&= \ldots = (\pi^{\frac{1}{2} })^n = \pi^{\frac{n}{2} } \bewende
\end{align*}
% subsection 710 (end)

\subsection[Beispiel: Kugelvolumen im $\mathds{R}^n$]{Beispiel} % (fold)
\label{sub:711}
Für das Kugelvolumen im $\mathds{R}^n$ gilt
\[
	v  \Big(B_{\mathds{R}^n} (0,R)\Big) = \frac{\pi^{\frac{n}{2} } \cdot R^n}{\Gamma\enbrace{\frac{n}{2}+1} }  \quad {\color{light_gray} \xrightarrow{n \to \infty} 0}
\]
\minisec{Beweis}
Übung (vgl. Beispiel \ref{sub:54}) \bewende
% subsection 711 (end)

\subsection[Beispiel: Euler'sche Betafunktion]{Beispiel} % (fold)
\label{sub:712}
Wir definieren die \Index{Euler'sche Betafunktion} $B :(0,\infty) \times (0,\infty) \to \mathds{R}_+$ durch 
\[
	B(x,y) := \int_{(0,1)}\! t^{x-1} (1-t)^{y-1}  \, \mathrm{d}t \stackrel{\text{Subs. }t=1/(1+u)}{=} \int_{(0,\infty)} \frac{u^{x-1}}{(1+u)^{x+y}}  \, \mathrm{d}u 
\]
{\color{light_gray} (Warum existiert das Integral?)} \\
Behauptung: 
\[
	B(x,y) = \frac{\Gamma(x) \cdot \Gamma(y)}{\Gamma(x+y)} 
\]
\minisec{Beweis}
Mit Substitution $t= R \cdot s, \diffd{t}{s} = R$
\begin{align*}
	\Gamma(z) = \int_{0} ^{\infty} \! t^{z-1} e^{-t}  \, \mathrm{d}t = \int_{0} ^{\infty} \! (R \cdot s)^{z-1} e^{-R s} \cdot R  \, \mathrm{d}s
	= R^z \int_{0} ^{\infty} \! s^{z-1} e^{-Rs}  \, \mathrm{d}s \tag{$R>0$}
\end{align*}
Dann gilt für $z= x+y, R=1+u$
\begin{align*}
	\leadsto \frac{1}{\Gamma(x,y)} &\cdot \int_{0} ^{\infty} \! t^{x+y-1} e^{-(1+u)t}  \, \mathrm{d}t = \frac{1}{(1+u)^{x+y}} \\
	\leadsto  B(x,y) &= \frac{1}{\Gamma(x+y)} \cdot \int_{0} ^{\infty} \! u^{x-1} \enbrace{\int_{0} ^{\infty} \! t^{x+y-1} e^{-(1+u)t}  \, \mathrm{d}t}  \, \mathrm{d}u \\
	&= \frac{1}{\Gamma(x+y)} \int_{0} ^{\infty} \!  \enbrace{ \int_{0} ^{\infty} \! u^{x-1} t^{x+y-1} e^{-(1+u)t}   \, \mathrm{d}t }  \, \mathrm{d}u   \\
	&= \frac{1}{\Gamma(x+y)} \cdot \int_{0} ^{\infty} \!  \enbrace{ \int_{0} ^{\infty} \! u^{x-1} t^{x+y-1} e^{-(1+u)t}  \, \mathrm{d}u }  \, \mathrm{d}t \\
	&= \frac{1}{\Gamma(x+y)} \cdot \int_{0} ^{\infty} \! t^{x+y-1} e^{-t} \cdot  \enbrace{\int_{0} ^{\infty} \! u^{x-1} e^{-ut}  \, \mathrm{d}u}  \, \mathrm{d}t \\
	\intertext{Substitution $ut=s$}
	&= \frac{\Gamma(x)}{\Gamma(x+y)} \cdot \int_{0} ^{\infty} \! t^{y-1} e^{-t} = \frac{\Gamma(x) \cdot \Gamma(y)}{\Gamma(x+y)} \bewende 
\end{align*}
% subsection 712 (end)

\subsection{Beispiel (Dirichlet)} % (fold)
\label{sub:713}
Sei $\triangle := \set[(x,y) \in (0,\infty) \times (0,\infty)]{x+y<1} $ das \Index{Standardsimplex} im $\mathds{R}^2$ und $f : \triangle \to \mathds{R}_+$ gegeben durch
$f(x,y) := x^{p-1} \cdot y^{q-1}$, $p,q >0$. Es gilt
\begin{align*}
	\int_{\triangle}\! f(x,y)  \, \mathrm{d}(x,y) =  \int_{(0,1)}\! \enbrace{\int_{(0,1-y)}\! f(x,y)  \, \mathrm{d}x }   \mathrm{d}y 
	&= \int_{(0,1)}\! y^{q-1} \enbrace{\int_{(0,1-y)}\! x^{p-1}  \, \mathrm{d}x }  \, \mathrm{d}y \\ 
	&= \int_{(0,1)}\! y^{q-1} \frac{1}{p} \cdot (1-y)^p  \, \mathrm{d}y \\
	&= \frac{1}{p} \cdot B(q,p+1) \\
	&= \frac{1}{p} \cdot \frac{\Gamma(q) \cdot \Gamma(p+1)}{\Gamma(q+p+1)} \\
	&= \frac{\Gamma(q) \cdot \Gamma(p)}{\Gamma(q+p+1)}  
\end{align*}
% subsection 713 (end)
% section exkurs_die_gammafunktion (end)
% section 7 (end)
\newpage

\section{Der Transformationssatz} % (fold)
\label{sec:8}
Erinnerung: $T : [a,b] \to [c,d]$ bijektiv, stetig differenzierbar, $f : [c,d] \to \mathds{R}$ stetig. Dann gilt 
\[
	\int_{a} ^{b} \! f\big(T(x)\big) \diffd{T}{x}(x) \, \mathrm{d}x  = \int_{T(a)} ^{T(b)} \! f(t)  \, \mathrm{d}t,
\]
also $\int_{[a,b]}\! f\big( T(x)\big) \abs{ \diffd{T}{x}(x)}  \, \mathrm{d}x = \int_{[c,d] = T([a,b])}\! f(t)  \, \mathrm{d}t$

\subsection{Transformationssatz} % (fold)
\label{sub:81}
Seien $U,V \subset \mathds{R}^n$ offen, $T : U \to V$ ein $C^1$-\Index{Diffeomorphismus}\footnote{$T$ stetig diff'bar, $T ^{-1}$ stetig diff'bar}; sei 
$f : V \to (-\infty, \infty]$ eine Funktion. Dann ist $f$ über $V$ integrierbar, genau dann wenn 
\[
	(f \circ T) \cdot \abs{\det(\mathrm{D}T)}
\]
über $U$ integrierbar ist. In diesem Fall gilt
\[
	\int_U\! (f \circ T)(x) \cdot \abs{\det \big(\mathrm{D} T (x)\big)}  \, \mathrm{d}x  = \int_V\! f(y)  \, \mathrm{d}y 
\]
\minisec{Beweis (Skizze)}
\begin{description}
	\item[Strategie:] Sei $V$ ein Quader und $f$ eine Treppenfunktion. 
	\begin{figure}[h]
		\centering{
		\begin{tikzpicture}
			\draw[thick] (7,0) grid (12,5);
			\node[below] at (9.5,0) {$V$};
		\end{tikzpicture}
		\caption{Veranschaulichung der Beweisstrategie von \ref{sub:81}}}
	\end{figure} 
	$T ^{-1}$ ist differenzierbar, d.h. für $y \in V$ gilt 
	\[
		T ^{-1} (y+ \xi) = T ^{-1}(y)+ \mathrm{D}(T ^{-1})(y) \xi + \varphi(\xi)
	\]
	,wo $\mathrm{D}(T ^{-1})(y) : \mathds{R}^n \to \mathds{R}^n$ linear (und invertierbar) ist und $\varphi: \mathds{R}^n \to \mathds{R}^n$ eine Abbildung mit 
	$\frac{1}{\norm[2]{\xi} } \varphi(\xi) \xrightarrow{\xi \to 0} 0 $. Setze $\tilde P_k := P_k - x_k$, $\tilde Q_k := Q_k - y_k$, dann
	\begin{align*}
		\tilde P_k = P_k - x_k = T ^{-1}(Q_k) - T ^{-1}(y_k) &\approx T ^{-1}(y_k)+ \mathrm{D}(T ^{-1})(y_k)(\tilde Q_k) - T ^{-1}(y_k) \\
		&= D(T ^{-1})(T(x_k))(\tilde Q_k) = \big(\mathrm{D}T(x_k)\big) ^{-1} (\tilde Q_k)
	\end{align*}
	Es folgt
	\begin{align*}
		v(P_k) = v(\tilde P_k) \approx v\big( (\mathrm{D} T(x_k)) ^{-1} (\tilde Q_k)\big) &\stackrel{\text{\ref{sub:516}}}{=} \abs{\det \big(\mathrm{D}T (x_k)\big) ^{-1}} 
		\cdot v(\tilde Q_k) \\
		&= \abs{ \det \D T(x_k)} ^{-1} \cdot v(Q_k) 
	\end{align*}
	also $\abs{\det \D T(x_k)} \cdot v(P_k) \approx v(Q_k) $. Daraus folgt
	\begin{align*}
		\int_V\! f(y)  \, \mathrm{d}y = \sum_{k} f(y_k) \cdot v(Q_k) &\stackrel{!}{\approx} \sum_{k} f\big(T(x_k)\big) \cdot \abs{\det \D T(x_k)} \cdot v(P_k)   \\
		&\stackrel{!}{\approx} \int_U\! f \circ T(x) \cdot \abs{\det \D T(x)}  \, \mathrm{d}x.
	\end{align*}
\end{description}
% subsection 81 (end)

\subsection[Lemma: Urbild einer Nullmenge unter einen $C^1$-Diffeomorphismus]{Lemma} % (fold)
\label{sub:82}
Sei $T, U,V$ wie in \ref{sub:81}, $N \subset V$ eine Nullmenge. Dann ist $T ^{-1}(N)$ eine Nullmenge.
\minisec{Beweisidee}
Man darf $N \subset K \subseteq V$ für ein kompaktes $K$ annehmen, benutze dann, dass $T ^{-1}\big|_K$ Lipschitz ist, sowie Bemerkung \ref{sub:511} ($N$ lässt sich durch
abzählbar viele Quader mit kleinem Gesamtvolumen überdecken). \bewende
% subsection 82 (end)

\subsection[Lemma: Abschätzungen für das Volumen eines Quaders mit kompaktem Urbild]{Lemma} % (fold)
\label{sub:83}
Sei $P \subset U$ eine kompakte Teilmenge, sodass $Q := T(P) \subset V$ ein kompakter Quader ist. Dann gilt:
\[
	\min_{x \in P} \abs{\det \D T(x)} \cdot v(P) \le v(Q)\le \max_{x \in P} \abs{\det \D T(x)} \cdot v(P)
\]
\minisec{Beweis benutzt:}
Corollar \ref{sub:516}, Lemma \ref{sub:82}, Kompaktheit von $P$ und Stetigkeit von $( x \mapsto \abs{\det \D T(x)} )$ \bewende
% subsection 83 (end)

\subsection[Proposition: Transformationssatz für Treppenfunktionen]{Proposition} % (fold)
\label{sub:84}
Der Transformationssatz \ref{sub:81} gilt für $f_V \in \mathcal{T}(\mathds{R}^n) $ mit $\supp f := \overline{\set[x \in V]{f(x)\not= 0}} \subset V $.\index{kompakter Träger}
\minisec{Beweis (Skizze)}
Wegen Linearität  und Lemma \ref{sub:82} genügt es den Satz für $f = \chi_Q$ zu beweisen, wo $Q \subset V$ ein kompakter Quader ist.
Integrierbarkeit von $(\chi_Q \circ T) \cdot \abs{\det \D T} $ folgt aus Kompaktheit von $T ^{-1}(Q)$ und Stetigkeit von ($x \mapsto \abs{\det \D T(x)} $).
Zu zeigen: $\int_{T ^{-1}(Q)}\! \abs{\det \D T(x)}  \, \mathrm{d}x  = \int_Q\! 1  \, \mathrm{d}y $

Sei $\varepsilon>0$. $\abs{\det \D T ^{-1}} ^{-1} $ ist gleichmäßig stetig auf $Q$, daher existieren $Q_1, \ldots Q_r \subset Q$ kompakte Quader, sodass 
\begin{enumerate}[(i)]
	\item $Q = \bigcup_{i=1}^r Q_i$
	\item $Q_i \cap Q_{i'}$ ist Nullmenge, falls $i \not= i'$
	\item $\max_{y \in Q_i} \abs{\det \D (T ^{-1})(y)} ^{-1} - \min_{y \in Q_i} \abs{\det \D (T ^{-1})(y)} ^{-1} \le \varepsilon$, \quad $i=1,\ldots ,n$
\end{enumerate}
Wegen\marginnote{Satz von der Umkehrabbildung} $\D( T ^{-1})( T(x)) = \D T(x) ^{-1}$ und $\det (A ^{-1}) = (\det A) ^{-1}$ (für $A \in M_n(\mathds{R})$), gilt 
\[
	\abs{\det \D (T ^{-1}) \big( T(x)\big) } ^{-1} =  \abs{\det \D T(x)}
\]
Für $P_i := T ^{-1}(Q_i)$ erhalten wir: 
\[
	\max_{x \in P_i} \abs{\det ( \D T(x))} - \min_{x \in P_i} \abs{\det ( \D T(x))} \stackrel{\text{(iii)}}{\le} \varepsilon \tag{$\star$}
\]
Dann gilt:
\[
	\abs{\int_{P_i}\! \abs{\det \D T(x)}  \, \mathrm{d}x  - v(Q_i)} \stackrel{\text{($\star$), \ref{sub:83}}}{\le} \varepsilon \cdot v(P_i)
\]
Es folgt
\begin{align*}
	\int_{T ^{-1}(Q)}\! \abs{\det \D T(x)}  \, \mathrm{d}x \stackrel{\text{(ii), \ref{sub:82}}}{=} \sum_{i=1}^{r} \int_{P_i}\! \abs{\det \D T(x)}  \, \mathrm{d}x 
	\underset{\substack{\sum_{i=1}^{r} \varepsilon \cdot v(P_i) \\= \varepsilon \cdot v(T ^{-1}(Q)) }}{=} \sum_{i=1}^{r} v(Q_i) = v(Q) = \int_{Q} \! 1  \, \mathrm{d}y 
\end{align*}  
$\varepsilon>0$ beliebig $\Rightarrow $ Behauptung. \bewende
% subsection 84 (end)

\subsection*{Beweis von \ref{sub:81} (Skizze)} % (fold)
%\label{sub:85}
Sei $f$ integrierbar über $V$. Konstruiere Treppenfunktionen $(\varphi_k)_{k \in \mathds{N}} \subset \mathcal{T}(\mathds{R}^n)$ mit
\begin{enumerate}[(i)]
	\item $\supp \varphi_k  \subset V$
	\item $\norm[1]{f_V- \varphi_k} \xrightarrow{k \to \infty} 0$
	\item Für eine Nullmenge $N \subset V$ gilt: $\varphi_k \xrightarrow{k \to \infty} f_V$ punktweise auf $V \setminus N$.\hfill(mit \ref{sub:63})
\end{enumerate}
Nach Proposition \ref{sub:84} gilt:
\begin{align*}
	\norm[1]{(\varphi_k \circ T) \cdot \abs{\det(\D T)} - (\varphi_l \circ T) \cdot \abs{\det(\D T)}\rule{0cm}{0.45cm} } &= \int_U \! \abs{\varphi_k \circ T- \varphi_l \circ T} \cdot 
	\abs{\det(\D T)}  \, \mathrm{d}x \\
	&\stackrel{\text{\ref{sub:84}}}{=} \int_V \abs{\varphi_k - \varphi_l}  \, \mathrm{d}y  \\
	&= \norm[1]{\varphi_k - \varphi_l} 
\end{align*}
$\Rightarrow \enbrace{ (\varphi_k \circ T) \cdot \abs{\det (\D T)} }_{k \in \mathds{N}} $ ist Cauchy-Folge bezüglich $\norm[1]{.} $. Weiter gilt 
\begin{align*}
	(\varphi_k \circ T) \cdot \abs{ \det (\D T)} \xrightarrow{k \to \infty} (f_V \circ T) \circ \abs{\det(\D T)} 
\end{align*}
punktweise fast überall, d.h. auf $U \setminus T^{-1}(N)$ ($T ^{-1}(N)$ ist Nullmenge nach \ref{sub:82}). Nach dem Satz von Riesz-Fischer (\ref{sub:63}) ist dann
$(f_V \circ T) \cdot \abs{\det (\D T)} $ integrierbar über $U$ und 
\[
	\int_U \!(f \circ T) \cdot \abs{ \det (\D T)}  \, \mathrm{d}x ) = \lim_{ k \to \infty} \int_U \! (\varphi_K \circ T) \cdot \abs{\det (\D T)}  \, \mathrm{d}x   
	\stackrel{\text{\ref{sub:84}}}{=} \lim_{ k \to \infty} \int_V \! \varphi_k  \, \mathrm{d}y \stackrel{\text{\ref{sub:63}}}{=} \int_V \! f  \, \mathrm{d}y.
\]
Rückrichtung $T \leadsto T ^{-1}$ \bewende
% subsection 85 (end)

\subsection[Beispiel: Affine Transformationen und Polarkoordinaten]{Beispiel} % (fold)
\label{sub:85}
\begin{enumerate}[(i)]
	\item Sei $A \in GL(n, \mathds{R})$ eine invertierbare $n \times n$-Matrix, $b \in \mathds{R}^n$. Sei $T : \mathds{R}^n \to \mathds{R}^n$ die 
	\Index{affine Transformation} $x \mapsto A x+ b$. Dann gilt $\D T(x) =A, x \in \mathds{R}^n$ (vgl. Ana \RM{2}) und $f$ ist über $K \subset \mathds{R}^n$ integrierbar
	genau dann, wenn $f \circ T$ über $T ^{-1} (K)$ integrierbar ist und es gilt
	\[
		\abs{\det A} \cdot  \int_{T ^{-1}(K)}\! f(Ax+b)  \, \mathrm{d}x = \int_K f(y)  \, \mathrm{d}y.
	\]
	Für $A \in O(n)$ (orthogonale Matrix) gilt 
	\[
		v\big( T(K)\big) = v(K)
	\]
	(\Index{Bewegungsinvarianz} des Lebesgue-Maßes).
	\item Polarkoordinaten in $\mathds{R}^2$ \\
	$P_2 : (0,\infty) \times (- \pi , \pi) \to \mathds{R}^2 \setminus (- \infty,0] \times \{0\}, \binom{r}{\varphi} \mapsto \binom{r \cdot \cos \varphi}{r \cdot \sin 
	\varphi}$ ist bijektiv. Es ist 
	\[
		\D P_2 \binom{r}{\varphi} = \begin{pmatrix}
			\cos \varphi & -r \sin \varphi \\
			\sin \varphi & r \cos \varphi
		\end{pmatrix} \tag{vgl. Ana \RM{2} Bsp. 8.6 (ii)}
	\]
	und $\det \D P_2 \binom{r}{\varphi} = r >0$. Daraus folgt $\D P_2 \binom{r}{\varphi} $ ist invertierbar $\Rightarrow P_2$ ist $C^1$-Diffeomorphismus. Für
	$f : \mathds{R}^2 \to \mathds{R}^2$ integrierbar gilt nun
	\begin{align*}
		\int_{\mathds{R}^2}\! f(x,y)  \, \mathrm{d}(x,y) = \int_{\mathds{R}^2 \setminus (- \infty, 0] \times \{0\}} f(x,y)  \, \mathrm{d}(x,y)
		&= \int_{(0, \infty) \times (-\pi , \pi )}\! f\big( P_2(r,\varphi)\big) \cdot r  \enspace \mathrm{d}(r,\varphi) \\ 
		&= \int_{(0,\infty)}\! r \enbrace{ \int_{(-\pi ,\pi )}\! f(P_2(r,\varphi))  \, \mathrm{d}\varphi}  \mathrm{d}r 
	\end{align*}
	\item Kugelkoordinaten im $\mathds{R}^3$ \\
	$P_3 : (0,\infty) \times (-\pi ,\pi )\times (- \frac{\pi}{2}, \frac{\pi}{2}) \to \mathds{R}^3 \setminus (-\infty, 0] \times \{0\} \times \mathds{R}$
	\[
		\begin{pmatrix}
			r \\
			\varphi \\
			\theta
		\end{pmatrix} \mapsto
		\begin{pmatrix}
			r \cdot \cos \varphi \cdot \cos \theta \\
			r \cdot \sin \varphi \cdot \cos \theta \\
			r \cdot \sin \theta
		\end{pmatrix}
	\]
	ist bijektiv. Es gilt
	\[
		\det \enbrace{\D P_3 \begin{pmatrix}
			r \\
			\varphi \\
			\theta
		\end{pmatrix}} = r^2 \cdot \cos \theta > 0
	\]
	für $r>0, \theta \in (- \frac{\pi}{2}, \frac{\pi}{2})$. Für $f$ integrierbar gilt
	\begin{align*}
		\int_{\mathds{R}^3} f(x,y,z)  \, \mathrm{d}(x,y,z) = \int_{(0,\infty)} \int_{(-\pi , \pi )} \int_{(- \frac{\pi}{2}, \frac{\pi}{2})}
		f\big( P_3(r, \varphi, \theta) \big) \cdot r^2  \cdot \cos \theta  \enspace \mathrm{d}\theta  \, \mathrm{d}\varphi  \, \mathrm{d}r
	\end{align*}
	\item 
	\[
		\int_{\mathds{R}^3} \chi_{B(0,R)}  \, \mathrm{d}(x,y,z) = \ldots = \frac{4}{3} \pi  R^3
	\]
\end{enumerate}
% subsection 85 (end)

\subsection{Corollar (Integration rotationssymetrischer Funktionen)} % (fold)
\label{sub:86}
Sei $g : \mathds{R}^n \to \mathds{R}$ eine Funktion, welche sich in der Form $g(x) = f(\norm[2]{x})$ schreiben lässt für eine Funktion $f : \mathds{R}_+ \to \mathds{R}$.
Dann ist $g$ über $\mathds{R}^n$ integrierbar genau dann, wenn $f(r) \cdot r^{n-1}$ über $(0, \infty)$ integrierbar ist. In diesem Fall gilt 
\[
	\int_{\mathds{R}^n} \! g(x)  \, \mathrm{d}x  = n \cdot v\big(B_{\mathds{R}^n}(0,1)\big) \cdot \int_{(0,\infty)}\! f(r) r^{n-1}  \, \mathrm{d}r
\]
\minisec{Beweis}
Für $n=2,3$ Übung. Für $n>3$ ohne Beweis. \index{rotationssymmetrische Funktion}
% subsection 86 (end)
% section 8 (end)
\newpage

\section{$\sigma$-Algebren und messbare Räume} % (fold)
\label{sec:9}

\subsection[Erinnerung: Topologien, Stetigkeit in topologischen Räumen]{Erinnerung} % (fold)
\label{sub:91}
Sei $X$ eine Menge. $\mathcal{T} \subset \mathcal{P}(X)$ heißt \Index{Topologie} auf $X$, falls 
\begin{enumerate}[(i)]
	\item $X, \emptyset \in \mathcal{T} $,
	\item $U_1, \ldots , U_n \in \mathcal{T} \Longrightarrow  \bigcap_{i=1}^n U_i \in \mathcal{T} $,
	\item $U_i \in \mathcal{T}, i \in I \Longrightarrow \bigcup_{i \in I} U_i \in \mathcal{T} $.
\end{enumerate} 
$(X,\mathcal{T})$ heißt \Index{topologischer Raum}, $U \in \mathcal{T} $ heißt \Index{offen}. Falls $(X_j, \mathcal{T}_j)$, $j=1,2$ topologischer Räume sind und
$f : X_1 \to X_2$ eine Abbildung, dann heißt $f$ \bet{stetig}\index{topologischer Raum!Stetigkeit}, falls gilt
\[
	U \in \mathcal{T}_2 \Longrightarrow  f ^{-1} (U) \in \mathcal{T}_1.
\] 
% subsection 91 (end)

\subsection[Definition: $\sigma$-Algebra, messbarer Raum]{Definition} % (fold)
\label{sub:92}
Sei $X$ eine Menge. $\Sigma  \subset \mathcal{P}(X)$ heißt \bet{$\sigma$-Algebra}\index{Sigma-Algebra@$\sigma$-Algebra} auf $X$, falls 
\begin{enumerate}[(i)]
	\item $X \in \Sigma$,
	\item $A \in \Sigma \Rightarrow X \setminus A \in \Sigma$,
	\item $A_n \in \Sigma, n \in \mathds{N} \Rightarrow \bigcup_{n \in N} A_n \in \Sigma$.
\end{enumerate}
$(X, \Sigma)$ heißt \Index{messbarer Raum}; die Elemente von $\Sigma$ heißen \Index{messbare Mengen}.  $(X_j, \Sigma_j)$ für $ j=1,2$ seien messbare Räume,
$f : X_1 \to X_2$ eine Abbildung. $f$ heißt \bet{messbar}\index{messbare Abbildung}, falls gilt
\[
	A \in \Sigma_2 \Longrightarrow f ^{-1}(A) \in \Sigma_1.
\]
% subsection 92 (end)

\subsection[Bemerkung: Schnitte und Komposition messbaren Abbildungen in $\sigma$-Algebren]{Bemerkung} % (fold)
\label{sub:93}
\begin{enumerate}[(i)]
	\item $\sigma$-Algebren sind abgeschlossen bezüglich abzählbarer Vereinigungen und Durchschnitten. \marginnote{de Morgan}
	\item $(X_j, \Sigma_j)$ für $j=1,2,3$, $f : X_1 \to X_2$, $g : X_2 \to X_3$ messbar. $\Rightarrow $ 
	$g \circ f : X_1 \to X_3$ messbar.
\end{enumerate}
% subsection 93 (end)

\subsection[Proposition: Kleinste $\sigma$-Algebra]{Proposition} % (fold)
\label{sub:94}
Sei $X$ eine Menge und $\mathcal{F} \subset \mathcal{P}(X)$ eine Menge von Teilmengen. Dann existiert eine kleinste $\sigma$-Algebra, die $\mathcal{F}$ enthält; diese nennen
wir die von $\mathcal{F}$ \bet{erzeugte $\sigma$-Algebra}\index{Sigma-Algebra@$\sigma$-Algebra!erzeugte} $\Sigma(\mathcal{F})$. 
\minisec{Beweis}
Setze 
\[
	\Omega := \set[\Gamma \subset \mathcal{P}(X)]{ \Gamma \text{ ist } \sigma \text{-Algebra und } \mathcal{F} \subset\rule{0cm}{0.35cm} \Gamma} 
\]
Dann ist $\Omega \not= \emptyset$ (denn $\mathcal{P} (X) \in \Omega$). Setze $\Sigma (\mathcal{F}) := \bigcap_{\Gamma \in \Omega} \Gamma \subset \mathcal{P} (X)$. Dann ist
$\mathcal{F} \subset \Sigma (\mathcal{F})$, denn $\mathcal{F} \subset \Gamma$ für $\Gamma \in \Omega$. Ist $\Gamma \subset \mathcal{P}(X)$ eine $\sigma$-Algebra, welche
$\mathcal{F}$ enthält, so ist $\Gamma  \in \Omega$ und $\Sigma(\mathcal{F}) \subset \Gamma$. Es bleibt zu zeigen: $\Sigma (\mathcal{F})$ ist $\sigma$-Algebra:
\begin{enumerate}[(i)]
	\item $X \in \Sigma(\mathcal{F})$, denn $X \in \Gamma$ für jedes $\Gamma \in \Omega$.
	\item $A \in \Sigma(\mathcal{F}) \Rightarrow A \in \Gamma$ für jedes $\Gamma \in \Omega$. Es folgt $X \setminus A \in \Gamma $für jedes $\Gamma  \in \Omega$.
	Also ist $X \setminus A \in \Sigma(\mathcal{F})$.
	\item ebenso. \bewende
\end{enumerate}
% subsection 94 (end)

\subsection[Bemerkung: \ref{sub:94} lässt sich analog auch auf Topologien übertragen]{Bemerkung} % (fold)
\label{sub:95}
Die entsprechende Aussage für topologische Räume beweist man analog. 
% subsection 95 (end)

\subsection[Definition: $\sigma$-Algebra der Borelmengen, Borel-messbar]{Definition} % (fold)
\label{sub:96}
\begin{enumerate}[(i)]
	\item Sei $(X,\mathcal{T})$ ein topologischer Raum. Wir nennen $\Sigma(\mathcal{T} )$ die \bet{$\sigma$-Algebra der Borelmengen}
	\index{Sigma-Algebra@$\sigma$-Algebra! der Borelmengen} auf $X$ und schreiben $\mathcal{B}_X$ für $\Sigma(\mathcal{T} )$.
	\item Seien $(X,\mathcal{T}_X)$, $(Y,\mathcal{T}_Y)$ topologische Räume, dann heißt $f : X \to Y$ \Index{Borel-messbar}, falls $f$ messbar ist bezüglich 
	$\mathcal{B}_X$ und $\mathcal{B}_Y$.
\end{enumerate}
% subsection 96 (end)

\subsection[Bemerkung, welche Mengen Borel sind]{Bemerkung} % (fold)
\label{sub:97}
Offene und abgeschlossene Mengen sind Borel und ebenso $F_\sigma$'s (abzählbare Vereinigungen von abgeschlossenen Mengen), und $G_\delta$'s (abzählbare Durchschnitte von 
offenen Mengen).
% subsection 97 (end)

\subsection[Proposition über Abbildungen zwischen topologischen und messbaren Räumen]{Proposition} % (fold)
\label{sub:98}
Sei $(X, \Sigma)$ ein messbarer Raum und $(Y, \mathcal{T} )$ ein topologischer Raum. Sei $f : X \to Y$ eine Abbildung. Dann gilt
\begin{enumerate}[(i)]
	\item $\Omega := \set[E \subset Y]{f ^{-1}(E) \in \Sigma} $ ist eine $\sigma$-Algebra auf $Y$.
	\item Falls $f ^{-1}(E) \in \Sigma$ für jedes $E \in \mathcal{T} $, so ist $f$ messbar (bezüglich. $\Sigma$ und $\mathcal{B}_Y$). Insbesondere: Falls 
	$(X, \mathcal{T}_X)$ ein topologischer Raum ist und $\Sigma = \mathcal{B}_X$ und falls $f$ stetig ist, so ist $f$ Borel.
	\item Falls $Y=\mathds{R}$ und $f ^{-1}\big( (\alpha, \infty)\big) \in \Sigma$ für jedes $\alpha \in \mathds{R}$, so ist $f$ messbar.
	\item Aussage (iii) gilt analog für $Y= (-\infty, \infty]$ und $(\alpha, \infty]$, wobei $\mathcal{T} $, die von $\mathcal{T}_\mathds{R}$ und $(\alpha, \infty]$, 
	$\alpha \in \mathds{R}$ erzeugte Topologie ist.
\end{enumerate}
\minisec{Beweis}
\begin{enumerate}[(i)]
	\item $f ^{-1}(Y) =X$, $f ^{-1}(Y \setminus E) = X \setminus f ^{-1}(E)$, $f ^{-1}( E_1 \cup E_2 \cup \ldots ) = f ^{-1}(E_1) \cup f ^{-1}(E_2) \cup \ldots $
	\item Setze $\Omega := \set[E \subset Y]{f ^{-1}(E) \in \Sigma}$ wie in (i). Dann gilt $\mathcal{T}  \subset \Omega$, also $\mathcal{B}_Y \subset \Omega$ nach (i).
	$\Rightarrow f ^{-1}(E) \in \Sigma$ für jedes $E \in \mathcal{B}_Y$ $\Rightarrow f $ ist messbar.
	\item Wieder $\Omega := \set[E \subset Y]{f ^{-1}(E) \in \Sigma} $. Für $\alpha \in \mathds{R}$ wähle $\alpha_0 < \alpha_1 < \ldots < \alpha$ mit $\alpha_n \xrightarrow{n \to \infty} \alpha$. Dann ist $(\alpha_k, \infty) \in \Omega$ für jedes $k$. Also auch $(-\infty, \alpha_n]$ (nach (i)) und 
	$\bigcup_{n \in \mathds{N}} (-\infty, \alpha_n] = (- \infty, \alpha) \in \Omega$. $\Rightarrow (\alpha, \beta) = (- \infty, \beta) \cap (\alpha, \infty) \in \Omega$.
	Also für jedes $U \in \mathcal{T}_\mathds{R}$ ist $U \in \Omega$, denn $U= \bigcup_{n \in \mathds{N}} (\gamma_n, \delta_n)$ (warum?).
	
	Daraus folgt $\mathcal{B}_Y \subset \Omega$, also ist $f$ messbar.
	\item Übung. \bewende
\end{enumerate}
% subsection 98 (end)

\subsection[Proposition: $\sup$ und $\limsup$ von messbaren Funktionen sind messbar]{Proposition} % (fold)
\label{sub:99}
Sei $(X, \Sigma)$ ein messbarer Raum und $f_n : X \to (-\infty, \infty]$ eine Folge messbarer Funktionen. Dann sind $g := \sup_{n \in \mathds{N}} f_n$ und
$h := \limsup _{n \to \infty} f_n$ messbar\footnote{$h(x) = \lim_k \sup_{n \ge k} f_n(x)$}.
\minisec{Beweis}
\[
	g ^{-1} \big( (\alpha, \infty]\big) = \underbrace{\bigcup_{k \in \mathds{N}} \underbrace{f_n ^{-1} \big( (\alpha, \infty] \big)}_{\in\Sigma}}_{\in \Sigma}
	\stackrel{\text{\ref{sub:98}}}{\Longrightarrow} g \text{ messbar}
\]
Ebenso für $\inf$ an Stelle von $\sup$. Weiter gilt 
\[
	h (x)= \lim_{k \to \infty} \sup_{n \ge k}f_n(x) = \underbrace{\inf_{k \ge 0} \underbrace{\sup_{n \ge k} f_n(x)}_{\text{messbar}}}_{\text{messbar}} \bewende
\]
% subsection 99 (end)

\subsection[Corollar: Punktweise Limiten und $\max$, $\min$ messbarer Funktionen sind messbar]{Corollar} % (fold)
\label{sub:910}
Seien $(X, \Sigma)$ und $f_n : X \to (-\infty, \infty]$ messbar für $n \in \mathds{N}$. 
\begin{enumerate}[(i)]
	\item Falls $f_n \xrightarrow{n \to \infty} f$ punktweise, so ist $f$ messbar.
	\item Falls $d,e : X\to (-\infty, \infty]$ messbar sind, so auch $\max \{d,e\}, \min \{d,e\} $, also auch $d_+$ und $ d_-$.
\end{enumerate}
\minisec{Beweis}
$f = \limsup_{n \to \infty} f_n$ ; $\max \{d,e \} = \sup \set{d,e,e,e, \ldots } $ benutze \ref{sub:99}.  Mit $d_+ = \max \{d,0\}, d_- = \max \{-d,0\}$ folgt der zweite 
Teil. \bewende
% subsection 910 (end)

\subsection[Definition: Einfache Funktion]{Definition} % (fold)
\label{sub:911}
Sei $(X, \Sigma)$ ein messbarer Raum. $f : X \to (-\infty, \infty)$ heißt \bet{einfach}\index{einfache Funktion}, falls $f(X)$ endlich ist. 
% subsection 911 (end)

\subsection[Bemerkung: Summenschreibweise für einfache Funktionen]{Bemerkung} % (fold)
\label{sub:912}
Sei $f$ einfach, dann ist $f = \sum_{i=1}^{n} \alpha_i \cdot \chi_{A_i}$ für $\set{\alpha_1, \ldots , \alpha_n} =f(X) $ und $A_i = f ^{-1} \big( \{\alpha\} \big)$.
$f$ ist messbar genau dann, wenn die $A_i$ messbar sind. 
% subsection 912 (end)

\subsection[Prop.: Approximation nichtnegativer, messbarer Funktionen durch einfache Funktionen]{Proposition} % (fold)
\label{sub:913}
Sei $(X, \Sigma)$ ein messbarer Raum, $f : X \to [0, \infty]$ messbar. Dann existieren einfache, messbare Funktionen $(f_n)_{n \in \mathds{N}}$ mit 
$0 \le f_1 \le f_2 \le \ldots $ und $f_n \xrightarrow{n \to \infty} f$ punktweise.
\minisec{Beweis}
Definiere $\varphi_n : [0,\infty] \to [0,\infty)$ wie folgt: 
\[
	\begin{tikzpicture}[xscale=1,yscale=0.7]
		%\draw[help lines] (0,0) grid (6,4);
		\draw[<->, thick] (0,4) node[left]{$n$} -- (0,0) -- (4,0) node[below]{$n$} -- (6,0);
	  	\foreach \x in {0,0.25,0.5,...,3.75,4}%
			   \draw[line width=0.5mm,Firebrick2](\x,\x)--(\x+0.25,\x);
		\draw[line width=0.5mm,Firebrick2] (4,4) --(6,4);
		\draw[decorate, decoration={brace,amplitude=5pt, mirror, raise=2pt},yshift=0pt, thick] (4,0) -- (4,4) %
		node[midway, right, xshift=0.2cm]{$n \cdot 2^n $ Stufen};
	\end{tikzpicture}
\]
Dann ist $\varphi_n$ messbar, also auch $\varphi_n \circ f =: f_n$. Dann gilt $f_n \nearrow f$ punktweise.\bewende
% subsection 913 (end)

\subsection[Definition: Maß und Maßraum]{Definition} % (fold)
\label{sub:914}
Sei $(X,\Sigma)$ ein messbarer Raum. Ein \Index{Maß} auf $(X, \Sigma)$  ist eine Abbildungen $\mu : \Sigma \to [0,\infty]$ mit $\mu(\emptyset)=0$, welche \Index{abzählbar additiv}
ist, d.h. Falls $A_i, \, i \in \mathds{N}$ paarweise disjunkte messbare Mengen sind, so gilt 
\[
	\mu \enbrace{\bigcup_{i \in \mathds{N}} A_i} = \sum_{i\in \mathds{N}} \mu(A_i).  
\]
$(X, \Sigma, \mu)$ heißt \Index{Maßraum}; \Index{Wahrscheinlichkeitsraum} falls $\mu(X)=1$. 
% subsection 914 (end)

\subsection[Proposition: Eigenschaften von Maßen]{Proposition} % (fold)
\label{sub:915}
Sei $(X,\Sigma, \mu)$ ein Maßraum. Dann gilt
\begin{enumerate}[(i)]
	\item $A,B \in \Sigma, A \subset B$ $\Longrightarrow \mu (A) \le \mu(B)$
	\item $A_i \in \Sigma, i \in \mathds{N}$, $A_0 \subset A_1 \subset A_2 \subset \ldots $, dann gilt $\mu (A_i) \xrightarrow{i \to \infty} \mu \enbrace{\bigcup_{i \in \mathds{N}} A_i} $.
	\item $A_i \in \Sigma, i\in \mathds{N}$, $A_0 \supset A_1 \supset A_2 \supset \ldots $ und $\mu(A_0) < \infty$, dann gilt $\mu (A_i) \xrightarrow{i \to \infty} \mu \enbrace{\bigcap_{i \in \mathds{N}} A_i} $.
\end{enumerate}
\minisec{Beweis}
\begin{enumerate}[(i)]
	\item \[
		\mu (B) =\mu \enbrace{\bigcup_{i \in \mathds{N}} A_i}  \stackrel{\text{\ref{sub:914}}}{=} \sum_{i \in \mathds{N}} \mu(A_i)  \ge \mu (A)
	\]
	mit $A_0 := A, A_1 := B \setminus A, A_2 := \emptyset, A_3 := \emptyset, \ldots $ paarweise disjunkt und messbar.
	\item Setze $B_0 := A_0, B_1 := A_1 \setminus A_0, B_2 := A_2 \setminus A_1, \ldots , B_{i+1} := A_{i+1} \setminus A_i$, dann sind die $B_i$ paarweise disjunkt und 
	messbar;
	\[
		\mu \enbrace{\bigcup_{i \in \mathds{N}} A_i} = \mu \enbrace{\bigcup_{i \in \mathds{N}} B_i} = \sum_{i \in \mathds{N}} \mu(B_i) = \lim_{ n \to \infty} 
		\sum_{i=1}^{n} \mu(B_i) = \lim_{ n \to \infty} \mu \enbrace{\bigcup_{i=1}^n B_i} = \lim_{ n \to \infty} \mu(A_n) 
	\]
	\item Übung. \bewende
\end{enumerate}
% subsection 915 (end)

\subsection[Beispiele für Maßräume]{Beispiele} % (fold)
\label{sub:916}
\begin{enumerate}[(i)]
	\item $(X, \mathcal{P}(X),  \mu)$ mit $\mu(E) = \begin{cases}
			0, &\text{ falls }E = \emptyset\\
			\infty, &\text{ falls } E \not= \emptyset
		\end{cases}$
	\item $(X, \mathcal{P}(X),  \mu)$ mit $\mu(E) = \begin{cases}
			\abs{E} , &\text{ falls }E \text{ endlich}\\
			\infty, &\text{ sonst } 
		\end{cases} $ \hfill (\Index{Zählmaß})
	\item $(X, \mathcal{P}(X), \delta_{x_0})$ mit $\delta_{x_0}(E) = \begin{cases}
		1, &\text{ falls }x_0 \in E\\
		0 , &\text{ sonst}
	\end{cases}$ ist W'raum für jedes $X$ und $x_0 \in X$.
	\item Sei $\mu$ das Zählmaß auf $(\mathds{N}, \mathcal{P}(\mathds{N}))$. Sei $A_i := \set{i, i+1, \ldots } $, $i \in \mathds{N}$. Dann gilt 
	$A_1 \supset A_1 \supset A_2 \supset \ldots $ und $\mu (A_i) = \infty$, aber $\mu \enbrace{\bigcap_{i \in \mathds{N}} A_i} = 0$. In Proposition \ref{sub:915} (iii) ist
	$\mu(A_0)< \infty$ also wesentlich!
	\item \begin{align*}
		\hat \Lambda &= \set[E \subset \mathds{R}^n]{\chi_E \text{ ist lokal integrierbar}} \\
		&= \set[E \subset \mathds{R}^n]{E \cap K \text{ messbar im Sinne von Abschnitt \ref{sec:6} für jedes kompakte }K  \subset \mathds{R}^n} 
	\end{align*}
	ist $\sigma$-Algebra auf $\mathds{R}^n$:
	\begin{enumerate}[a)]
		\item $\emptyset, X \in \hat \Lambda$ denn für jedes kompakte $K$ ist $\chi_X$ über $K$ integrierbar.
		\item Sei $E \subset \mathds{R}^n, E \in \hat \Lambda$ und $K \subset \mathds{R}^n$ kompakt: 
		$\chi_{(\mathds{R}^n \setminus E) \cap K} = \chi_{K \setminus (K \cap E)} = \chi_{K} - \chi_{K \cap E}$ ist integrierbar. Also 
		$\mathds{R}^n \setminus E \in \hat \Lambda$.
		\item abzählbare Vereinigungen ähnlich (mit Beppo-Levi).
	\end{enumerate}
	$\hat \Lambda$ enthält die Borelmengen auf $\mathds{R}^n$ (vgl. \ref{sub:68}). $\lambda : \hat \Lambda \to [0,\infty]  $ gegeben durch 
	\[
		\lambda (E) := \begin{cases}
			v(E), &\text{ falls } \chi_E \text{ integrierbar}\\
			\infty, &\text{ sonst }
		\end{cases}
	\]
	ist ein Maß auf $\mathds{R}^n$, das \Index{Lebesgue-Maß}  (vgl. Bemerkung \ref{sub:68}). 
	\item $f \in \mathcal{L}^1(\mathds{R}^n) \Rightarrow $ $f$ ist messbar im Sinne von Abschnitt \ref{sec:9}
	(Benutze \ref{sub:98} und Corollar \ref{sub:64}, $f$ ist fast überall punktweiser Limes von Treppenfunktionen.)
\end{enumerate}
% subsection 916 (end)

\subsection[{Definition: Integrierbarkeit, $\norm[1]{.} $ in Maßräumen}]{Definition} % (fold)
\label{sub:917}
Sei $(X,\Sigma, \mu)$ ein Maßraum. 
\begin{enumerate}[(i)]
	\item Für eine einfache, messbare Funktion $s : X \to [0,\infty)$, $s = \sum_{i=1}^{k} \alpha_i \cdot \chi_{A_i}$ definieren wir
	\[
		\int\! s  \, \mathrm{d}\mu := \sum_{i=1}^{k} \alpha_i \cdot \mu(A_i)
	\]
	\item Für $f : X \to [0,\infty]$ messbar definieren wir \marginnote{vgl. \ref{sub:913}}
	\[
		\int\! f  \, \mathrm{d}\mu := \sup \set[\int\! s  \, \mathrm{d}\mu]{s \text{ ist einfach, messbar}, 0 \le s \le f} 
	\]
	\item Falls $f$ nicht notwendig positiv ist, aber $\int\! f_-  \, \mathrm{d}\mu < \infty$, so setzen wir $\int\! f  \, \mathrm{d}\mu := \int\! f_+  \, \mathrm{d}\mu -
	\int\! f_-  \, \mathrm{d}\mu $.
	\item Wir schreiben $\mathcal{L}^1(\mu) := \set[f : X \to \mathds{R}]{\int\! \abs{f}  \, \mathrm{d}\mu < \infty }$. $\mathcal{L}^1(\mu)$ ist Vektorraum und die
	Abbildung $f \mapsto\int\! f  \, \mathrm{d}\mu$ ist lineares Funktional auf $\mathcal{L}^1(\mu)$.
	\item Für $f \in \mathcal{L}^1(\mu)$ setzen wir $\norm[1]{f} = \int\! \abs{f}  \, \mathrm{d}\mu$. $\norm[1]{.} $ ist Halbnorm.
\end{enumerate}
% subsection 917 (end)

\subsection[Satz: Das Lebesgue-Maß definiert das Lebesgue-Integral!]{Satz} % (fold)
\label{sub:918}
Für $f \in \mathcal{L}^1(\mathds{R}^n)$ gilt
\[
	\int\! f(x)  \, \mathrm{d}x  = \int\! f  \, \mathrm{d}\lambda 
\]
\minisec{Beweis}
Klar für einfache Funktionen (Linearität). Benutze dann \ref{sub:913} und Beppo-Levi für $f_+$ und $f_-$. \bewende
% subsection 918 (end)

\subsection[Bemerkung: Aussage von Satz \ref{sub:918}]{Bemerkung} % (fold)
\label{sub:919}
Wir haben vorher das Lebesgue-Integral definiert und daraus das Lebesgue-Maß erhalten. Der Satz sagt, dass man umgekehrt das Integral von dem Maß erhält.
% subsection 919 (end)

\subsection[Definition: Quotientenvektorraum $L^1(\mu)$]{Definition} % (fold)
\label{sub:920}
$(X, \Sigma, \mu)$ Maßraum. Wir definieren wieder $\mathcal{N} := \set[f \in \mathcal{L}^1(\mu)]{\int\! \abs{f}  \, \mathrm{d}\mu = 0 } $ und $L^1(\mu) := \nicefrac{\mathcal{L}^1(\mu)}{\mathcal{N}}$.
% subsection 920 (end)
\bigskip
Die folgenden Sätze beweist man wie in Abschnitt \ref{sec:6}

\subsection[Satz: $L^1(\mu)$ ist ein Banachraum]{Satz} % (fold)
\label{sub:921}
$(L^1(\mu), \norm[1]{.} )$ ist vollständig normierter Vektorraum (Banachraum).
% subsection 921 (end)

\subsection{Satz (von Lebesgue über monotone Konvergenz)} % (fold)
\label{sub:922}
Sei $(X, \Sigma, \mu)$ ein Maßraum und $f_n : X \to [0,\infty]$, $n \in \mathds{N}$ eine Folge messbarer Funktionen mit $0 \le f_0 \le f_1 \le \ldots $. Dann ist 
$f := \lim_{ n \to \infty} f_n : X \to [0,\infty]$ messbar und es gilt
\[
	\int\! f_n  \, \mathrm{d}\mu \xrightarrow{n \to \infty} \int\! f  \, \mathrm{d}\mu.
\]
% subsection 922 (end)

\subsection{Lemma von Fatou} % (fold)
\label{sub:923}
Sei $(X, \Sigma, \mu)$ ein Maßraum und $f_n : X \to  [0,\infty]$, $n \in \mathds{N}$, eine Folge messbarer Funktionen. Dann gilt 
\[
	\int\! \liminf_{n \to \infty} f_n  \, \mathrm{d}\mu \le \liminf_{n \to \infty} \int\! f_n  \, \mathrm{d}\mu
\]
% subsection 923 (end)

\subsection{Satz (von Lebesgue über dominierte Konvergenz)} % (fold)
\label{sub:924}
Sei $(X, \Sigma,  \mu)$ ein Maßraum und $f_n : X \to (-\infty,\infty]$, $n \in \mathds{N} $, eine Folge messbarer Funktionen, welche punktweise gegen 
$f : X  \to (-\infty, \infty]$ konvergiert. Es sei $g : X \to [0,\infty]$ messbar mit $\int\! g  \, \mathrm{d}\mu < \infty$; es gelte $\abs{f_n} \le g $.
Dann ist $f := \lim_{ n \to \infty} f_n$ integrierbar und es gilt
\[
	\int\! f_n  \, \mathrm{d}\mu \xrightarrow{n \to \infty} \int\! f  \, \mathrm{d}\mu.
\]
% subsection 924 (end)

Haben gesehen: Maß $\mu \leadsto$ lineares Funktional auf $\mathcal{L}^1(\mu)$. Unter gewissen Umständen lassen sich umgekehrt Funktionale auf Funktionsräumen durch Integrale
darstellen.

\subsection*{Erinnerung} % (fold)
\label{sub:erinnerung}
Ein topologischer Raum $(X, \mathcal{T} )$ heißt \Index{lokalkompakt} und \Index{Hausdorff} falls gilt:
\begin{enumerate}[(i)]
	\item Jeder Punkt besitzt eine kompakte Umgebung
	\item Je zwei Punkte lassen sich durch offene Mengen trennen.
\end{enumerate}
D.h. 
\begin{enumerate}[(i)]
	\item $\forall x \in X : \exists K \text{ kompakt }\supset V \text{ offen } \ni x$
	\item $\forall x \not= y \in X : \exists V,W \in \mathcal{T}: x \in V, y \in W, V \cap W = \emptyset$
\end{enumerate} 
$C_c(X) := \set[f : X \to \mathds{R}]{\supp f \text{ kompakt}} $ 
% subsection erinnerung (end)

\subsection{Satz (Darstellungssatz von Riesz)} % (fold)
\label{sub:925}
Sei $(X,\mathcal{T})$ ein lokalkompakter Hausdorffraum und sei $\Phi : C_c(X) \to \mathds{R}$ ein positives $(f \ge 0 \Rightarrow \Phi(f)\ge 0)$ lineares Funktional.
Dann existieren eine $\sigma$-Algebra $\Sigma$ auf $X$ mit $\mathcal{B}_X \subset \Sigma$ und ein Maß $\mu$ auf $(X, \Sigma)$ mit folgenden Eigenschaften:
\begin{enumerate}[(i)]
	\item $\Phi(f) = \int\! f  \, \mathrm{d}\mu$ für $f \in C_c(X)$.
	\item $\mu(K)< \infty$ für $K$ kompakt. 
	\item Für $E \in \Sigma$ gilt $\mu(E)= \inf\set[\mu(V)]{E \subset V \in \mathcal{T} } $
	\item Für $E \in \Sigma$ mit $\mu(E)< \infty$ gilt $\mu(E) = \sup \set[\mu(K)]{K \subseteq E \text{ kompakt}} $
	\item Falls $E \in \Sigma$, $\mu(E)=0$, $A \subseteq E$ dann gilt $A \in \Sigma$.
\end{enumerate} 
Außerdem ist $\mu$ durch diese Eigenschaften eindeutig bestimmt. \bewende(ohne Beweis)
\minisec{Bemerkung}
\begin{enumerate}[a)]
	\item (iii) und (iv): $\mu$ ist \Index{reguläres Borelmaß} (von außen und von innen regulär) 
	
	(v): $\mu$ ist "{}vollständig"{}, d.h. Teilmengen von Nullmengen sind wieder Nullmengen.
	\item Der Satz erlaubt es, das Lebesgue-Maß/Integral aus dem Regelintegral für stetige Funktionen mit kompaktem Träger zu gewinnen.
\end{enumerate}
% subsection 925 (end)
% section 9 (end)
\newpage

\section{Ausblick} % (fold)
\label{sec:10}
\begin{description}
	\item[lokalkompakte Gruppen:] $\leadsto \exists$ (links)invariantes reguläres Borelmaß $\lambda$. $\lambda (g \cdot E) = \lambda (E)$ $\leadsto L^1(G,\lambda) $.
	Faltung $ \leadsto$ Fourieranalysis
	\item[Dynamische Systeme:] $X$ topologischer Raum, $\alpha : X \to X$ bijektiv, stetig Borel. Existenz von invarianten Borelmaßen? $\leadsto$ Ergodentheorie.
	\item[Mannigfaltigkeiten:] $M = \bigcup_{i \in I} U_i$ wo $U_i \approx \mathds{R}^n$
	\item[Räume von Funktionen und Abbildungen zwischen ihnen] $L^2(\mu) := \set[f :X \to \mathds{R}]{\int \abs{f}^2  \, \mathrm{d}\mu < \infty } $
	Skalarprodukt auf $L^2$ $\leadsto $ Hilbertraum, unendlich dimensionaler euklidischer Vektorraum.
\end{description}
% section 10 (end)





\cleardoubleoddemptypage
\pagenumbering{Alph}
\setcounter{page}{1}
\printindex
\listoffigures
\end{document}
